\documentclass[11pt,a4paper]{article}

% Biên dịch bằng pdfLaTeX (chạy 2 lần để có mục lục):
%   python -m src.v2.eval.xuat_so_lieu_bao_cao
%   cd docs/v2/bao_cao && pdflatex bao_cao_v2.tex && pdflatex bao_cao_v2.tex
% Mọi con số lấy qua \so{khoa} từ so_lieu.tex (sinh tự động từ results/v2). Không gõ số tay.
\usepackage[utf8]{inputenc}
\usepackage[T5]{fontenc}
\usepackage[vietnamese]{babel}
\usepackage[a4paper,left=25mm,right=22mm,top=22mm,bottom=22mm]{geometry}
\usepackage{lmodern}
\usepackage{graphicx}
\usepackage{booktabs}
\usepackage{tabularx}
\usepackage{array}
\usepackage{multirow}
\usepackage{enumitem}
\usepackage{amsmath}
\usepackage[table]{xcolor}
\usepackage{float}
\usepackage{caption}
\usepackage{fancyhdr}
\usepackage{tikz}
\usetikzlibrary{positioning,arrows.meta,shapes.geometric}
\usepackage{pgfplots}
\pgfplotsset{compat=1.18}
\PassOptionsToPackage{obeyspaces,spaces}{url}
\usepackage{hyperref}
\usepackage{microtype}

\definecolor{navy}{HTML}{17365D}
\definecolor{nenxanh}{HTML}{EAF2F8}
\definecolor{nenxam}{HTML}{F3F4F6}
\definecolor{chuxam}{HTML}{666666}
% Màu 3 hệ trong biểu đồ (thứ tự cố định: v1, LLM API, v2)
\definecolor{mauvmot}{HTML}{257DE0}
\definecolor{maullm}{HTML}{F47230}
\definecolor{mauvhai}{HTML}{00B894}

\hypersetup{
  colorlinks=true, linkcolor=navy, urlcolor=navy, citecolor=navy,
  pdftitle={SafeViChi v2 - Báo cáo vòng chung kết},
  pdfauthor={Nhóm SafeViChi}
}
\setlength{\parindent}{0pt}
\setlength{\parskip}{0.5em}
\setlist{itemsep=0.2em,topsep=0.2em}
\captionsetup{font=small,labelfont=bf}
\renewcommand{\arraystretch}{1.15}
\newcolumntype{Y}{>{\raggedright\arraybackslash}X}
\newcolumntype{C}{>{\centering\arraybackslash}X}

% ---------- Số liệu sinh tự động ----------
\makeatletter
\newcommand{\soset}[2]{\expandafter\def\csname so@#1\endcsname{#2}}
\newcommand{\so}[1]{\ifcsname so@#1\endcsname\csname so@#1\endcsname\else
  \textbf{\textcolor{red}{??}}\PackageWarning{so}{Thieu so lieu: #1}\fi}
\makeatother
% Tệp sinh tự động bằng `python -m src.v2.eval.xuat_so_lieu_bao_cao`. Không sửa tay.
% Nguồn (sha256 12 ký tự đầu):
%   data/v2/interim/chon_gan_lai_stats.json  ec2ffd0e9165
%   data/v2/interim/merged_stats.json  ea9ca5b73f3c
%   data/v2/interim/sinh_llm_loc_stats.json  03026021a523
%   data/v2/interim/sinh_llm_stats.json  cf07668235b0
%   data/v2/interim/that_gan_lai_stats.json  58502d66ee37
%   data/v2/test_stats.json  400a5b585c10
%   data/v2/train_stats.json  042c72f18f5a
%   demo/model-manifest.json  bb7db6c14f3f
%   results/v2/chi_phi_du_lieu/chi_phi.json  444438f46f07
%   results/v2/dev_selection/dev_summary.json  eba74eff750a
%   results/v2/do_thiet_bi/cpu_laptop.json  1b5ef4ab5fa8
%   results/v2/do_thiet_bi/trinh_duyet.json  71cd51a66c5c
%   results/v2/dong_thuan_nhan/metrics.json  c5192197b320
%   results/v2/kiem_tra_ro_ri/bao_cao.json  bb7cb1685f9b
%   results/v2/llm_api/test_metrics.json  1cbf24b234aa
%   results/v2/llm_api/test_perturb_C1_metrics.json  c80bed6119a6
%   results/v2/llm_api/test_perturb_C2_metrics.json  d38aa545db56
%   results/v2/official_eval/dualorc_mmbert_base/benchmark.json  766a6fdfbc94
%   results/v2/official_eval/dualorc_mmbert_base/calibration.json  a882c90cc1d8
%   results/v2/official_eval/dualorc_mmbert_base/config.json  64528aa5e799
%   results/v2/official_eval/dualorc_mmbert_base/dev_context_checks.json  fdabcedaed97
%   results/v2/official_eval/dualorc_mmbert_base/dev_norm_stress_tests.json  bc069a3ad5e7
%   results/v2/official_eval/dualorc_mmbert_base/dev_stress_tests.json  0c118a4879fa
%   results/v2/official_eval/dualorc_mmbert_base/environment.json  6e296cee27b3
%   results/v2/official_eval/dualorc_mmbert_base/test/bench_cpu.json  7f705f3e9640
%   results/v2/official_eval/dualorc_mmbert_base/test/int8_comparison.json  1c3c2eac2e9b
%   results/v2/official_eval/dualorc_mmbert_base/test/real_test_all_predictions.jsonl  dfe1d08a9b8a
%   results/v2/official_eval/dualorc_mmbert_base/test/real_test_norm_official_matrix.json  83269fd6528e
%   results/v2/official_eval/dualorc_mmbert_base/test/real_test_official_matrix.json  1a7f54368af1
%   results/v2/so_sanh_3_he/metrics.json  4437476f71f5
%   v1/results/step3/test/matrix.json  deb27cc2872a
\soset{C0.v1.n}{334}
\soset{C0.v1.bin}{31,2}
\soset{C0.v1.p}{100,0}
\soset{C0.v1.r}{18,5}
\soset{C0.v1.fpr}{0,0\%}
\soset{C0.v1.at.r}{100,0\%}
\soset{C0.v1.xp.r}{73,6\%}
\soset{C0.v1.qr.r}{0,0\%}
\soset{C0.v1.dd.r}{8,1\%}
\soset{C0.v1.ld.r}{0,0\%}
\soset{C0.llm.n}{334}
\soset{C0.llm.bin}{89,2}
\soset{C0.llm.p}{98,5}
\soset{C0.llm.r}{81,5}
\soset{C0.llm.fpr}{3,1\%}
\soset{C0.llm.macro}{77,2}
\soset{C0.llm.bcp}{83,2}
\soset{C0.llm.bcr}{66,3}
\soset{C0.llm.bcf}{70,7}
\soset{C0.llm.bctu}{65,5}
\soset{C0.llm.nbc}{194}
\soset{C0.llm.at.p}{67,9}
\soset{C0.llm.at.r}{96,9}
\soset{C0.llm.at.f}{79,8}
\soset{C0.llm.at.n}{96}
\soset{C0.llm.xp.p}{81,7}
\soset{C0.llm.xp.r}{92,5}
\soset{C0.llm.xp.f}{86,7}
\soset{C0.llm.xp.n}{53}
\soset{C0.llm.qr.p}{88,5}
\soset{C0.llm.qr.r}{45,1}
\soset{C0.llm.qr.f}{59,7}
\soset{C0.llm.qr.n}{51}
\soset{C0.llm.dd.p}{80,5}
\soset{C0.llm.dd.r}{53,2}
\soset{C0.llm.dd.f}{64,1}
\soset{C0.llm.dd.n}{62}
\soset{C0.llm.ld.p}{97,1}
\soset{C0.llm.ld.r}{94,4}
\soset{C0.llm.ld.f}{95,8}
\soset{C0.llm.ld.n}{72}
\soset{C0.llm.macro.lo}{72,4}
\soset{C0.llm.macro.hi}{81,6}
\soset{C0.llm.bin.lo}{85,9}
\soset{C0.llm.bin.hi}{92,0}
\soset{C0.llm.cm.at.at}{93}
\soset{C0.llm.cm.at.xp}{1}
\soset{C0.llm.cm.at.qr}{0}
\soset{C0.llm.cm.at.dd}{2}
\soset{C0.llm.cm.at.ld}{0}
\soset{C0.llm.cm.xp.at}{4}
\soset{C0.llm.cm.xp.xp}{49}
\soset{C0.llm.cm.xp.qr}{0}
\soset{C0.llm.cm.xp.dd}{0}
\soset{C0.llm.cm.xp.ld}{0}
\soset{C0.llm.cm.qr.at}{15}
\soset{C0.llm.cm.qr.xp}{5}
\soset{C0.llm.cm.qr.qr}{23}
\soset{C0.llm.cm.qr.dd}{6}
\soset{C0.llm.cm.qr.ld}{2}
\soset{C0.llm.cm.dd.at}{21}
\soset{C0.llm.cm.dd.xp}{5}
\soset{C0.llm.cm.dd.qr}{3}
\soset{C0.llm.cm.dd.dd}{33}
\soset{C0.llm.cm.dd.ld}{0}
\soset{C0.llm.cm.ld.at}{4}
\soset{C0.llm.cm.ld.xp}{0}
\soset{C0.llm.cm.ld.qr}{0}
\soset{C0.llm.cm.ld.dd}{0}
\soset{C0.llm.cm.ld.ld}{68}
\soset{C0.llm.xp.bcf}{74,1}
\soset{C0.llm.qr.bcf}{67,5}
\soset{C0.llm.dd.bcf}{59,5}
\soset{C0.llm.ld.bcf}{76,8}
\soset{C0.v2.n}{334}
\soset{C0.v2.bin}{92,9}
\soset{C0.v2.p}{95,6}
\soset{C0.v2.r}{90,3}
\soset{C0.v2.fpr}{10,4\%}
\soset{C0.v2.macro}{86,5}
\soset{C0.v2.bcp}{85,7}
\soset{C0.v2.bcr}{84,6}
\soset{C0.v2.bcf}{83,0}
\soset{C0.v2.bctu}{79,3}
\soset{C0.v2.nbc}{215}
\soset{C0.v2.at.p}{78,9}
\soset{C0.v2.at.r}{89,6}
\soset{C0.v2.at.f}{83,9}
\soset{C0.v2.at.n}{96}
\soset{C0.v2.xp.p}{98,0}
\soset{C0.v2.xp.r}{90,6}
\soset{C0.v2.xp.f}{94,1}
\soset{C0.v2.xp.n}{53}
\soset{C0.v2.qr.p}{87,8}
\soset{C0.v2.qr.r}{70,6}
\soset{C0.v2.qr.f}{78,3}
\soset{C0.v2.qr.n}{51}
\soset{C0.v2.dd.p}{79,7}
\soset{C0.v2.dd.r}{82,3}
\soset{C0.v2.dd.f}{81,0}
\soset{C0.v2.dd.n}{62}
\soset{C0.v2.ld.p}{95,8}
\soset{C0.v2.ld.r}{94,4}
\soset{C0.v2.ld.f}{95,1}
\soset{C0.v2.ld.n}{72}
\soset{C0.v2.macro.lo}{82,5}
\soset{C0.v2.macro.hi}{89,9}
\soset{C0.v2.bin.lo}{90,2}
\soset{C0.v2.bin.hi}{95,1}
\soset{C0.v2.cm.at.at}{86}
\soset{C0.v2.cm.at.xp}{0}
\soset{C0.v2.cm.at.qr}{1}
\soset{C0.v2.cm.at.dd}{6}
\soset{C0.v2.cm.at.ld}{3}
\soset{C0.v2.cm.xp.at}{2}
\soset{C0.v2.cm.xp.xp}{48}
\soset{C0.v2.cm.xp.qr}{3}
\soset{C0.v2.cm.xp.dd}{0}
\soset{C0.v2.cm.xp.ld}{0}
\soset{C0.v2.cm.qr.at}{8}
\soset{C0.v2.cm.qr.xp}{0}
\soset{C0.v2.cm.qr.qr}{36}
\soset{C0.v2.cm.qr.dd}{7}
\soset{C0.v2.cm.qr.ld}{0}
\soset{C0.v2.cm.dd.at}{9}
\soset{C0.v2.cm.dd.xp}{1}
\soset{C0.v2.cm.dd.qr}{1}
\soset{C0.v2.cm.dd.dd}{51}
\soset{C0.v2.cm.dd.ld}{0}
\soset{C0.v2.cm.ld.at}{4}
\soset{C0.v2.cm.ld.xp}{0}
\soset{C0.v2.cm.ld.qr}{0}
\soset{C0.v2.cm.ld.dd}{0}
\soset{C0.v2.cm.ld.ld}{68}
\soset{C0.v2.xp.bcf}{91,8}
\soset{C0.v2.qr.bcf}{78,1}
\soset{C0.v2.dd.bcf}{76,4}
\soset{C0.v2.ld.bcf}{84,7}
\soset{C0.bac.v1.n}{112}
\soset{C0.bac.v1.bin}{31,9}
\soset{C0.bac.v1.p}{100,0}
\soset{C0.bac.v1.r}{19,0}
\soset{C0.bac.v1.fpr}{0,0\%}
\soset{C0.bac.v1.at.r}{100,0\%}
\soset{C0.bac.v1.xp.r}{68,4\%}
\soset{C0.bac.v1.qr.r}{0,0\%}
\soset{C0.bac.v1.dd.r}{9,1\%}
\soset{C0.bac.v1.ld.r}{0,0\%}
\soset{C0.bac.llm.n}{112}
\soset{C0.bac.llm.bin}{88,9}
\soset{C0.bac.llm.p}{98,5}
\soset{C0.bac.llm.r}{81,0}
\soset{C0.bac.llm.fpr}{3,0\%}
\soset{C0.bac.llm.macro}{75,6}
\soset{C0.bac.llm.bcp}{90,0}
\soset{C0.bac.llm.bcr}{71,3}
\soset{C0.bac.llm.bcf}{75,9}
\soset{C0.bac.llm.bctu}{71,5}
\soset{C0.bac.llm.nbc}{64}
\soset{C0.bac.llm.at.p}{68,1}
\soset{C0.bac.llm.at.r}{97,0}
\soset{C0.bac.llm.at.f}{80,0}
\soset{C0.bac.llm.at.n}{33}
\soset{C0.bac.llm.xp.p}{89,5}
\soset{C0.bac.llm.xp.r}{89,5}
\soset{C0.bac.llm.xp.f}{89,5}
\soset{C0.bac.llm.xp.n}{19}
\soset{C0.bac.llm.qr.p}{80,0}
\soset{C0.bac.llm.qr.r}{28,6}
\soset{C0.bac.llm.qr.f}{42,1}
\soset{C0.bac.llm.qr.n}{14}
\soset{C0.bac.llm.dd.p}{81,2}
\soset{C0.bac.llm.dd.r}{59,1}
\soset{C0.bac.llm.dd.f}{68,4}
\soset{C0.bac.llm.dd.n}{22}
\soset{C0.bac.llm.ld.p}{96,0}
\soset{C0.bac.llm.ld.r}{100,0}
\soset{C0.bac.llm.ld.f}{98,0}
\soset{C0.bac.llm.ld.n}{24}
\soset{C0.bac.v2.n}{112}
\soset{C0.bac.v2.bin}{92,9}
\soset{C0.bac.v2.p}{94,7}
\soset{C0.bac.v2.r}{91,1}
\soset{C0.bac.v2.fpr}{12,1\%}
\soset{C0.bac.v2.macro}{86,4}
\soset{C0.bac.v2.bcp}{86,1}
\soset{C0.bac.v2.bcr}{85,9}
\soset{C0.bac.v2.bcf}{83,9}
\soset{C0.bac.v2.bctu}{80,0}
\soset{C0.bac.v2.nbc}{72}
\soset{C0.bac.v2.at.p}{80,6}
\soset{C0.bac.v2.at.r}{87,9}
\soset{C0.bac.v2.at.f}{84,1}
\soset{C0.bac.v2.at.n}{33}
\soset{C0.bac.v2.xp.p}{94,1}
\soset{C0.bac.v2.xp.r}{84,2}
\soset{C0.bac.v2.xp.f}{88,9}
\soset{C0.bac.v2.xp.n}{19}
\soset{C0.bac.v2.qr.p}{84,6}
\soset{C0.bac.v2.qr.r}{78,6}
\soset{C0.bac.v2.qr.f}{81,5}
\soset{C0.bac.v2.qr.n}{14}
\soset{C0.bac.v2.dd.p}{81,8}
\soset{C0.bac.v2.dd.r}{81,8}
\soset{C0.bac.v2.dd.f}{81,8}
\soset{C0.bac.v2.dd.n}{22}
\soset{C0.bac.v2.ld.p}{95,8}
\soset{C0.bac.v2.ld.r}{95,8}
\soset{C0.bac.v2.ld.f}{95,8}
\soset{C0.bac.v2.ld.n}{24}
\soset{C0.trung.v1.n}{111}
\soset{C0.trung.v1.bin}{24,4}
\soset{C0.trung.v1.p}{100,0}
\soset{C0.trung.v1.r}{13,9}
\soset{C0.trung.v1.fpr}{0,0\%}
\soset{C0.trung.v1.at.r}{100,0\%}
\soset{C0.trung.v1.xp.r}{64,7\%}
\soset{C0.trung.v1.qr.r}{0,0\%}
\soset{C0.trung.v1.dd.r}{0,0\%}
\soset{C0.trung.v1.ld.r}{0,0\%}
\soset{C0.trung.llm.n}{111}
\soset{C0.trung.llm.bin}{88,1}
\soset{C0.trung.llm.p}{98,4}
\soset{C0.trung.llm.r}{79,8}
\soset{C0.trung.llm.fpr}{3,1\%}
\soset{C0.trung.llm.macro}{72,8}
\soset{C0.trung.llm.bcp}{75,8}
\soset{C0.trung.llm.bcr}{61,5}
\soset{C0.trung.llm.bcf}{64,9}
\soset{C0.trung.llm.bctu}{59,7}
\soset{C0.trung.llm.nbc}{63}
\soset{C0.trung.llm.at.p}{66,0}
\soset{C0.trung.llm.at.r}{96,9}
\soset{C0.trung.llm.at.f}{78,5}
\soset{C0.trung.llm.at.n}{32}
\soset{C0.trung.llm.xp.p}{68,2}
\soset{C0.trung.llm.xp.r}{88,2}
\soset{C0.trung.llm.xp.f}{76,9}
\soset{C0.trung.llm.xp.n}{17}
\soset{C0.trung.llm.qr.p}{90,9}
\soset{C0.trung.llm.qr.r}{52,6}
\soset{C0.trung.llm.qr.f}{66,7}
\soset{C0.trung.llm.qr.n}{19}
\soset{C0.trung.llm.dd.p}{75,0}
\soset{C0.trung.llm.dd.r}{35,3}
\soset{C0.trung.llm.dd.f}{48,0}
\soset{C0.trung.llm.dd.n}{17}
\soset{C0.trung.llm.ld.p}{100,0}
\soset{C0.trung.llm.ld.r}{88,5}
\soset{C0.trung.llm.ld.f}{93,9}
\soset{C0.trung.llm.ld.n}{26}
\soset{C0.trung.v2.n}{111}
\soset{C0.trung.v2.bin}{94,0}
\soset{C0.trung.v2.p}{98,6}
\soset{C0.trung.v2.r}{89,9}
\soset{C0.trung.v2.fpr}{3,1\%}
\soset{C0.trung.v2.macro}{84,5}
\soset{C0.trung.v2.bcp}{90,0}
\soset{C0.trung.v2.bcr}{86,4}
\soset{C0.trung.v2.bcf}{85,6}
\soset{C0.trung.v2.bctu}{81,7}
\soset{C0.trung.v2.nbc}{71}
\soset{C0.trung.v2.at.p}{79,5}
\soset{C0.trung.v2.at.r}{96,9}
\soset{C0.trung.v2.at.f}{87,3}
\soset{C0.trung.v2.at.n}{32}
\soset{C0.trung.v2.xp.p}{100,0}
\soset{C0.trung.v2.xp.r}{94,1}
\soset{C0.trung.v2.xp.f}{97,0}
\soset{C0.trung.v2.xp.n}{17}
\soset{C0.trung.v2.qr.p}{81,8}
\soset{C0.trung.v2.qr.r}{47,4}
\soset{C0.trung.v2.qr.f}{60,0}
\soset{C0.trung.v2.qr.n}{19}
\soset{C0.trung.v2.dd.p}{77,8}
\soset{C0.trung.v2.dd.r}{82,3}
\soset{C0.trung.v2.dd.f}{80,0}
\soset{C0.trung.v2.dd.n}{17}
\soset{C0.trung.v2.ld.p}{96,3}
\soset{C0.trung.v2.ld.r}{100,0}
\soset{C0.trung.v2.ld.f}{98,1}
\soset{C0.trung.v2.ld.n}{26}
\soset{C0.nam.v1.n}{111}
\soset{C0.nam.v1.bin}{36,7}
\soset{C0.nam.v1.p}{100,0}
\soset{C0.nam.v1.r}{22,5}
\soset{C0.nam.v1.fpr}{0,0\%}
\soset{C0.nam.v1.at.r}{100,0\%}
\soset{C0.nam.v1.xp.r}{88,2\%}
\soset{C0.nam.v1.qr.r}{0,0\%}
\soset{C0.nam.v1.dd.r}{13,0\%}
\soset{C0.nam.v1.ld.r}{0,0\%}
\soset{C0.nam.llm.n}{111}
\soset{C0.nam.llm.bin}{90,5}
\soset{C0.nam.llm.p}{98,5}
\soset{C0.nam.llm.r}{83,8}
\soset{C0.nam.llm.fpr}{3,2\%}
\soset{C0.nam.llm.macro}{81,0}
\soset{C0.nam.llm.bcp}{83,7}
\soset{C0.nam.llm.bcr}{66,0}
\soset{C0.nam.llm.bcf}{71,2}
\soset{C0.nam.llm.bctu}{65,0}
\soset{C0.nam.llm.nbc}{67}
\soset{C0.nam.llm.at.p}{69,8}
\soset{C0.nam.llm.at.r}{96,8}
\soset{C0.nam.llm.at.f}{81,1}
\soset{C0.nam.llm.at.n}{31}
\soset{C0.nam.llm.xp.p}{89,5}
\soset{C0.nam.llm.xp.r}{100,0}
\soset{C0.nam.llm.xp.f}{94,4}
\soset{C0.nam.llm.xp.n}{17}
\soset{C0.nam.llm.qr.p}{90,0}
\soset{C0.nam.llm.qr.r}{50,0}
\soset{C0.nam.llm.qr.f}{64,3}
\soset{C0.nam.llm.qr.n}{18}
\soset{C0.nam.llm.dd.p}{82,3}
\soset{C0.nam.llm.dd.r}{60,9}
\soset{C0.nam.llm.dd.f}{70,0}
\soset{C0.nam.llm.dd.n}{23}
\soset{C0.nam.llm.ld.p}{95,5}
\soset{C0.nam.llm.ld.r}{95,5}
\soset{C0.nam.llm.ld.f}{95,5}
\soset{C0.nam.llm.ld.n}{22}
\soset{C0.nam.v2.n}{111}
\soset{C0.nam.v2.bin}{91,7}
\soset{C0.nam.v2.p}{93,5}
\soset{C0.nam.v2.r}{90,0}
\soset{C0.nam.v2.fpr}{16,1\%}
\soset{C0.nam.v2.macro}{87,9}
\soset{C0.nam.v2.bcp}{81,2}
\soset{C0.nam.v2.bcr}{81,7}
\soset{C0.nam.v2.bcf}{79,7}
\soset{C0.nam.v2.bctu}{76,3}
\soset{C0.nam.v2.nbc}{72}
\soset{C0.nam.v2.at.p}{76,5}
\soset{C0.nam.v2.at.r}{83,9}
\soset{C0.nam.v2.at.f}{80,0}
\soset{C0.nam.v2.at.n}{31}
\soset{C0.nam.v2.xp.p}{100,0}
\soset{C0.nam.v2.xp.r}{94,1}
\soset{C0.nam.v2.xp.f}{97,0}
\soset{C0.nam.v2.xp.n}{17}
\soset{C0.nam.v2.qr.p}{94,1}
\soset{C0.nam.v2.qr.r}{88,9}
\soset{C0.nam.v2.qr.f}{91,4}
\soset{C0.nam.v2.qr.n}{18}
\soset{C0.nam.v2.dd.p}{79,2}
\soset{C0.nam.v2.dd.r}{82,6}
\soset{C0.nam.v2.dd.f}{80,8}
\soset{C0.nam.v2.dd.n}{23}
\soset{C0.nam.v2.ld.p}{95,0}
\soset{C0.nam.v2.ld.r}{86,4}
\soset{C0.nam.v2.ld.f}{90,5}
\soset{C0.nam.v2.ld.n}{22}
\soset{C0.llm.gap.macro}{8,3}
\soset{C0.llm.gap.bin}{2,4}
\soset{C0.llm.gap.fpr}{0,2\%}
\soset{C0.v2.gap.macro}{3,5}
\soset{C0.v2.gap.bin}{2,3}
\soset{C0.v2.gap.fpr}{13,0\%}
\soset{C0.v1.gap.bin}{12,3}
\soset{C1.v1.n}{668}
\soset{C1.v1.bin}{31,0}
\soset{C1.v1.p}{96,7}
\soset{C1.v1.r}{18,5}
\soset{C1.v1.fpr}{1,6\%}
\soset{C1.v1.at.r}{98,4\%}
\soset{C1.v1.xp.r}{67,9\%}
\soset{C1.v1.qr.r}{2,0\%}
\soset{C1.v1.dd.r}{10,5\%}
\soset{C1.v1.ld.r}{0,7\%}
\soset{C1.llm.n}{668}
\soset{C1.llm.bin}{74,5}
\soset{C1.llm.p}{93,3}
\soset{C1.llm.r}{62,0}
\soset{C1.llm.fpr}{10,9\%}
\soset{C1.llm.macro}{58,5}
\soset{C1.llm.bcp}{73,4}
\soset{C1.llm.bcr}{50,2}
\soset{C1.llm.bcf}{56,1}
\soset{C1.llm.bctu}{49,8}
\soset{C1.llm.nbc}{295}
\soset{C1.llm.at.p}{48,6}
\soset{C1.llm.at.r}{89,1}
\soset{C1.llm.at.f}{62,9}
\soset{C1.llm.at.n}{192}
\soset{C1.llm.xp.p}{82,2}
\soset{C1.llm.xp.r}{69,8}
\soset{C1.llm.xp.f}{75,5}
\soset{C1.llm.xp.n}{106}
\soset{C1.llm.qr.p}{71,9}
\soset{C1.llm.qr.r}{22,6}
\soset{C1.llm.qr.f}{34,3}
\soset{C1.llm.qr.n}{102}
\soset{C1.llm.dd.p}{69,4}
\soset{C1.llm.dd.r}{20,2}
\soset{C1.llm.dd.f}{31,2}
\soset{C1.llm.dd.n}{124}
\soset{C1.llm.ld.p}{84,8}
\soset{C1.llm.ld.r}{93,1}
\soset{C1.llm.ld.f}{88,7}
\soset{C1.llm.ld.n}{144}
\soset{C1.llm.macro.lo}{54,4}
\soset{C1.llm.macro.hi}{62,1}
\soset{C1.llm.bin.lo}{70,6}
\soset{C1.llm.bin.hi}{77,6}
\soset{C1.llm.cm.at.at}{171}
\soset{C1.llm.cm.at.xp}{3}
\soset{C1.llm.cm.at.qr}{3}
\soset{C1.llm.cm.at.dd}{1}
\soset{C1.llm.cm.at.ld}{14}
\soset{C1.llm.cm.xp.at}{30}
\soset{C1.llm.cm.xp.xp}{74}
\soset{C1.llm.cm.xp.qr}{2}
\soset{C1.llm.cm.xp.dd}{0}
\soset{C1.llm.cm.xp.ld}{0}
\soset{C1.llm.cm.qr.at}{62}
\soset{C1.llm.cm.qr.xp}{3}
\soset{C1.llm.cm.qr.qr}{23}
\soset{C1.llm.cm.qr.dd}{8}
\soset{C1.llm.cm.qr.ld}{6}
\soset{C1.llm.cm.dd.at}{81}
\soset{C1.llm.cm.dd.xp}{10}
\soset{C1.llm.cm.dd.qr}{4}
\soset{C1.llm.cm.dd.dd}{25}
\soset{C1.llm.cm.dd.ld}{4}
\soset{C1.llm.cm.ld.at}{8}
\soset{C1.llm.cm.ld.xp}{0}
\soset{C1.llm.cm.ld.qr}{0}
\soset{C1.llm.cm.ld.dd}{2}
\soset{C1.llm.cm.ld.ld}{134}
\soset{C1.llm.xp.bcf}{54,2}
\soset{C1.llm.qr.bcf}{53,8}
\soset{C1.llm.dd.bcf}{35,0}
\soset{C1.llm.ld.bcf}{64,5}
\soset{C1.v2.n}{668}
\soset{C1.v2.bin}{90,3}
\soset{C1.v2.p}{95,1}
\soset{C1.v2.r}{85,9}
\soset{C1.v2.fpr}{10,9\%}
\soset{C1.v2.macro}{78,6}
\soset{C1.v2.bcp}{82,5}
\soset{C1.v2.bcr}{76,9}
\soset{C1.v2.bcf}{76,5}
\soset{C1.v2.bctu}{72,2}
\soset{C1.v2.nbc}{409}
\soset{C1.v2.at.p}{71,9}
\soset{C1.v2.at.r}{89,1}
\soset{C1.v2.at.f}{79,5}
\soset{C1.v2.at.n}{192}
\soset{C1.v2.xp.p}{85,8}
\soset{C1.v2.xp.r}{91,5}
\soset{C1.v2.xp.f}{88,6}
\soset{C1.v2.xp.n}{106}
\soset{C1.v2.qr.p}{83,9}
\soset{C1.v2.qr.r}{51,0}
\soset{C1.v2.qr.f}{63,4}
\soset{C1.v2.qr.n}{102}
\soset{C1.v2.dd.p}{70,8}
\soset{C1.v2.dd.r}{68,5}
\soset{C1.v2.dd.f}{69,7}
\soset{C1.v2.dd.n}{124}
\soset{C1.v2.ld.p}{94,8}
\soset{C1.v2.ld.r}{88,9}
\soset{C1.v2.ld.f}{91,8}
\soset{C1.v2.ld.n}{144}
\soset{C1.v2.macro.lo}{75,3}
\soset{C1.v2.macro.hi}{81,8}
\soset{C1.v2.bin.lo}{88,1}
\soset{C1.v2.bin.hi}{92,4}
\soset{C1.v2.cm.at.at}{171}
\soset{C1.v2.cm.at.xp}{4}
\soset{C1.v2.cm.at.qr}{3}
\soset{C1.v2.cm.at.dd}{8}
\soset{C1.v2.cm.at.ld}{6}
\soset{C1.v2.cm.xp.at}{5}
\soset{C1.v2.cm.xp.xp}{97}
\soset{C1.v2.cm.xp.qr}{3}
\soset{C1.v2.cm.xp.dd}{1}
\soset{C1.v2.cm.xp.ld}{0}
\soset{C1.v2.cm.qr.at}{21}
\soset{C1.v2.cm.qr.xp}{4}
\soset{C1.v2.cm.qr.qr}{52}
\soset{C1.v2.cm.qr.dd}{25}
\soset{C1.v2.cm.qr.ld}{0}
\soset{C1.v2.cm.dd.at}{26}
\soset{C1.v2.cm.dd.xp}{8}
\soset{C1.v2.cm.dd.qr}{4}
\soset{C1.v2.cm.dd.dd}{85}
\soset{C1.v2.cm.dd.ld}{1}
\soset{C1.v2.cm.ld.at}{15}
\soset{C1.v2.cm.ld.xp}{0}
\soset{C1.v2.cm.ld.qr}{0}
\soset{C1.v2.cm.ld.dd}{1}
\soset{C1.v2.cm.ld.ld}{128}
\soset{C1.v2.xp.bcf}{85,2}
\soset{C1.v2.qr.bcf}{66,0}
\soset{C1.v2.dd.bcf}{73,4}
\soset{C1.v2.ld.bcf}{78,6}
\soset{C1.bac.v1.n}{224}
\soset{C1.bac.v1.bin}{30,1}
\soset{C1.bac.v1.p}{100,0}
\soset{C1.bac.v1.r}{17,7}
\soset{C1.bac.v1.fpr}{0,0\%}
\soset{C1.bac.v1.at.r}{100,0\%}
\soset{C1.bac.v1.xp.r}{68,4\%}
\soset{C1.bac.v1.qr.r}{0,0\%}
\soset{C1.bac.v1.dd.r}{4,5\%}
\soset{C1.bac.v1.ld.r}{0,0\%}
\soset{C1.bac.llm.n}{224}
\soset{C1.bac.llm.bin}{71,8}
\soset{C1.bac.llm.p}{90,4}
\soset{C1.bac.llm.r}{59,5}
\soset{C1.bac.llm.fpr}{15,2\%}
\soset{C1.bac.llm.macro}{55,5}
\soset{C1.bac.llm.bcp}{77,7}
\soset{C1.bac.llm.bcr}{53,8}
\soset{C1.bac.llm.bcf}{60,1}
\soset{C1.bac.llm.bctu}{55,8}
\soset{C1.bac.llm.nbc}{94}
\soset{C1.bac.llm.at.p}{46,7}
\soset{C1.bac.llm.at.r}{84,9}
\soset{C1.bac.llm.at.f}{60,2}
\soset{C1.bac.llm.at.n}{66}
\soset{C1.bac.llm.xp.p}{96,0}
\soset{C1.bac.llm.xp.r}{63,2}
\soset{C1.bac.llm.xp.f}{76,2}
\soset{C1.bac.llm.xp.n}{38}
\soset{C1.bac.llm.qr.p}{55,6}
\soset{C1.bac.llm.qr.r}{17,9}
\soset{C1.bac.llm.qr.f}{27,0}
\soset{C1.bac.llm.qr.n}{28}
\soset{C1.bac.llm.dd.p}{57,1}
\soset{C1.bac.llm.dd.r}{18,2}
\soset{C1.bac.llm.dd.f}{27,6}
\soset{C1.bac.llm.dd.n}{44}
\soset{C1.bac.llm.ld.p}{80,4}
\soset{C1.bac.llm.ld.r}{93,8}
\soset{C1.bac.llm.ld.f}{86,5}
\soset{C1.bac.llm.ld.n}{48}
\soset{C1.bac.v2.n}{224}
\soset{C1.bac.v2.bin}{92,5}
\soset{C1.bac.v2.p}{95,9}
\soset{C1.bac.v2.r}{89,2}
\soset{C1.bac.v2.fpr}{9,1\%}
\soset{C1.bac.v2.macro}{81,6}
\soset{C1.bac.v2.bcp}{83,7}
\soset{C1.bac.v2.bcr}{77,3}
\soset{C1.bac.v2.bcf}{76,9}
\soset{C1.bac.v2.bctu}{72,2}
\soset{C1.bac.v2.nbc}{141}
\soset{C1.bac.v2.at.p}{77,9}
\soset{C1.bac.v2.at.r}{90,9}
\soset{C1.bac.v2.at.f}{83,9}
\soset{C1.bac.v2.at.n}{66}
\soset{C1.bac.v2.xp.p}{91,4}
\soset{C1.bac.v2.xp.r}{84,2}
\soset{C1.bac.v2.xp.f}{87,7}
\soset{C1.bac.v2.xp.n}{38}
\soset{C1.bac.v2.qr.p}{88,9}
\soset{C1.bac.v2.qr.r}{57,1}
\soset{C1.bac.v2.qr.f}{69,6}
\soset{C1.bac.v2.qr.n}{28}
\soset{C1.bac.v2.dd.p}{71,1}
\soset{C1.bac.v2.dd.r}{72,7}
\soset{C1.bac.v2.dd.f}{71,9}
\soset{C1.bac.v2.dd.n}{44}
\soset{C1.bac.v2.ld.p}{93,9}
\soset{C1.bac.v2.ld.r}{95,8}
\soset{C1.bac.v2.ld.f}{94,8}
\soset{C1.bac.v2.ld.n}{48}
\soset{C1.trung.v1.n}{222}
\soset{C1.trung.v1.bin}{22,2}
\soset{C1.trung.v1.p}{90,9}
\soset{C1.trung.v1.r}{12,7}
\soset{C1.trung.v1.fpr}{3,1\%}
\soset{C1.trung.v1.at.r}{96,9\%}
\soset{C1.trung.v1.xp.r}{55,9\%}
\soset{C1.trung.v1.qr.r}{0,0\%}
\soset{C1.trung.v1.dd.r}{2,9\%}
\soset{C1.trung.v1.ld.r}{0,0\%}
\soset{C1.trung.llm.n}{222}
\soset{C1.trung.llm.bin}{74,3}
\soset{C1.trung.llm.p}{94,2}
\soset{C1.trung.llm.r}{61,4}
\soset{C1.trung.llm.fpr}{9,4\%}
\soset{C1.trung.llm.macro}{57,8}
\soset{C1.trung.llm.bcp}{72,9}
\soset{C1.trung.llm.bcr}{49,2}
\soset{C1.trung.llm.bcf}{55,2}
\soset{C1.trung.llm.bctu}{48,0}
\soset{C1.trung.llm.nbc}{97}
\soset{C1.trung.llm.at.p}{48,7}
\soset{C1.trung.llm.at.r}{90,6}
\soset{C1.trung.llm.at.f}{63,4}
\soset{C1.trung.llm.at.n}{64}
\soset{C1.trung.llm.xp.p}{71,4}
\soset{C1.trung.llm.xp.r}{58,8}
\soset{C1.trung.llm.xp.f}{64,5}
\soset{C1.trung.llm.xp.n}{34}
\soset{C1.trung.llm.qr.p}{87,5}
\soset{C1.trung.llm.qr.r}{36,8}
\soset{C1.trung.llm.qr.f}{51,8}
\soset{C1.trung.llm.qr.n}{38}
\soset{C1.trung.llm.dd.p}{80,0}
\soset{C1.trung.llm.dd.r}{11,8}
\soset{C1.trung.llm.dd.f}{20,5}
\soset{C1.trung.llm.dd.n}{34}
\soset{C1.trung.llm.ld.p}{87,0}
\soset{C1.trung.llm.ld.r}{90,4}
\soset{C1.trung.llm.ld.f}{88,7}
\soset{C1.trung.llm.ld.n}{52}
\soset{C1.trung.v2.n}{222}
\soset{C1.trung.v2.bin}{89,7}
\soset{C1.trung.v2.p}{97,8}
\soset{C1.trung.v2.r}{82,9}
\soset{C1.trung.v2.fpr}{4,7\%}
\soset{C1.trung.v2.macro}{73,0}
\soset{C1.trung.v2.bcp}{83,1}
\soset{C1.trung.v2.bcr}{78,4}
\soset{C1.trung.v2.bcf}{78,1}
\soset{C1.trung.v2.bctu}{73,7}
\soset{C1.trung.v2.nbc}{131}
\soset{C1.trung.v2.at.p}{69,3}
\soset{C1.trung.v2.at.r}{95,3}
\soset{C1.trung.v2.at.f}{80,3}
\soset{C1.trung.v2.at.n}{64}
\soset{C1.trung.v2.xp.p}{83,8}
\soset{C1.trung.v2.xp.r}{91,2}
\soset{C1.trung.v2.xp.f}{87,3}
\soset{C1.trung.v2.xp.n}{34}
\soset{C1.trung.v2.qr.p}{68,8}
\soset{C1.trung.v2.qr.r}{28,9}
\soset{C1.trung.v2.qr.f}{40,7}
\soset{C1.trung.v2.qr.n}{38}
\soset{C1.trung.v2.dd.p}{65,6}
\soset{C1.trung.v2.dd.r}{61,8}
\soset{C1.trung.v2.dd.f}{63,6}
\soset{C1.trung.v2.dd.n}{34}
\soset{C1.trung.v2.ld.p}{95,9}
\soset{C1.trung.v2.ld.r}{90,4}
\soset{C1.trung.v2.ld.f}{93,1}
\soset{C1.trung.v2.ld.n}{52}
\soset{C1.nam.v1.n}{222}
\soset{C1.nam.v1.bin}{39,8}
\soset{C1.nam.v1.p}{97,6}
\soset{C1.nam.v1.r}{25,0}
\soset{C1.nam.v1.fpr}{1,6\%}
\soset{C1.nam.v1.at.r}{98,4\%}
\soset{C1.nam.v1.xp.r}{79,4\%}
\soset{C1.nam.v1.qr.r}{5,6\%}
\soset{C1.nam.v1.dd.r}{21,7\%}
\soset{C1.nam.v1.ld.r}{2,3\%}
\soset{C1.nam.llm.n}{222}
\soset{C1.nam.llm.bin}{77,3}
\soset{C1.nam.llm.p}{95,4}
\soset{C1.nam.llm.r}{65,0}
\soset{C1.nam.llm.fpr}{8,1\%}
\soset{C1.nam.llm.macro}{60,2}
\soset{C1.nam.llm.bcp}{70,0}
\soset{C1.nam.llm.bcr}{47,8}
\soset{C1.nam.llm.bcf}{53,3}
\soset{C1.nam.llm.bctu}{46,2}
\soset{C1.nam.llm.nbc}{104}
\soset{C1.nam.llm.at.p}{50,4}
\soset{C1.nam.llm.at.r}{91,9}
\soset{C1.nam.llm.at.f}{65,1}
\soset{C1.nam.llm.at.n}{62}
\soset{C1.nam.llm.xp.p}{81,1}
\soset{C1.nam.llm.xp.r}{88,2}
\soset{C1.nam.llm.xp.f}{84,5}
\soset{C1.nam.llm.xp.n}{34}
\soset{C1.nam.llm.qr.p}{57,1}
\soset{C1.nam.llm.qr.r}{11,1}
\soset{C1.nam.llm.qr.f}{18,6}
\soset{C1.nam.llm.qr.n}{36}
\soset{C1.nam.llm.dd.p}{76,5}
\soset{C1.nam.llm.dd.r}{28,3}
\soset{C1.nam.llm.dd.f}{41,3}
\soset{C1.nam.llm.dd.n}{46}
\soset{C1.nam.llm.ld.p}{87,5}
\soset{C1.nam.llm.ld.r}{95,5}
\soset{C1.nam.llm.ld.f}{91,3}
\soset{C1.nam.llm.ld.n}{44}
\soset{C1.nam.v2.n}{222}
\soset{C1.nam.v2.bin}{88,7}
\soset{C1.nam.v2.p}{92,0}
\soset{C1.nam.v2.r}{85,6}
\soset{C1.nam.v2.fpr}{19,4\%}
\soset{C1.nam.v2.macro}{80,2}
\soset{C1.nam.v2.bcp}{80,7}
\soset{C1.nam.v2.bcr}{75,1}
\soset{C1.nam.v2.bcf}{74,6}
\soset{C1.nam.v2.bctu}{70,8}
\soset{C1.nam.v2.nbc}{137}
\soset{C1.nam.v2.at.p}{68,5}
\soset{C1.nam.v2.at.r}{80,7}
\soset{C1.nam.v2.at.f}{74,1}
\soset{C1.nam.v2.at.n}{62}
\soset{C1.nam.v2.xp.p}{82,9}
\soset{C1.nam.v2.xp.r}{100,0}
\soset{C1.nam.v2.xp.f}{90,7}
\soset{C1.nam.v2.xp.n}{34}
\soset{C1.nam.v2.qr.p}{89,3}
\soset{C1.nam.v2.qr.r}{69,4}
\soset{C1.nam.v2.qr.f}{78,1}
\soset{C1.nam.v2.qr.n}{36}
\soset{C1.nam.v2.dd.p}{74,4}
\soset{C1.nam.v2.dd.r}{69,6}
\soset{C1.nam.v2.dd.f}{71,9}
\soset{C1.nam.v2.dd.n}{46}
\soset{C1.nam.v2.ld.p}{94,6}
\soset{C1.nam.v2.ld.r}{79,5}
\soset{C1.nam.v2.ld.f}{86,4}
\soset{C1.nam.v2.ld.n}{44}
\soset{C1.llm.gap.macro}{4,7}
\soset{C1.llm.gap.bin}{5,6}
\soset{C1.llm.gap.fpr}{7,1\%}
\soset{C1.v2.gap.macro}{8,6}
\soset{C1.v2.gap.bin}{3,8}
\soset{C1.v2.gap.fpr}{14,7\%}
\soset{C1.v1.gap.bin}{17,6}
\soset{C2.v1.n}{668}
\soset{C2.v1.bin}{39,5}
\soset{C2.v1.p}{93,7}
\soset{C2.v1.r}{25,0}
\soset{C2.v1.fpr}{4,2\%}
\soset{C2.v1.at.r}{95,8\%}
\soset{C2.v1.xp.r}{67,9\%}
\soset{C2.v1.qr.r}{9,8\%}
\soset{C2.v1.dd.r}{25,0\%}
\soset{C2.v1.ld.r}{4,2\%}
\soset{C2.llm.n}{668}
\soset{C2.llm.bin}{77,9}
\soset{C2.llm.p}{97,5}
\soset{C2.llm.r}{64,9}
\soset{C2.llm.fpr}{4,2\%}
\soset{C2.llm.macro}{62,7}
\soset{C2.llm.bcp}{44,2}
\soset{C2.llm.bcr}{28,4}
\soset{C2.llm.bcf}{32,2}
\soset{C2.llm.bctu}{28,1}
\soset{C2.llm.nbc}{309}
\soset{C2.llm.at.p}{52,4}
\soset{C2.llm.at.r}{95,8}
\soset{C2.llm.at.f}{67,8}
\soset{C2.llm.at.n}{192}
\soset{C2.llm.xp.p}{81,0}
\soset{C2.llm.xp.r}{72,6}
\soset{C2.llm.xp.f}{76,6}
\soset{C2.llm.xp.n}{106}
\soset{C2.llm.qr.p}{94,3}
\soset{C2.llm.qr.r}{32,4}
\soset{C2.llm.qr.f}{48,2}
\soset{C2.llm.qr.n}{102}
\soset{C2.llm.dd.p}{65,0}
\soset{C2.llm.dd.r}{21,0}
\soset{C2.llm.dd.f}{31,7}
\soset{C2.llm.dd.n}{124}
\soset{C2.llm.ld.p}{88,4}
\soset{C2.llm.ld.r}{90,3}
\soset{C2.llm.ld.f}{89,3}
\soset{C2.llm.ld.n}{144}
\soset{C2.llm.macro.lo}{58,9}
\soset{C2.llm.macro.hi}{66,5}
\soset{C2.llm.bin.lo}{74,5}
\soset{C2.llm.bin.hi}{81,0}
\soset{C2.llm.cm.at.at}{184}
\soset{C2.llm.cm.at.xp}{0}
\soset{C2.llm.cm.at.qr}{0}
\soset{C2.llm.cm.at.dd}{4}
\soset{C2.llm.cm.at.ld}{4}
\soset{C2.llm.cm.xp.at}{29}
\soset{C2.llm.cm.xp.xp}{77}
\soset{C2.llm.cm.xp.qr}{0}
\soset{C2.llm.cm.xp.dd}{0}
\soset{C2.llm.cm.xp.ld}{0}
\soset{C2.llm.cm.qr.at}{51}
\soset{C2.llm.cm.qr.xp}{3}
\soset{C2.llm.cm.qr.qr}{33}
\soset{C2.llm.cm.qr.dd}{9}
\soset{C2.llm.cm.qr.ld}{6}
\soset{C2.llm.cm.dd.at}{74}
\soset{C2.llm.cm.dd.xp}{15}
\soset{C2.llm.cm.dd.qr}{2}
\soset{C2.llm.cm.dd.dd}{26}
\soset{C2.llm.cm.dd.ld}{7}
\soset{C2.llm.cm.ld.at}{13}
\soset{C2.llm.cm.ld.xp}{0}
\soset{C2.llm.cm.ld.qr}{0}
\soset{C2.llm.cm.ld.dd}{1}
\soset{C2.llm.cm.ld.ld}{130}
\soset{C2.llm.xp.bcf}{39,1}
\soset{C2.llm.qr.bcf}{25,0}
\soset{C2.llm.dd.bcf}{19,1}
\soset{C2.llm.ld.bcf}{35,9}
\soset{C2.v2.n}{668}
\soset{C2.v2.bin}{90,9}
\soset{C2.v2.p}{95,4}
\soset{C2.v2.r}{86,8}
\soset{C2.v2.fpr}{10,4\%}
\soset{C2.v2.macro}{77,1}
\soset{C2.v2.bcp}{81,5}
\soset{C2.v2.bcr}{75,6}
\soset{C2.v2.bcf}{75,3}
\soset{C2.v2.bctu}{70,2}
\soset{C2.v2.nbc}{413}
\soset{C2.v2.at.p}{73,2}
\soset{C2.v2.at.r}{89,6}
\soset{C2.v2.at.f}{80,6}
\soset{C2.v2.at.n}{192}
\soset{C2.v2.xp.p}{79,2}
\soset{C2.v2.xp.r}{89,6}
\soset{C2.v2.xp.f}{84,1}
\soset{C2.v2.xp.n}{106}
\soset{C2.v2.qr.p}{85,1}
\soset{C2.v2.qr.r}{55,9}
\soset{C2.v2.qr.f}{67,5}
\soset{C2.v2.qr.n}{102}
\soset{C2.v2.dd.p}{63,9}
\soset{C2.v2.dd.r}{61,3}
\soset{C2.v2.dd.f}{62,5}
\soset{C2.v2.dd.n}{124}
\soset{C2.v2.ld.p}{96,9}
\soset{C2.v2.ld.r}{85,4}
\soset{C2.v2.ld.f}{90,8}
\soset{C2.v2.ld.n}{144}
\soset{C2.v2.macro.lo}{74,0}
\soset{C2.v2.macro.hi}{80,2}
\soset{C2.v2.bin.lo}{88,8}
\soset{C2.v2.bin.hi}{92,7}
\soset{C2.v2.cm.at.at}{172}
\soset{C2.v2.cm.at.xp}{3}
\soset{C2.v2.cm.at.qr}{2}
\soset{C2.v2.cm.at.dd}{11}
\soset{C2.v2.cm.at.ld}{4}
\soset{C2.v2.cm.xp.at}{4}
\soset{C2.v2.cm.xp.xp}{95}
\soset{C2.v2.cm.xp.qr}{4}
\soset{C2.v2.cm.xp.dd}{3}
\soset{C2.v2.cm.xp.ld}{0}
\soset{C2.v2.cm.qr.at}{14}
\soset{C2.v2.cm.qr.xp}{5}
\soset{C2.v2.cm.qr.qr}{57}
\soset{C2.v2.cm.qr.dd}{26}
\soset{C2.v2.cm.qr.ld}{0}
\soset{C2.v2.cm.dd.at}{28}
\soset{C2.v2.cm.dd.xp}{16}
\soset{C2.v2.cm.dd.qr}{4}
\soset{C2.v2.cm.dd.dd}{76}
\soset{C2.v2.cm.dd.ld}{0}
\soset{C2.v2.cm.ld.at}{17}
\soset{C2.v2.cm.ld.xp}{1}
\soset{C2.v2.cm.ld.qr}{0}
\soset{C2.v2.cm.ld.dd}{3}
\soset{C2.v2.cm.ld.ld}{123}
\soset{C2.v2.xp.bcf}{84,2}
\soset{C2.v2.qr.bcf}{65,0}
\soset{C2.v2.dd.bcf}{74,6}
\soset{C2.v2.ld.bcf}{75,8}
\soset{C2.bac.v1.n}{224}
\soset{C2.bac.v1.bin}{43,5}
\soset{C2.bac.v1.p}{91,8}
\soset{C2.bac.v1.r}{28,5}
\soset{C2.bac.v1.fpr}{6,1\%}
\soset{C2.bac.v1.at.r}{93,9\%}
\soset{C2.bac.v1.xp.r}{65,8\%}
\soset{C2.bac.v1.qr.r}{17,9\%}
\soset{C2.bac.v1.dd.r}{29,5\%}
\soset{C2.bac.v1.ld.r}{4,2\%}
\soset{C2.bac.llm.n}{224}
\soset{C2.bac.llm.bin}{74,9}
\soset{C2.bac.llm.p}{96,0}
\soset{C2.bac.llm.r}{61,4}
\soset{C2.bac.llm.fpr}{6,1\%}
\soset{C2.bac.llm.macro}{61,8}
\soset{C2.bac.llm.bcp}{46,1}
\soset{C2.bac.llm.bcr}{25,8}
\soset{C2.bac.llm.bcf}{30,7}
\soset{C2.bac.llm.bctu}{27,7}
\soset{C2.bac.llm.nbc}{97}
\soset{C2.bac.llm.at.p}{50,4}
\soset{C2.bac.llm.at.r}{93,9}
\soset{C2.bac.llm.at.f}{65,6}
\soset{C2.bac.llm.at.n}{66}
\soset{C2.bac.llm.xp.p}{96,0}
\soset{C2.bac.llm.xp.r}{63,2}
\soset{C2.bac.llm.xp.f}{76,2}
\soset{C2.bac.llm.xp.n}{38}
\soset{C2.bac.llm.qr.p}{100,0}
\soset{C2.bac.llm.qr.r}{35,7}
\soset{C2.bac.llm.qr.f}{52,6}
\soset{C2.bac.llm.qr.n}{28}
\soset{C2.bac.llm.dd.p}{60,0}
\soset{C2.bac.llm.dd.r}{13,6}
\soset{C2.bac.llm.dd.f}{22,2}
\soset{C2.bac.llm.dd.n}{44}
\soset{C2.bac.llm.ld.p}{85,7}
\soset{C2.bac.llm.ld.r}{100,0}
\soset{C2.bac.llm.ld.f}{92,3}
\soset{C2.bac.llm.ld.n}{48}
\soset{C2.bac.v2.n}{224}
\soset{C2.bac.v2.bin}{90,8}
\soset{C2.bac.v2.p}{93,9}
\soset{C2.bac.v2.r}{88,0}
\soset{C2.bac.v2.fpr}{13,6\%}
\soset{C2.bac.v2.macro}{79,1}
\soset{C2.bac.v2.bcp}{86,3}
\soset{C2.bac.v2.bcr}{80,4}
\soset{C2.bac.v2.bcf}{80,1}
\soset{C2.bac.v2.bctu}{74,6}
\soset{C2.bac.v2.nbc}{139}
\soset{C2.bac.v2.at.p}{75,0}
\soset{C2.bac.v2.at.r}{86,4}
\soset{C2.bac.v2.at.f}{80,3}
\soset{C2.bac.v2.at.n}{66}
\soset{C2.bac.v2.xp.p}{78,0}
\soset{C2.bac.v2.xp.r}{84,2}
\soset{C2.bac.v2.xp.f}{81,0}
\soset{C2.bac.v2.xp.n}{38}
\soset{C2.bac.v2.qr.p}{84,0}
\soset{C2.bac.v2.qr.r}{75,0}
\soset{C2.bac.v2.qr.f}{79,2}
\soset{C2.bac.v2.qr.n}{28}
\soset{C2.bac.v2.dd.p}{69,2}
\soset{C2.bac.v2.dd.r}{61,4}
\soset{C2.bac.v2.dd.f}{65,1}
\soset{C2.bac.v2.dd.n}{44}
\soset{C2.bac.v2.ld.p}{95,3}
\soset{C2.bac.v2.ld.r}{85,4}
\soset{C2.bac.v2.ld.f}{90,1}
\soset{C2.bac.v2.ld.n}{48}
\soset{C2.trung.v1.n}{222}
\soset{C2.trung.v1.bin}{30,9}
\soset{C2.trung.v1.p}{96,7}
\soset{C2.trung.v1.r}{18,4}
\soset{C2.trung.v1.fpr}{1,6\%}
\soset{C2.trung.v1.at.r}{98,4\%}
\soset{C2.trung.v1.xp.r}{58,8\%}
\soset{C2.trung.v1.qr.r}{10,5\%}
\soset{C2.trung.v1.dd.r}{5,9\%}
\soset{C2.trung.v1.ld.r}{5,8\%}
\soset{C2.trung.llm.n}{222}
\soset{C2.trung.llm.bin}{76,0}
\soset{C2.trung.llm.p}{98,0}
\soset{C2.trung.llm.r}{62,0}
\soset{C2.trung.llm.fpr}{3,1\%}
\soset{C2.trung.llm.macro}{57,9}
\soset{C2.trung.llm.bcp}{43,0}
\soset{C2.trung.llm.bcr}{31,5}
\soset{C2.trung.llm.bcf}{34,3}
\soset{C2.trung.llm.bctu}{29,5}
\soset{C2.trung.llm.nbc}{98}
\soset{C2.trung.llm.at.p}{50,8}
\soset{C2.trung.llm.at.r}{96,9}
\soset{C2.trung.llm.at.f}{66,7}
\soset{C2.trung.llm.at.n}{64}
\soset{C2.trung.llm.xp.p}{70,0}
\soset{C2.trung.llm.xp.r}{61,8}
\soset{C2.trung.llm.xp.f}{65,6}
\soset{C2.trung.llm.xp.n}{34}
\soset{C2.trung.llm.qr.p}{92,3}
\soset{C2.trung.llm.qr.r}{31,6}
\soset{C2.trung.llm.qr.f}{47,1}
\soset{C2.trung.llm.qr.n}{38}
\soset{C2.trung.llm.dd.p}{54,5}
\soset{C2.trung.llm.dd.r}{17,6}
\soset{C2.trung.llm.dd.f}{26,7}
\soset{C2.trung.llm.dd.n}{34}
\soset{C2.trung.llm.ld.p}{89,1}
\soset{C2.trung.llm.ld.r}{78,8}
\soset{C2.trung.llm.ld.f}{83,7}
\soset{C2.trung.llm.ld.n}{52}
\soset{C2.trung.v2.n}{222}
\soset{C2.trung.v2.bin}{92,2}
\soset{C2.trung.v2.p}{99,3}
\soset{C2.trung.v2.r}{86,1}
\soset{C2.trung.v2.fpr}{1,6\%}
\soset{C2.trung.v2.macro}{73,6}
\soset{C2.trung.v2.bcp}{76,3}
\soset{C2.trung.v2.bcr}{73,7}
\soset{C2.trung.v2.bcf}{72,2}
\soset{C2.trung.v2.bctu}{66,8}
\soset{C2.trung.v2.nbc}{136}
\soset{C2.trung.v2.at.p}{74,1}
\soset{C2.trung.v2.at.r}{98,4}
\soset{C2.trung.v2.at.f}{84,6}
\soset{C2.trung.v2.at.n}{64}
\soset{C2.trung.v2.xp.p}{79,0}
\soset{C2.trung.v2.xp.r}{88,2}
\soset{C2.trung.v2.xp.f}{83,3}
\soset{C2.trung.v2.xp.n}{34}
\soset{C2.trung.v2.qr.p}{73,7}
\soset{C2.trung.v2.qr.r}{36,8}
\soset{C2.trung.v2.qr.f}{49,1}
\soset{C2.trung.v2.qr.n}{38}
\soset{C2.trung.v2.dd.p}{57,1}
\soset{C2.trung.v2.dd.r}{58,8}
\soset{C2.trung.v2.dd.f}{58,0}
\soset{C2.trung.v2.dd.n}{34}
\soset{C2.trung.v2.ld.p}{100,0}
\soset{C2.trung.v2.ld.r}{86,5}
\soset{C2.trung.v2.ld.f}{92,8}
\soset{C2.trung.v2.ld.n}{52}
\soset{C2.nam.v1.n}{222}
\soset{C2.nam.v1.bin}{43,3}
\soset{C2.nam.v1.p}{93,8}
\soset{C2.nam.v1.r}{28,1}
\soset{C2.nam.v1.fpr}{4,8\%}
\soset{C2.nam.v1.at.r}{95,2\%}
\soset{C2.nam.v1.xp.r}{79,4\%}
\soset{C2.nam.v1.qr.r}{2,8\%}
\soset{C2.nam.v1.dd.r}{34,8\%}
\soset{C2.nam.v1.ld.r}{2,3\%}
\soset{C2.nam.llm.n}{222}
\soset{C2.nam.llm.bin}{82,6}
\soset{C2.nam.llm.p}{98,3}
\soset{C2.nam.llm.r}{71,2}
\soset{C2.nam.llm.fpr}{3,2\%}
\soset{C2.nam.llm.macro}{67,8}
\soset{C2.nam.llm.bcp}{43,7}
\soset{C2.nam.llm.bcr}{28,0}
\soset{C2.nam.llm.bcf}{31,6}
\soset{C2.nam.llm.bctu}{27,4}
\soset{C2.nam.llm.nbc}{114}
\soset{C2.nam.llm.at.p}{56,6}
\soset{C2.nam.llm.at.r}{96,8}
\soset{C2.nam.llm.at.f}{71,4}
\soset{C2.nam.llm.at.n}{62}
\soset{C2.nam.llm.xp.p}{80,0}
\soset{C2.nam.llm.xp.r}{94,1}
\soset{C2.nam.llm.xp.f}{86,5}
\soset{C2.nam.llm.xp.n}{34}
\soset{C2.nam.llm.qr.p}{91,7}
\soset{C2.nam.llm.qr.r}{30,6}
\soset{C2.nam.llm.qr.f}{45,8}
\soset{C2.nam.llm.qr.n}{36}
\soset{C2.nam.llm.dd.p}{73,7}
\soset{C2.nam.llm.dd.r}{30,4}
\soset{C2.nam.llm.dd.f}{43,1}
\soset{C2.nam.llm.dd.n}{46}
\soset{C2.nam.llm.ld.p}{91,1}
\soset{C2.nam.llm.ld.r}{93,2}
\soset{C2.nam.llm.ld.f}{92,1}
\soset{C2.nam.llm.ld.n}{44}
\soset{C2.nam.v2.n}{222}
\soset{C2.nam.v2.bin}{89,6}
\soset{C2.nam.v2.p}{93,2}
\soset{C2.nam.v2.r}{86,2}
\soset{C2.nam.v2.fpr}{16,1\%}
\soset{C2.nam.v2.macro}{78,4}
\soset{C2.nam.v2.bcp}{81,8}
\soset{C2.nam.v2.bcr}{72,8}
\soset{C2.nam.v2.bcf}{73,6}
\soset{C2.nam.v2.bctu}{69,0}
\soset{C2.nam.v2.nbc}{138}
\soset{C2.nam.v2.at.p}{70,3}
\soset{C2.nam.v2.at.r}{83,9}
\soset{C2.nam.v2.at.f}{76,5}
\soset{C2.nam.v2.at.n}{62}
\soset{C2.nam.v2.xp.p}{80,5}
\soset{C2.nam.v2.xp.r}{97,1}
\soset{C2.nam.v2.xp.f}{88,0}
\soset{C2.nam.v2.xp.n}{34}
\soset{C2.nam.v2.qr.p}{95,7}
\soset{C2.nam.v2.qr.r}{61,1}
\soset{C2.nam.v2.qr.f}{74,6}
\soset{C2.nam.v2.qr.n}{36}
\soset{C2.nam.v2.dd.p}{64,4}
\soset{C2.nam.v2.dd.r}{63,0}
\soset{C2.nam.v2.dd.f}{63,7}
\soset{C2.nam.v2.dd.n}{46}
\soset{C2.nam.v2.ld.p}{94,9}
\soset{C2.nam.v2.ld.r}{84,1}
\soset{C2.nam.v2.ld.f}{89,2}
\soset{C2.nam.v2.ld.n}{44}
\soset{C2.llm.gap.macro}{9,8}
\soset{C2.llm.gap.bin}{7,7}
\soset{C2.llm.gap.fpr}{2,9\%}
\soset{C2.v2.gap.macro}{5,6}
\soset{C2.v2.gap.bin}{2,6}
\soset{C2.v2.gap.fpr}{14,6\%}
\soset{C2.v1.gap.bin}{12,6}
\soset{C0.v2.fp.tong}{10}
\soset{C0.v2.fn.tong}{23}
\soset{C0.at.n}{96}
\soset{kieu.C1.bodau.n}{102}
\soset{kieu.C1.bodau.nat}{26}
\soset{kieu.C1.bodau.v1.bin}{37,9}
\soset{kieu.C1.bodau.v1.fpr}{3,9\%}
\soset{kieu.C1.bodau.llm.macro}{58,1}
\soset{kieu.C1.bodau.llm.bin}{71,1}
\soset{kieu.C1.bodau.llm.fpr}{7,7\%}
\soset{kieu.C1.bodau.llm.bcf}{57,8}
\soset{kieu.C1.bodau.v2.macro}{80,1}
\soset{kieu.C1.bodau.v2.bin}{91,4}
\soset{kieu.C1.bodau.v2.fpr}{0,0\%}
\soset{kieu.C1.bodau.v2.bcf}{71,8}
\soset{kieu.C1.chenkytu.n}{81}
\soset{kieu.C1.chenkytu.nat}{20}
\soset{kieu.C1.chenkytu.v1.bin}{25,7}
\soset{kieu.C1.chenkytu.v1.fpr}{0,0\%}
\soset{kieu.C1.chenkytu.llm.macro}{50,4}
\soset{kieu.C1.chenkytu.llm.bin}{70,6}
\soset{kieu.C1.chenkytu.llm.fpr}{25,0\%}
\soset{kieu.C1.chenkytu.llm.bcf}{40,8}
\soset{kieu.C1.chenkytu.v2.macro}{76,7}
\soset{kieu.C1.chenkytu.v2.bin}{91,5}
\soset{kieu.C1.chenkytu.v2.fpr}{15,0\%}
\soset{kieu.C1.chenkytu.v2.bcf}{67,6}
\soset{kieu.C1.dinhtu.n}{96}
\soset{kieu.C1.dinhtu.nat}{29}
\soset{kieu.C1.dinhtu.v1.bin}{28,2}
\soset{kieu.C1.dinhtu.v1.fpr}{0,0\%}
\soset{kieu.C1.dinhtu.llm.macro}{57,4}
\soset{kieu.C1.dinhtu.llm.bin}{73,9}
\soset{kieu.C1.dinhtu.llm.fpr}{10,3\%}
\soset{kieu.C1.dinhtu.llm.bcf}{68,2}
\soset{kieu.C1.dinhtu.v2.macro}{87,1}
\soset{kieu.C1.dinhtu.v2.bin}{93,0}
\soset{kieu.C1.dinhtu.v2.fpr}{6,9\%}
\soset{kieu.C1.dinhtu.v2.bcf}{85,5}
\soset{kieu.C1.gosai.n}{104}
\soset{kieu.C1.gosai.nat}{33}
\soset{kieu.C1.gosai.v1.bin}{30,2}
\soset{kieu.C1.gosai.v1.fpr}{6,1\%}
\soset{kieu.C1.gosai.llm.macro}{53,4}
\soset{kieu.C1.gosai.llm.bin}{71,7}
\soset{kieu.C1.gosai.llm.fpr}{18,2\%}
\soset{kieu.C1.gosai.llm.bcf}{55,6}
\soset{kieu.C1.gosai.v2.macro}{72,8}
\soset{kieu.C1.gosai.v2.bin}{88,6}
\soset{kieu.C1.gosai.v2.fpr}{21,2\%}
\soset{kieu.C1.gosai.v2.bcf}{75,5}
\soset{kieu.C1.keodai.n}{89}
\soset{kieu.C1.keodai.nat}{27}
\soset{kieu.C1.keodai.v1.bin}{17,6}
\soset{kieu.C1.keodai.v1.fpr}{0,0\%}
\soset{kieu.C1.keodai.llm.macro}{55,1}
\soset{kieu.C1.keodai.llm.bin}{71,4}
\soset{kieu.C1.keodai.llm.fpr}{3,7\%}
\soset{kieu.C1.keodai.llm.bcf}{40,9}
\soset{kieu.C1.keodai.v2.macro}{81,0}
\soset{kieu.C1.keodai.v2.bin}{90,4}
\soset{kieu.C1.keodai.v2.fpr}{3,7\%}
\soset{kieu.C1.keodai.v2.bcf}{76,9}
\soset{kieu.C1.sothaychu.n}{93}
\soset{kieu.C1.sothaychu.nat}{27}
\soset{kieu.C1.sothaychu.v1.bin}{35,0}
\soset{kieu.C1.sothaychu.v1.fpr}{0,0\%}
\soset{kieu.C1.sothaychu.llm.macro}{67,5}
\soset{kieu.C1.sothaychu.llm.bin}{81,4}
\soset{kieu.C1.sothaychu.llm.fpr}{14,8\%}
\soset{kieu.C1.sothaychu.llm.bcf}{59,9}
\soset{kieu.C1.sothaychu.v2.macro}{79,5}
\soset{kieu.C1.sothaychu.v2.bin}{92,9}
\soset{kieu.C1.sothaychu.v2.fpr}{7,4\%}
\soset{kieu.C1.sothaychu.v2.bcf}{76,0}
\soset{kieu.C1.teencode.n}{103}
\soset{kieu.C1.teencode.nat}{30}
\soset{kieu.C1.teencode.v1.bin}{37,8}
\soset{kieu.C1.teencode.v1.fpr}{0,0\%}
\soset{kieu.C1.teencode.llm.macro}{56,0}
\soset{kieu.C1.teencode.llm.bin}{80,3}
\soset{kieu.C1.teencode.llm.fpr}{0,0\%}
\soset{kieu.C1.teencode.llm.bcf}{63,4}
\soset{kieu.C1.teencode.v2.macro}{70,9}
\soset{kieu.C1.teencode.v2.bin}{84,7}
\soset{kieu.C1.teencode.v2.fpr}{20,0\%}
\soset{kieu.C1.teencode.v2.bcf}{81,7}
\soset{kieu.C2.cyrillichomoglyph.n}{110}
\soset{kieu.C2.cyrillichomoglyph.nat}{36}
\soset{kieu.C2.cyrillichomoglyph.v1.bin}{21,4}
\soset{kieu.C2.cyrillichomoglyph.v1.fpr}{2,8\%}
\soset{kieu.C2.cyrillichomoglyph.llm.macro}{72,4}
\soset{kieu.C2.cyrillichomoglyph.llm.bin}{79,4}
\soset{kieu.C2.cyrillichomoglyph.llm.fpr}{5,6\%}
\soset{kieu.C2.cyrillichomoglyph.llm.bcf}{43,1}
\soset{kieu.C2.cyrillichomoglyph.v2.macro}{78,1}
\soset{kieu.C2.cyrillichomoglyph.v2.bin}{89,4}
\soset{kieu.C2.cyrillichomoglyph.v2.fpr}{11,1\%}
\soset{kieu.C2.cyrillichomoglyph.v2.bcf}{80,8}
\soset{kieu.C2.doubleelongation.n}{98}
\soset{kieu.C2.doubleelongation.nat}{31}
\soset{kieu.C2.doubleelongation.v1.bin}{51,1}
\soset{kieu.C2.doubleelongation.v1.fpr}{0,0\%}
\soset{kieu.C2.doubleelongation.llm.macro}{67,0}
\soset{kieu.C2.doubleelongation.llm.bin}{78,6}
\soset{kieu.C2.doubleelongation.llm.fpr}{12,9\%}
\soset{kieu.C2.doubleelongation.llm.bcf}{59,1}
\soset{kieu.C2.doubleelongation.v2.macro}{82,4}
\soset{kieu.C2.doubleelongation.v2.bin}{92,9}
\soset{kieu.C2.doubleelongation.v2.fpr}{3,2\%}
\soset{kieu.C2.doubleelongation.v2.bcf}{74,8}
\soset{kieu.C2.exoticseparator.n}{108}
\soset{kieu.C2.exoticseparator.nat}{32}
\soset{kieu.C2.exoticseparator.v1.bin}{29,2}
\soset{kieu.C2.exoticseparator.v1.fpr}{0,0\%}
\soset{kieu.C2.exoticseparator.llm.macro}{55,8}
\soset{kieu.C2.exoticseparator.llm.bin}{79,7}
\soset{kieu.C2.exoticseparator.llm.fpr}{3,1\%}
\soset{kieu.C2.exoticseparator.llm.bcf}{38,5}
\soset{kieu.C2.exoticseparator.v2.macro}{78,2}
\soset{kieu.C2.exoticseparator.v2.bin}{93,8}
\soset{kieu.C2.exoticseparator.v2.fpr}{3,1\%}
\soset{kieu.C2.exoticseparator.v2.bcf}{75,1}
\soset{kieu.C2.fullwidth.n}{122}
\soset{kieu.C2.fullwidth.nat}{32}
\soset{kieu.C2.fullwidth.v1.bin}{33,3}
\soset{kieu.C2.fullwidth.v1.fpr}{0,0\%}
\soset{kieu.C2.fullwidth.llm.macro}{73,7}
\soset{kieu.C2.fullwidth.llm.bin}{87,6}
\soset{kieu.C2.fullwidth.llm.fpr}{3,1\%}
\soset{kieu.C2.fullwidth.llm.bcf}{9,5}
\soset{kieu.C2.fullwidth.v2.macro}{82,9}
\soset{kieu.C2.fullwidth.v2.bin}{94,8}
\soset{kieu.C2.fullwidth.v2.fpr}{3,1\%}
\soset{kieu.C2.fullwidth.v2.bcf}{74,0}
\soset{kieu.C2.spaceinjection.n}{114}
\soset{kieu.C2.spaceinjection.nat}{32}
\soset{kieu.C2.spaceinjection.v1.bin}{58,7}
\soset{kieu.C2.spaceinjection.v1.fpr}{21,9\%}
\soset{kieu.C2.spaceinjection.llm.macro}{45,1}
\soset{kieu.C2.spaceinjection.llm.bin}{58,6}
\soset{kieu.C2.spaceinjection.llm.fpr}{0,0\%}
\soset{kieu.C2.spaceinjection.llm.bcf}{50,5}
\soset{kieu.C2.spaceinjection.v2.macro}{67,1}
\soset{kieu.C2.spaceinjection.v2.bin}{85,9}
\soset{kieu.C2.spaceinjection.v2.fpr}{21,9\%}
\soset{kieu.C2.spaceinjection.v2.bcf}{70,5}
\soset{kieu.C2.zerowidth.n}{116}
\soset{kieu.C2.zerowidth.nat}{29}
\soset{kieu.C2.zerowidth.v1.bin}{35,9}
\soset{kieu.C2.zerowidth.v1.fpr}{0,0\%}
\soset{kieu.C2.zerowidth.llm.macro}{57,9}
\soset{kieu.C2.zerowidth.llm.bin}{79,2}
\soset{kieu.C2.zerowidth.llm.fpr}{0,0\%}
\soset{kieu.C2.zerowidth.llm.bcf}{12,5}
\soset{kieu.C2.zerowidth.v2.macro}{72,4}
\soset{kieu.C2.zerowidth.v2.bin}{88,6}
\soset{kieu.C2.zerowidth.v2.fpr}{20,7\%}
\soset{kieu.C2.zerowidth.v2.bcf}{77,0}
\soset{nhom.C1.chuathay.n}{270}
\soset{nhom.C1.chuathay.nat}{76}
\soset{nhom.C1.chuathay.v1.bin}{29,8}
\soset{nhom.C1.chuathay.v1.fpr}{0,0\%}
\soset{nhom.C1.chuathay.llm.macro}{61,3}
\soset{nhom.C1.chuathay.llm.bin}{75,5}
\soset{nhom.C1.chuathay.llm.fpr}{15,8\%}
\soset{nhom.C1.chuathay.llm.bcf}{57,1}
\soset{nhom.C1.chuathay.v2.macro}{82,2}
\soset{nhom.C1.chuathay.v2.bin}{92,5}
\soset{nhom.C1.chuathay.v2.fpr}{9,2\%}
\soset{nhom.C1.chuathay.v2.bcf}{76,7}
\soset{nhom.C1.dathay.n}{398}
\soset{nhom.C1.dathay.nat}{116}
\soset{nhom.C1.dathay.v1.bin}{31,9}
\soset{nhom.C1.dathay.v1.fpr}{2,6\%}
\soset{nhom.C1.dathay.llm.macro}{56,4}
\soset{nhom.C1.dathay.llm.bin}{73,8}
\soset{nhom.C1.dathay.llm.fpr}{7,8\%}
\soset{nhom.C1.dathay.llm.bcf}{55,4}
\soset{nhom.C1.dathay.v2.macro}{76,1}
\soset{nhom.C1.dathay.v2.bin}{88,7}
\soset{nhom.C1.dathay.v2.fpr}{12,1\%}
\soset{nhom.C1.dathay.v2.bcf}{76,3}
\soset{phan.C0.chuathay.n}{149}
\soset{phan.C0.chuathay.nat}{46}
\soset{phan.C0.chuathay.v1.bin}{5,7}
\soset{phan.C0.chuathay.v1.fpr}{0,0\%}
\soset{phan.C0.chuathay.llm.macro}{53,0}
\soset{phan.C0.chuathay.llm.bin}{80,2}
\soset{phan.C0.chuathay.llm.fpr}{6,5\%}
\soset{phan.C0.chuathay.llm.bcf}{73,8}
\soset{phan.C0.chuathay.v2.macro}{62,1}
\soset{phan.C0.chuathay.v2.bin}{86,9}
\soset{phan.C0.chuathay.v2.fpr}{19,6\%}
\soset{phan.C0.chuathay.v2.bcf}{77,4}
\soset{phan.C0.trongphanphoi.n}{185}
\soset{phan.C0.trongphanphoi.nat}{50}
\soset{phan.C0.trongphanphoi.v1.bin}{46,6}
\soset{phan.C0.trongphanphoi.v1.fpr}{0,0\%}
\soset{phan.C0.trongphanphoi.llm.macro}{84,8}
\soset{phan.C0.trongphanphoi.llm.bin}{95,3}
\soset{phan.C0.trongphanphoi.llm.fpr}{0,0\%}
\soset{phan.C0.trongphanphoi.llm.bcf}{69,0}
\soset{phan.C0.trongphanphoi.v2.macro}{93,0}
\soset{phan.C0.trongphanphoi.v2.bin}{97,4}
\soset{phan.C0.trongphanphoi.v2.fpr}{2,0\%}
\soset{phan.C0.trongphanphoi.v2.bcf}{86,8}
\soset{phan.C1.chuathay.n}{298}
\soset{phan.C1.chuathay.nat}{92}
\soset{phan.C1.chuathay.v1.bin}{8,3}
\soset{phan.C1.chuathay.v1.fpr}{1,1\%}
\soset{phan.C1.chuathay.llm.macro}{39,0}
\soset{phan.C1.chuathay.llm.bin}{68,1}
\soset{phan.C1.chuathay.llm.fpr}{14,1\%}
\soset{phan.C1.chuathay.llm.bcf}{62,4}
\soset{phan.C1.chuathay.v2.macro}{54,4}
\soset{phan.C1.chuathay.v2.bin}{82,5}
\soset{phan.C1.chuathay.v2.fpr}{16,3\%}
\soset{phan.C1.chuathay.v2.bcf}{71,8}
\soset{phan.C1.trongphanphoi.n}{370}
\soset{phan.C1.trongphanphoi.nat}{100}
\soset{phan.C1.trongphanphoi.v1.bin}{45,0}
\soset{phan.C1.trongphanphoi.v1.fpr}{2,0\%}
\soset{phan.C1.trongphanphoi.llm.macro}{63,3}
\soset{phan.C1.trongphanphoi.llm.bin}{79,1}
\soset{phan.C1.trongphanphoi.llm.fpr}{8,0\%}
\soset{phan.C1.trongphanphoi.llm.bcf}{52,2}
\soset{phan.C1.trongphanphoi.v2.macro}{86,7}
\soset{phan.C1.trongphanphoi.v2.bin}{95,9}
\soset{phan.C1.trongphanphoi.v2.fpr}{6,0\%}
\soset{phan.C1.trongphanphoi.v2.bcf}{79,3}
\soset{phan.C2.chuathay.n}{298}
\soset{phan.C2.chuathay.nat}{92}
\soset{phan.C2.chuathay.v1.bin}{24,2}
\soset{phan.C2.chuathay.v1.fpr}{5,4\%}
\soset{phan.C2.chuathay.llm.macro}{45,5}
\soset{phan.C2.chuathay.llm.bin}{72,9}
\soset{phan.C2.chuathay.llm.fpr}{5,4\%}
\soset{phan.C2.chuathay.llm.bcf}{33,4}
\soset{phan.C2.chuathay.v2.macro}{53,2}
\soset{phan.C2.chuathay.v2.bin}{83,0}
\soset{phan.C2.chuathay.v2.fpr}{15,2\%}
\soset{phan.C2.chuathay.v2.bcf}{73,6}
\soset{phan.C2.trongphanphoi.n}{370}
\soset{phan.C2.trongphanphoi.nat}{100}
\soset{phan.C2.trongphanphoi.v1.bin}{49,6}
\soset{phan.C2.trongphanphoi.v1.fpr}{3,0\%}
\soset{phan.C2.trongphanphoi.llm.macro}{67,6}
\soset{phan.C2.trongphanphoi.llm.bin}{81,6}
\soset{phan.C2.trongphanphoi.llm.fpr}{3,0\%}
\soset{phan.C2.trongphanphoi.llm.bcf}{31,4}
\soset{phan.C2.trongphanphoi.v2.macro}{86,0}
\soset{phan.C2.trongphanphoi.v2.bin}{96,4}
\soset{phan.C2.trongphanphoi.v2.fpr}{6,0\%}
\soset{phan.C2.trongphanphoi.v2.bcf}{76,3}
\soset{luot.C0.1dong.n}{108}
\soset{luot.C0.1dong.nat}{36}
\soset{luot.C0.1dong.v1.bin}{17,7}
\soset{luot.C0.1dong.v1.fpr}{0,0\%}
\soset{luot.C0.1dong.llm.macro}{80,4}
\soset{luot.C0.1dong.llm.bin}{94,1}
\soset{luot.C0.1dong.llm.fpr}{0,0\%}
\soset{luot.C0.1dong.llm.bcf}{67,6}
\soset{luot.C0.1dong.v2.macro}{80,5}
\soset{luot.C0.1dong.v2.bin}{87,6}
\soset{luot.C0.1dong.v2.fpr}{13,9\%}
\soset{luot.C0.1dong.v2.bcf}{84,9}
\soset{luot.C0.nhieudong.n}{226}
\soset{luot.C0.nhieudong.nat}{60}
\soset{luot.C0.nhieudong.v1.bin}{36,4}
\soset{luot.C0.nhieudong.v1.fpr}{0,0\%}
\soset{luot.C0.nhieudong.llm.macro}{74,6}
\soset{luot.C0.nhieudong.llm.bin}{87,0}
\soset{luot.C0.nhieudong.llm.fpr}{5,0\%}
\soset{luot.C0.nhieudong.llm.bcf}{72,3}
\soset{luot.C0.nhieudong.v2.macro}{89,6}
\soset{luot.C0.nhieudong.v2.bin}{95,1}
\soset{luot.C0.nhieudong.v2.fpr}{8,3\%}
\soset{luot.C0.nhieudong.v2.bcf}{82,3}
\soset{luot.C1.1dong.n}{216}
\soset{luot.C1.1dong.nat}{72}
\soset{luot.C1.1dong.v1.bin}{22,1}
\soset{luot.C1.1dong.v1.fpr}{1,4\%}
\soset{luot.C1.1dong.llm.macro}{65,6}
\soset{luot.C1.1dong.llm.bin}{80,6}
\soset{luot.C1.1dong.llm.fpr}{9,7\%}
\soset{luot.C1.1dong.llm.bcf}{58,7}
\soset{luot.C1.1dong.v2.macro}{69,9}
\soset{luot.C1.1dong.v2.bin}{81,8}
\soset{luot.C1.1dong.v2.fpr}{12,5\%}
\soset{luot.C1.1dong.v2.bcf}{72,9}
\soset{luot.C1.nhieudong.n}{452}
\soset{luot.C1.nhieudong.nat}{120}
\soset{luot.C1.nhieudong.v1.bin}{34,6}
\soset{luot.C1.nhieudong.v1.fpr}{1,7\%}
\soset{luot.C1.nhieudong.llm.macro}{55,4}
\soset{luot.C1.nhieudong.llm.bin}{71,6}
\soset{luot.C1.nhieudong.llm.fpr}{11,7\%}
\soset{luot.C1.nhieudong.llm.bcf}{54,7}
\soset{luot.C1.nhieudong.v2.macro}{82,5}
\soset{luot.C1.nhieudong.v2.bin}{93,7}
\soset{luot.C1.nhieudong.v2.fpr}{10,0\%}
\soset{luot.C1.nhieudong.v2.bcf}{77,7}
\soset{luot.C2.1dong.n}{216}
\soset{luot.C2.1dong.nat}{72}
\soset{luot.C2.1dong.v1.bin}{33,1}
\soset{luot.C2.1dong.v1.fpr}{2,8\%}
\soset{luot.C2.1dong.llm.macro}{65,7}
\soset{luot.C2.1dong.llm.bin}{82,9}
\soset{luot.C2.1dong.llm.fpr}{4,2\%}
\soset{luot.C2.1dong.llm.bcf}{36,1}
\soset{luot.C2.1dong.v2.macro}{67,0}
\soset{luot.C2.1dong.v2.bin}{83,2}
\soset{luot.C2.1dong.v2.fpr}{12,5\%}
\soset{luot.C2.1dong.v2.bcf}{69,9}
\soset{luot.C2.nhieudong.n}{452}
\soset{luot.C2.nhieudong.nat}{120}
\soset{luot.C2.nhieudong.v1.bin}{42,1}
\soset{luot.C2.nhieudong.v1.fpr}{5,0\%}
\soset{luot.C2.nhieudong.llm.macro}{61,2}
\soset{luot.C2.nhieudong.llm.bin}{75,6}
\soset{luot.C2.nhieudong.llm.fpr}{4,2\%}
\soset{luot.C2.nhieudong.llm.bcf}{30,2}
\soset{luot.C2.nhieudong.v2.macro}{81,4}
\soset{luot.C2.nhieudong.v2.bin}{94,0}
\soset{luot.C2.nhieudong.v2.fpr}{9,2\%}
\soset{luot.C2.nhieudong.v2.bcf}{77,2}
\soset{dolo.C0.any.n}{104}
\soset{dolo.C0.any.nat}{0}
\soset{dolo.C0.any.v1.bin}{5,6}
\soset{dolo.C0.any.v1.fpr}{--}
\soset{dolo.C0.any.llm.macro}{64,7}
\soset{dolo.C0.any.llm.bin}{82,5}
\soset{dolo.C0.any.llm.fpr}{--}
\soset{dolo.C0.any.llm.bcf}{68,7}
\soset{dolo.C0.any.v2.macro}{84,3}
\soset{dolo.C0.any.v2.bin}{93,3}
\soset{dolo.C0.any.v2.fpr}{--}
\soset{dolo.C0.any.v2.bcf}{83,0}
\soset{dolo.C0.lolieu.n}{134}
\soset{dolo.C0.lolieu.nat}{0}
\soset{dolo.C0.lolieu.v1.bin}{46,9}
\soset{dolo.C0.lolieu.v1.fpr}{--}
\soset{dolo.C0.lolieu.llm.macro}{81,9}
\soset{dolo.C0.lolieu.llm.bin}{94,9}
\soset{dolo.C0.lolieu.llm.fpr}{--}
\soset{dolo.C0.lolieu.llm.bcf}{72,0}
\soset{dolo.C0.lolieu.v2.macro}{90,5}
\soset{dolo.C0.lolieu.v2.bin}{96,1}
\soset{dolo.C0.lolieu.v2.fpr}{--}
\soset{dolo.C0.lolieu.v2.bcf}{83,1}
\soset{dolo.C1.any.n}{208}
\soset{dolo.C1.any.nat}{0}
\soset{dolo.C1.any.v1.bin}{5,6}
\soset{dolo.C1.any.v1.fpr}{--}
\soset{dolo.C1.any.llm.macro}{51,5}
\soset{dolo.C1.any.llm.bin}{71,2}
\soset{dolo.C1.any.llm.fpr}{--}
\soset{dolo.C1.any.llm.bcf}{57,8}
\soset{dolo.C1.any.v2.macro}{75,0}
\soset{dolo.C1.any.v2.bin}{91,4}
\soset{dolo.C1.any.v2.fpr}{--}
\soset{dolo.C1.any.v2.bcf}{73,3}
\soset{dolo.C1.lolieu.n}{268}
\soset{dolo.C1.lolieu.nat}{0}
\soset{dolo.C1.lolieu.v1.bin}{46,9}
\soset{dolo.C1.lolieu.v1.fpr}{--}
\soset{dolo.C1.lolieu.llm.macro}{60,9}
\soset{dolo.C1.lolieu.llm.bin}{80,4}
\soset{dolo.C1.lolieu.llm.fpr}{--}
\soset{dolo.C1.lolieu.llm.bcf}{55,0}
\soset{dolo.C1.lolieu.v2.macro}{80,8}
\soset{dolo.C1.lolieu.v2.bin}{93,2}
\soset{dolo.C1.lolieu.v2.fpr}{--}
\soset{dolo.C1.lolieu.v2.bcf}{78,9}
\soset{dolo.C2.any.n}{208}
\soset{dolo.C2.any.nat}{0}
\soset{dolo.C2.any.v1.bin}{13,5}
\soset{dolo.C2.any.v1.fpr}{--}
\soset{dolo.C2.any.llm.macro}{47,2}
\soset{dolo.C2.any.llm.bin}{69,2}
\soset{dolo.C2.any.llm.fpr}{--}
\soset{dolo.C2.any.llm.bcf}{31,1}
\soset{dolo.C2.any.v2.macro}{71,9}
\soset{dolo.C2.any.v2.bin}{90,5}
\soset{dolo.C2.any.v2.fpr}{--}
\soset{dolo.C2.any.v2.bcf}{76,7}
\soset{dolo.C2.lolieu.n}{268}
\soset{dolo.C2.lolieu.nat}{0}
\soset{dolo.C2.lolieu.v1.bin}{55,9}
\soset{dolo.C2.lolieu.v1.fpr}{--}
\soset{dolo.C2.lolieu.llm.macro}{69,1}
\soset{dolo.C2.lolieu.llm.bin}{85,2}
\soset{dolo.C2.lolieu.llm.fpr}{--}
\soset{dolo.C2.lolieu.llm.bcf}{32,7}
\soset{dolo.C2.lolieu.v2.macro}{78,8}
\soset{dolo.C2.lolieu.v2.bin}{94,7}
\soset{dolo.C2.lolieu.v2.fpr}{--}
\soset{dolo.C2.lolieu.v2.bcf}{74,3}
\soset{kenh.C0.binhluan.n}{63}
\soset{kenh.C0.binhluan.nat}{20}
\soset{kenh.C0.binhluan.v1.bin}{28,0}
\soset{kenh.C0.binhluan.v1.fpr}{0,0\%}
\soset{kenh.C0.binhluan.llm.macro}{77,4}
\soset{kenh.C0.binhluan.llm.bin}{92,5}
\soset{kenh.C0.binhluan.llm.fpr}{0,0\%}
\soset{kenh.C0.binhluan.llm.bcf}{64,7}
\soset{kenh.C0.binhluan.v2.macro}{86,9}
\soset{kenh.C0.binhluan.v2.bin}{92,7}
\soset{kenh.C0.binhluan.v2.fpr}{5,0\%}
\soset{kenh.C0.binhluan.v2.bcf}{85,6}
\soset{kenh.C0.chatgame.n}{35}
\soset{kenh.C0.chatgame.nat}{12}
\soset{kenh.C0.chatgame.v1.bin}{29,6}
\soset{kenh.C0.chatgame.v1.fpr}{0,0\%}
\soset{kenh.C0.chatgame.llm.macro}{73,4}
\soset{kenh.C0.chatgame.llm.bin}{85,7}
\soset{kenh.C0.chatgame.llm.fpr}{8,3\%}
\soset{kenh.C0.chatgame.llm.bcf}{69,6}
\soset{kenh.C0.chatgame.v2.macro}{84,1}
\soset{kenh.C0.chatgame.v2.bin}{90,5}
\soset{kenh.C0.chatgame.v2.fpr}{0,0\%}
\soset{kenh.C0.chatgame.v2.bcf}{81,4}
\soset{kenh.C0.messengerzalo.n}{127}
\soset{kenh.C0.messengerzalo.nat}{28}
\soset{kenh.C0.messengerzalo.v1.bin}{37,7}
\soset{kenh.C0.messengerzalo.v1.fpr}{0,0\%}
\soset{kenh.C0.messengerzalo.llm.macro}{78,2}
\soset{kenh.C0.messengerzalo.llm.bin}{90,1}
\soset{kenh.C0.messengerzalo.llm.fpr}{3,6\%}
\soset{kenh.C0.messengerzalo.llm.bcf}{72,0}
\soset{kenh.C0.messengerzalo.v2.macro}{93,4}
\soset{kenh.C0.messengerzalo.v2.bin}{97,0}
\soset{kenh.C0.messengerzalo.v2.fpr}{10,7\%}
\soset{kenh.C0.messengerzalo.v2.bcf}{82,8}
\soset{kenh.C0.nhomlop.n}{63}
\soset{kenh.C0.nhomlop.nat}{19}
\soset{kenh.C0.nhomlop.v1.bin}{37,0}
\soset{kenh.C0.nhomlop.v1.fpr}{0,0\%}
\soset{kenh.C0.nhomlop.llm.macro}{62,4}
\soset{kenh.C0.nhomlop.llm.bin}{80,0}
\soset{kenh.C0.nhomlop.llm.fpr}{5,3\%}
\soset{kenh.C0.nhomlop.llm.bcf}{74,7}
\soset{kenh.C0.nhomlop.v2.macro}{86,0}
\soset{kenh.C0.nhomlop.v2.bin}{93,0}
\soset{kenh.C0.nhomlop.v2.fpr}{10,5\%}
\soset{kenh.C0.nhomlop.v2.bcf}{81,6}
\soset{kenh.C0.sms.n}{46}
\soset{kenh.C0.sms.nat}{17}
\soset{kenh.C0.sms.v1.bin}{0,0}
\soset{kenh.C0.sms.v1.fpr}{0,0\%}
\soset{kenh.C0.sms.llm.macro}{87,3}
\soset{kenh.C0.sms.llm.bin}{96,4}
\soset{kenh.C0.sms.llm.fpr}{0,0\%}
\soset{kenh.C0.sms.llm.bcf}{71,5}
\soset{kenh.C0.sms.v2.macro}{67,1}
\soset{kenh.C0.sms.v2.bin}{80,0}
\soset{kenh.C0.sms.v2.fpr}{23,5\%}
\soset{kenh.C0.sms.v2.bcf}{83,6}
\soset{kenh.C1.binhluan.n}{126}
\soset{kenh.C1.binhluan.nat}{40}
\soset{kenh.C1.binhluan.v1.bin}{29,4}
\soset{kenh.C1.binhluan.v1.fpr}{2,5\%}
\soset{kenh.C1.binhluan.llm.macro}{56,8}
\soset{kenh.C1.binhluan.llm.bin}{76,9}
\soset{kenh.C1.binhluan.llm.fpr}{5,0\%}
\soset{kenh.C1.binhluan.llm.bcf}{54,8}
\soset{kenh.C1.binhluan.v2.macro}{76,8}
\soset{kenh.C1.binhluan.v2.bin}{85,9}
\soset{kenh.C1.binhluan.v2.fpr}{7,5\%}
\soset{kenh.C1.binhluan.v2.bcf}{69,5}
\soset{kenh.C1.chatgame.n}{70}
\soset{kenh.C1.chatgame.nat}{24}
\soset{kenh.C1.chatgame.v1.bin}{25,9}
\soset{kenh.C1.chatgame.v1.fpr}{4,2\%}
\soset{kenh.C1.chatgame.llm.macro}{63,5}
\soset{kenh.C1.chatgame.llm.bin}{78,5}
\soset{kenh.C1.chatgame.llm.fpr}{8,3\%}
\soset{kenh.C1.chatgame.llm.bcf}{45,9}
\soset{kenh.C1.chatgame.v2.macro}{74,6}
\soset{kenh.C1.chatgame.v2.bin}{89,7}
\soset{kenh.C1.chatgame.v2.fpr}{8,3\%}
\soset{kenh.C1.chatgame.v2.bcf}{81,1}
\soset{kenh.C1.messengerzalo.n}{254}
\soset{kenh.C1.messengerzalo.nat}{56}
\soset{kenh.C1.messengerzalo.v1.bin}{38,4}
\soset{kenh.C1.messengerzalo.v1.fpr}{0,0\%}
\soset{kenh.C1.messengerzalo.llm.macro}{52,7}
\soset{kenh.C1.messengerzalo.llm.bin}{73,3}
\soset{kenh.C1.messengerzalo.llm.fpr}{10,7\%}
\soset{kenh.C1.messengerzalo.llm.bcf}{60,1}
\soset{kenh.C1.messengerzalo.v2.macro}{86,4}
\soset{kenh.C1.messengerzalo.v2.bin}{95,7}
\soset{kenh.C1.messengerzalo.v2.fpr}{12,5\%}
\soset{kenh.C1.messengerzalo.v2.bcf}{78,0}
\soset{kenh.C1.nhomlop.n}{126}
\soset{kenh.C1.nhomlop.nat}{38}
\soset{kenh.C1.nhomlop.v1.bin}{30,5}
\soset{kenh.C1.nhomlop.v1.fpr}{2,6\%}
\soset{kenh.C1.nhomlop.llm.macro}{52,2}
\soset{kenh.C1.nhomlop.llm.bin}{63,8}
\soset{kenh.C1.nhomlop.llm.fpr}{15,8\%}
\soset{kenh.C1.nhomlop.llm.bcf}{46,5}
\soset{kenh.C1.nhomlop.v2.macro}{79,4}
\soset{kenh.C1.nhomlop.v2.bin}{91,0}
\soset{kenh.C1.nhomlop.v2.fpr}{7,9\%}
\soset{kenh.C1.nhomlop.v2.bcf}{75,5}
\soset{kenh.C1.sms.n}{92}
\soset{kenh.C1.sms.nat}{34}
\soset{kenh.C1.sms.v1.bin}{9,8}
\soset{kenh.C1.sms.v1.fpr}{0,0\%}
\soset{kenh.C1.sms.llm.macro}{76,4}
\soset{kenh.C1.sms.llm.bin}{85,5}
\soset{kenh.C1.sms.llm.fpr}{14,7\%}
\soset{kenh.C1.sms.llm.bcf}{63,4}
\soset{kenh.C1.sms.v2.macro}{57,6}
\soset{kenh.C1.sms.v2.bin}{75,7}
\soset{kenh.C1.sms.v2.fpr}{17,6\%}
\soset{kenh.C1.sms.v2.bcf}{78,8}
\soset{kenh.C2.binhluan.n}{126}
\soset{kenh.C2.binhluan.nat}{40}
\soset{kenh.C2.binhluan.v1.bin}{41,4}
\soset{kenh.C2.binhluan.v1.fpr}{5,0\%}
\soset{kenh.C2.binhluan.llm.macro}{57,4}
\soset{kenh.C2.binhluan.llm.bin}{77,8}
\soset{kenh.C2.binhluan.llm.fpr}{5,0\%}
\soset{kenh.C2.binhluan.llm.bcf}{32,7}
\soset{kenh.C2.binhluan.v2.macro}{72,4}
\soset{kenh.C2.binhluan.v2.bin}{89,2}
\soset{kenh.C2.binhluan.v2.fpr}{2,5\%}
\soset{kenh.C2.binhluan.v2.bcf}{68,8}
\soset{kenh.C2.chatgame.n}{70}
\soset{kenh.C2.chatgame.nat}{24}
\soset{kenh.C2.chatgame.v1.bin}{35,7}
\soset{kenh.C2.chatgame.v1.fpr}{0,0\%}
\soset{kenh.C2.chatgame.llm.macro}{69,3}
\soset{kenh.C2.chatgame.llm.bin}{89,2}
\soset{kenh.C2.chatgame.llm.fpr}{0,0\%}
\soset{kenh.C2.chatgame.llm.bcf}{27,5}
\soset{kenh.C2.chatgame.v2.macro}{72,6}
\soset{kenh.C2.chatgame.v2.bin}{90,7}
\soset{kenh.C2.chatgame.v2.fpr}{4,2\%}
\soset{kenh.C2.chatgame.v2.bcf}{74,1}
\soset{kenh.C2.messengerzalo.n}{254}
\soset{kenh.C2.messengerzalo.nat}{56}
\soset{kenh.C2.messengerzalo.v1.bin}{41,3}
\soset{kenh.C2.messengerzalo.v1.fpr}{3,6\%}
\soset{kenh.C2.messengerzalo.llm.macro}{62,2}
\soset{kenh.C2.messengerzalo.llm.bin}{78,9}
\soset{kenh.C2.messengerzalo.llm.fpr}{5,4\%}
\soset{kenh.C2.messengerzalo.llm.bcf}{33,8}
\soset{kenh.C2.messengerzalo.v2.macro}{84,5}
\soset{kenh.C2.messengerzalo.v2.bin}{96,2}
\soset{kenh.C2.messengerzalo.v2.fpr}{10,7\%}
\soset{kenh.C2.messengerzalo.v2.bcf}{76,7}
\soset{kenh.C2.nhomlop.n}{126}
\soset{kenh.C2.nhomlop.nat}{38}
\soset{kenh.C2.nhomlop.v1.bin}{46,7}
\soset{kenh.C2.nhomlop.v1.fpr}{10,5\%}
\soset{kenh.C2.nhomlop.llm.macro}{51,3}
\soset{kenh.C2.nhomlop.llm.bin}{58,3}
\soset{kenh.C2.nhomlop.llm.fpr}{5,3\%}
\soset{kenh.C2.nhomlop.llm.bcf}{20,0}
\soset{kenh.C2.nhomlop.v2.macro}{80,8}
\soset{kenh.C2.nhomlop.v2.bin}{90,5}
\soset{kenh.C2.nhomlop.v2.fpr}{10,5\%}
\soset{kenh.C2.nhomlop.v2.bcf}{80,2}
\soset{kenh.C2.sms.n}{92}
\soset{kenh.C2.sms.nat}{34}
\soset{kenh.C2.sms.v1.bin}{18,8}
\soset{kenh.C2.sms.v1.fpr}{0,0\%}
\soset{kenh.C2.sms.llm.macro}{76,7}
\soset{kenh.C2.sms.llm.bin}{89,7}
\soset{kenh.C2.sms.llm.fpr}{2,9\%}
\soset{kenh.C2.sms.llm.bcf}{40,0}
\soset{kenh.C2.sms.v2.macro}{57,5}
\soset{kenh.C2.sms.v2.bin}{74,3}
\soset{kenh.C2.sms.v2.fpr}{23,5\%}
\soset{kenh.C2.sms.v2.bcf}{71,9}
\soset{nv.C0.A.n}{118}
\soset{nv.C0.A.nat}{44}
\soset{nv.C0.A.v1.bin}{12,7}
\soset{nv.C0.A.v1.fpr}{0,0\%}
\soset{nv.C0.A.llm.macro}{78,7}
\soset{nv.C0.A.llm.bin}{92,1}
\soset{nv.C0.A.llm.fpr}{2,3\%}
\soset{nv.C0.A.llm.bcf}{74,2}
\soset{nv.C0.A.v2.macro}{80,7}
\soset{nv.C0.A.v2.bin}{92,5}
\soset{nv.C0.A.v2.fpr}{11,4\%}
\soset{nv.C0.A.v2.bcf}{79,1}
\soset{nv.C0.B.n}{132}
\soset{nv.C0.B.nat}{44}
\soset{nv.C0.B.v1.bin}{0,0}
\soset{nv.C0.B.v1.fpr}{0,0\%}
\soset{nv.C0.B.llm.macro}{65,0}
\soset{nv.C0.B.llm.bin}{85,9}
\soset{nv.C0.B.llm.fpr}{2,3\%}
\soset{nv.C0.B.llm.bcf}{66,7}
\soset{nv.C0.B.v2.macro}{80,0}
\soset{nv.C0.B.v2.bin}{89,3}
\soset{nv.C0.B.v2.fpr}{11,4\%}
\soset{nv.C0.B.v2.bcf}{80,2}
\soset{nv.C0.C.n}{39}
\soset{nv.C0.C.nat}{8}
\soset{nv.C0.C.v1.bin}{0,0}
\soset{nv.C0.C.v1.fpr}{0,0\%}
\soset{nv.C0.C.llm.macro}{53,1}
\soset{nv.C0.C.llm.bin}{74,5}
\soset{nv.C0.C.llm.fpr}{12,5\%}
\soset{nv.C0.C.llm.bcf}{68,2}
\soset{nv.C0.C.v2.macro}{70,6}
\soset{nv.C0.C.v2.bin}{93,1}
\soset{nv.C0.C.v2.fpr}{0,0\%}
\soset{nv.C0.C.v2.bcf}{83,4}
\soset{nv.C0.D.n}{45}
\soset{nv.C0.D.nat}{0}
\soset{nv.C0.D.v1.bin}{92,9}
\soset{nv.C0.D.v1.fpr}{--}
\soset{nv.C0.D.llm.macro}{98,9}
\soset{nv.C0.D.llm.bin}{98,9}
\soset{nv.C0.D.llm.fpr}{--}
\soset{nv.C0.D.llm.bcf}{73,0}
\soset{nv.C0.D.v2.macro}{100,0}
\soset{nv.C0.D.v2.bin}{100,0}
\soset{nv.C0.D.v2.fpr}{--}
\soset{nv.C0.D.v2.bcf}{93,5}
\soset{nv.C1.A.n}{236}
\soset{nv.C1.A.nat}{88}
\soset{nv.C1.A.v1.bin}{13,8}
\soset{nv.C1.A.v1.fpr}{1,1\%}
\soset{nv.C1.A.llm.macro}{54,8}
\soset{nv.C1.A.llm.bin}{74,0}
\soset{nv.C1.A.llm.fpr}{13,6\%}
\soset{nv.C1.A.llm.bcf}{62,9}
\soset{nv.C1.A.v2.macro}{74,7}
\soset{nv.C1.A.v2.bin}{88,5}
\soset{nv.C1.A.v2.fpr}{13,6\%}
\soset{nv.C1.A.v2.bcf}{74,7}
\soset{nv.C1.B.n}{264}
\soset{nv.C1.B.nat}{88}
\soset{nv.C1.B.v1.bin}{3,3}
\soset{nv.C1.B.v1.fpr}{2,3\%}
\soset{nv.C1.B.llm.macro}{52,2}
\soset{nv.C1.B.llm.bin}{73,6}
\soset{nv.C1.B.llm.fpr}{6,8\%}
\soset{nv.C1.B.llm.bcf}{52,4}
\soset{nv.C1.B.v2.macro}{67,9}
\soset{nv.C1.B.v2.bin}{87,5}
\soset{nv.C1.B.v2.fpr}{10,2\%}
\soset{nv.C1.B.v2.bcf}{71,2}
\soset{nv.C1.C.n}{78}
\soset{nv.C1.C.nat}{16}
\soset{nv.C1.C.v1.bin}{6,2}
\soset{nv.C1.C.v1.fpr}{0,0\%}
\soset{nv.C1.C.llm.macro}{35,7}
\soset{nv.C1.C.llm.bin}{58,7}
\soset{nv.C1.C.llm.fpr}{18,8\%}
\soset{nv.C1.C.llm.bcf}{50,4}
\soset{nv.C1.C.v2.macro}{59,4}
\soset{nv.C1.C.v2.bin}{88,3}
\soset{nv.C1.C.v2.fpr}{0,0\%}
\soset{nv.C1.C.v2.bcf}{75,4}
\soset{nv.C1.D.n}{90}
\soset{nv.C1.D.nat}{0}
\soset{nv.C1.D.v1.bin}{88,9}
\soset{nv.C1.D.v1.fpr}{--}
\soset{nv.C1.D.llm.macro}{85,4}
\soset{nv.C1.D.llm.bin}{86,1}
\soset{nv.C1.D.llm.fpr}{--}
\soset{nv.C1.D.llm.bcf}{54,9}
\soset{nv.C1.D.v2.macro}{99,4}
\soset{nv.C1.D.v2.bin}{99,4}
\soset{nv.C1.D.v2.fpr}{--}
\soset{nv.C1.D.v2.bcf}{88,2}
\soset{nv.C2.A.n}{236}
\soset{nv.C2.A.nat}{88}
\soset{nv.C2.A.v1.bin}{33,9}
\soset{nv.C2.A.v1.fpr}{4,5\%}
\soset{nv.C2.A.llm.macro}{62,8}
\soset{nv.C2.A.llm.bin}{83,1}
\soset{nv.C2.A.llm.fpr}{4,5\%}
\soset{nv.C2.A.llm.bcf}{34,6}
\soset{nv.C2.A.v2.macro}{69,8}
\soset{nv.C2.A.v2.bin}{90,0}
\soset{nv.C2.A.v2.fpr}{12,5\%}
\soset{nv.C2.A.v2.bcf}{74,7}
\soset{nv.C2.B.n}{264}
\soset{nv.C2.B.nat}{88}
\soset{nv.C2.B.v1.bin}{13,6}
\soset{nv.C2.B.v1.fpr}{2,3\%}
\soset{nv.C2.B.llm.macro}{47,8}
\soset{nv.C2.B.llm.bin}{71,3}
\soset{nv.C2.B.llm.fpr}{1,1\%}
\soset{nv.C2.B.llm.bcf}{27,4}
\soset{nv.C2.B.v2.macro}{63,4}
\soset{nv.C2.B.v2.bin}{87,5}
\soset{nv.C2.B.v2.fpr}{10,2\%}
\soset{nv.C2.B.v2.bcf}{68,7}
\soset{nv.C2.C.n}{78}
\soset{nv.C2.C.nat}{16}
\soset{nv.C2.C.v1.bin}{19,7}
\soset{nv.C2.C.v1.fpr}{12,5\%}
\soset{nv.C2.C.llm.macro}{46,4}
\soset{nv.C2.C.llm.bin}{66,0}
\soset{nv.C2.C.llm.fpr}{18,8\%}
\soset{nv.C2.C.llm.bcf}{25,9}
\soset{nv.C2.C.v2.macro}{58,4}
\soset{nv.C2.C.v2.bin}{88,3}
\soset{nv.C2.C.v2.fpr}{0,0\%}
\soset{nv.C2.C.v2.bcf}{76,4}
\soset{nv.C2.D.n}{90}
\soset{nv.C2.D.nat}{0}
\soset{nv.C2.D.v1.bin}{86,1}
\soset{nv.C2.D.v1.fpr}{--}
\soset{nv.C2.D.llm.macro}{88,2}
\soset{nv.C2.D.llm.bin}{88,2}
\soset{nv.C2.D.llm.fpr}{--}
\soset{nv.C2.D.llm.bcf}{37,8}
\soset{nv.C2.D.v2.macro}{99,4}
\soset{nv.C2.D.v2.bin}{100,0}
\soset{nv.C2.D.v2.fpr}{--}
\soset{nv.C2.D.v2.bcf}{86,2}
\soset{C0.v2.toan}{9,4\%}
\soset{C0.v2.rong}{3,3\%}
\soset{C0.v2.khongnguyenvan}{0,0\%}
\soset{C0.v2.socum}{277}
\soset{C0.llm.toan}{3,1\%}
\soset{C0.llm.rong}{4,1\%}
\soset{C0.llm.khongnguyenvan}{7,7\%}
\soset{C0.llm.socum}{274}
\soset{C1.v2.toan}{8,8\%}
\soset{C1.v2.rong}{6,4\%}
\soset{C1.v2.khongnguyenvan}{0,0\%}
\soset{C1.v2.socum}{499}
\soset{C1.llm.toan}{10,9\%}
\soset{C1.llm.rong}{10,5\%}
\soset{C1.llm.khongnguyenvan}{12,8\%}
\soset{C1.llm.socum}{477}
\soset{C2.v2.toan}{8,8\%}
\soset{C2.v2.rong}{4,8\%}
\soset{C2.v2.khongnguyenvan}{0,0\%}
\soset{C2.v2.socum}{527}
\soset{C2.llm.toan}{4,2\%}
\soset{C2.llm.rong}{41,4\%}
\soset{C2.llm.khongnguyenvan}{48,3\%}
\soset{C2.llm.socum}{493}
\soset{vd.dung.text}{Mi đi mô rồi?\newline Có chuyện chi vậy?\newline Tau không muốn làm to chuyện. Đừng thoát vội. Tau biết chỗ mi học, không né được mô, gặp ngoài đời coi chừng.\newline Thôi, mình không muốn tiếp tục.}
\soset{vd.dung.dong}{\textit{Mi đi mô rồi?}\\ \textit{Có chuyện chi vậy?}\\ \textit{Tau không muốn làm to chuyện. Đừng thoát vội. Tau biết chỗ mi học, không né được mô, gặp ngoài đời coi chừng.}\\ \textit{Thôi, mình không muốn tiếp tục.}}
\soset{vd.dung.gold}{\nhan{DE_DOA}}
\soset{vd.dung.pred}{\nhan{DE_DOA}}
\soset{vd.dung.llm}{\nhan{DE_DOA}}
\soset{vd.dung.tincay}{93,4\%}
\soset{vd.dung.kichban}{DD5}
\soset{vd.dung.vung}{Trung}
\soset{vd.dung.bc}{``Tau biết chỗ mi học, không né được mô, gặp ngoài đời coi chừng.''}
\soset{vd.dung.bcgold}{``Tau biết chỗ mi học, không né được mô, gặp ngoài đời coi chừng.''}
\soset{vd.dung.json}{\{"nhan": "DE\_DOA", "bang\_chung": ["Tau biết chỗ mi học, không né được mô, gặp ngoài đời coi chừng."]\}}
\soset{vd.otp.text}{Mã OTP xác nhận đăng nhập vừa gửi tới điện thoại của bạn. Nếu bạn đang đăng nhập trên thiết bị mới, vui lòng nhập mã trên màn hình ứng dụng để tiếp tục.}
\soset{vd.otp.dong}{\textit{Mã OTP xác nhận đăng nhập vừa gửi tới điện thoại của bạn. Nếu bạn đang đăng nhập trên thiết bị mới, vui lòng nhập mã trên màn hình ứng dụng để tiếp tục.}}
\soset{vd.otp.gold}{\nhan{AN_TOAN}}
\soset{vd.otp.pred}{\nhan{LUA_DAO}}
\soset{vd.otp.llm}{\nhan{AN_TOAN}}
\soset{vd.otp.tincay}{91,9\%}
\soset{vd.otp.kichban}{an\_toan\_chung}
\soset{vd.otp.vung}{Bắc}
\soset{vd.otp.bc}{``Mã OTP xác nhận đăng nhập vừa gửi'', ``nhập mã trên màn hình ứng''}
\soset{vd.otp.bcgold}{(rỗng)}
\soset{vd.otp.json}{\{"nhan": "LUA\_DAO", "bang\_chung": ["Mã OTP xác nhận đăng nhập vừa gửi", "nhập mã trên màn hình ứng"]\}}
\soset{vd.baoco.text}{Mọi người đừng lan truyền ảnh vụ cãi nhau nữa.\newline Vậy có cần báo cô không?\newline Có, nói cô đi. Để t gửi cho cô giải quyết, đừng tự ra gặp nhóm kia.}
\soset{vd.baoco.dong}{\textit{Mọi người đừng lan truyền ảnh vụ cãi nhau nữa.}\\ \textit{Vậy có cần báo cô không?}\\ \textit{Có, nói cô đi. Để t gửi cho cô giải quyết, đừng tự ra gặp nhóm kia.}}
\soset{vd.baoco.gold}{\nhan{AN_TOAN}}
\soset{vd.baoco.pred}{\nhan{DE_DOA}}
\soset{vd.baoco.llm}{\nhan{AN_TOAN}}
\soset{vd.baoco.tincay}{93,7\%}
\soset{vd.baoco.kichban}{DD6}
\soset{vd.baoco.vung}{Bắc}
\soset{vd.baoco.bc}{``Để t gửi cho cô giải quyết, đừng tự ra gặp nhóm kia.''}
\soset{vd.baoco.bcgold}{(rỗng)}
\soset{vd.baoco.json}{\{"nhan": "DE\_DOA", "bang\_chung": ["Để t gửi cho cô giải quyết, đừng tự ra gặp nhóm kia."]\}}
\soset{vd.ghiam.text}{Đứa nào gửi ghi âm cho cô thì chắc đang đọc tin này đấy nhỉ\newline Ai gửi thì kệ người ta đi\newline Ừ kệ thôi, giỏi thế cơ mà. Mai đi học nhớ đi đường đông người nhé}
\soset{vd.ghiam.dong}{\textit{Đứa nào gửi ghi âm cho cô thì chắc đang đọc tin này đấy nhỉ}\\ \textit{Ai gửi thì kệ người ta đi}\\ \textit{Ừ kệ thôi, giỏi thế cơ mà. Mai đi học nhớ đi đường đông người nhé}}
\soset{vd.ghiam.gold}{\nhan{DE_DOA}}
\soset{vd.ghiam.pred}{\nhan{AN_TOAN}}
\soset{vd.ghiam.llm}{\nhan{AN_TOAN}}
\soset{vd.ghiam.tincay}{87,4\%}
\soset{vd.ghiam.kichban}{DD6}
\soset{vd.ghiam.vung}{Bắc}
\soset{vd.ghiam.bc}{(rỗng)}
\soset{vd.ghiam.bcgold}{``Mai đi học nhớ đi đường đông người nhé''}
\soset{vd.ghiam.json}{\{"nhan": "AN\_TOAN", "bang\_chung": []\}}
\soset{vd.vay.text}{Thiếu tiền trả phòng trọ hả? Vay online 5 triệu chỉ cần CCCD, không hỏi phụ huynh. Nhắn tui, duyệt trong buổi.}
\soset{vd.vay.dong}{\textit{Thiếu tiền trả phòng trọ hả? Vay online 5 triệu chỉ cần CCCD, không hỏi phụ huynh. Nhắn tui, duyệt trong buổi.}}
\soset{vd.vay.gold}{\nhan{LUA_DAO}}
\soset{vd.vay.pred}{\nhan{AN_TOAN}}
\soset{vd.vay.llm}{\nhan{LUA_DAO}}
\soset{vd.vay.tincay}{92,8\%}
\soset{vd.vay.kichban}{LD7}
\soset{vd.vay.vung}{Nam}
\soset{vd.vay.bc}{(rỗng)}
\soset{vd.vay.bcgold}{``Vay online 5 triệu chỉ cần CCCD'', ``duyệt trong buổi''}
\soset{vd.vay.json}{\{"nhan": "AN\_TOAN", "bang\_chung": []\}}
\soset{vd.nham.text}{Chị lỡ chuyển nhầm 2 triệu vô tài khoản em rồi, em kiểm tra giùm chị nha. Em chuyển lại vô [STK] này giùm chị, chị cảm ơn em nhiều lắm}
\soset{vd.nham.dong}{\textit{Chị lỡ chuyển nhầm 2 triệu vô tài khoản em rồi, em kiểm tra giùm chị nha. Em chuyển lại vô [STK] này giùm chị, chị cảm ơn em nhiều lắm}}
\soset{vd.nham.gold}{\nhan{LUA_DAO}}
\soset{vd.nham.pred}{\nhan{AN_TOAN}}
\soset{vd.nham.llm}{\nhan{LUA_DAO}}
\soset{vd.nham.tincay}{83,4\%}
\soset{vd.nham.kichban}{moi}
\soset{vd.nham.vung}{Nam}
\soset{vd.nham.bc}{(rỗng)}
\soset{vd.nham.bcgold}{``Chị lỡ chuyển nhầm 2 triệu vô tài khoản em rồi'', ``Em chuyển lại vô [STK] này giùm chị''}
\soset{vd.nham.json}{\{"nhan": "AN\_TOAN", "bang\_chung": []\}}
\soset{vd.kiemsoat.text}{Tau thấy mi vừa theo dõi lại người yêu cũ. Bỏ theo dõi liền đi, rồi chụp màn hình gửi tau. Đừng bắt tau phải vào tài khoản mi tự kiểm tra.}
\soset{vd.kiemsoat.dong}{\textit{Tau thấy mi vừa theo dõi lại người yêu cũ. Bỏ theo dõi liền đi, rồi chụp màn hình gửi tau. Đừng bắt tau phải vào tài khoản mi tự kiểm tra.}}
\soset{vd.kiemsoat.gold}{\nhan{QUAY_ROI}}
\soset{vd.kiemsoat.pred}{\nhan{AN_TOAN}}
\soset{vd.kiemsoat.llm}{\nhan{QUAY_ROI}}
\soset{vd.kiemsoat.tincay}{92,5\%}
\soset{vd.kiemsoat.kichban}{QR6}
\soset{vd.kiemsoat.vung}{Trung}
\soset{vd.kiemsoat.bc}{(rỗng)}
\soset{vd.kiemsoat.bcgold}{``Bỏ theo dõi liền đi, rồi chụp màn hình gửi tau'', ``Đừng bắt tau phải vào tài khoản mi tự kiểm tra''}
\soset{vd.kiemsoat.json}{\{"nhan": "AN\_TOAN", "bang\_chung": []\}}
\soset{dl.goc.tong}{54.454}
\soset{dl.goc.vihos}{11.056}
\soset{dl.goc.vihsd}{33.398}
\soset{dl.goc.vithsd}{10.000}
\soset{dl.khacnhau}{30.580}
\soset{dl.khotest}{6.724}
\soset{dl.khotrain}{23.856}
\soset{dl.ganlai}{10.062}
\soset{dl.ganlai.nguyhai}{5.691}
\soset{dl.ganlai.blacklist}{1.276}
\soset{dl.ganlai.tukhoa}{95}
\soset{dl.ganlai.ngaunhien}{3.000}
\soset{dl.that.giu}{8.649}
\soset{dl.that.bo}{1.413}
\soset{dl.that.khopgoc}{7.674}
\soset{dl.that.bo.hate.antoan}{865}
\soset{dl.that.bo.clean.xucpham}{475}
\soset{dl.that.bo.clean.dedoa}{41}
\soset{dl.that.bo.clean.quayroi}{19}
\soset{dl.that.bo.clean.luadao}{13}
\soset{dl.sinh.tong}{10.557}
\soset{dl.sinh.giu}{9.625}
\soset{dl.sinh.loai}{932}
\soset{dl.sinh.tilegiu}{91,2\%}
\soset{dl.sinh.at.tilegiu}{93,4\%}
\soset{dl.sinh.dd.tilegiu}{92,3\%}
\soset{dl.sinh.ld.tilegiu}{96,1\%}
\soset{dl.sinh.qr.tilegiu}{73,8\%}
\soset{dl.sinh.xp.tilegiu}{98,3\%}
\soset{dl.sinh.loai.raxp}{424}
\soset{dl.sinh.qrba.tilegiu}{53,3\%}
\soset{dl.sinh.dongtb}{4,2}
\soset{tr.tong}{13.392}
\soset{tr.at.that}{3.967}
\soset{tr.at.sinh}{2.000}
\soset{tr.at}{5.967}
\soset{tr.xp.that}{2.527}
\soset{tr.xp.sinh}{470}
\soset{tr.xp}{2.997}
\soset{tr.qr.that}{62}
\soset{tr.qr.sinh}{1.098}
\soset{tr.qr}{1.160}
\soset{tr.dd.that}{156}
\soset{tr.dd.sinh}{1.361}
\soset{tr.dd}{1.517}
\soset{tr.ld.that}{0}
\soset{tr.ld.sinh}{1.751}
\soset{tr.ld}{1.751}
\soset{tr.that}{6.712}
\soset{tr.sinh}{6.680}
\soset{tr.motdong}{8.044}
\soset{te.tong}{334}
\soset{te.at}{96}
\soset{te.xp}{53}
\soset{te.qr}{51}
\soset{te.dd}{62}
\soset{te.ld}{72}
\soset{te.bac}{112}
\soset{te.trung}{111}
\soset{te.nam}{111}
\soset{te.chuathay}{149}
\soset{te.trongpp}{185}
\soset{te.chuathay.xp}{3}
\soset{te.chon.trung}{273}
\soset{te.chon.nguoi}{52}
\soset{te.chon.llm}{9}
\soset{te.nguoiviet}{4}
\soset{te.C1.n}{668}
\soset{te.C1.sokieu}{7}
\soset{te.C1.tiletu}{60,0\%}
\soset{te.C1.bodau}{102}
\soset{te.C1.chenkytu}{81}
\soset{te.C1.dinhtu}{96}
\soset{te.C1.gosai}{104}
\soset{te.C1.keodai}{89}
\soset{te.C1.sothaychu}{93}
\soset{te.C1.teencode}{103}
\soset{te.C2.n}{668}
\soset{te.C2.sokieu}{6}
\soset{te.C2.tiletu}{60,0\%}
\soset{te.C2.cyrillichomoglyph}{110}
\soset{te.C2.doubleelongation}{98}
\soset{te.C2.exoticseparator}{108}
\soset{te.C2.fullwidth}{122}
\soset{te.C2.spaceinjection}{114}
\soset{te.C2.zerowidth}{116}
\soset{te.C1.dathay}{398}
\soset{te.C1.chuathay}{270}
\soset{rr.trunggan}{0}
\soset{rr.nguong}{0,8}
\soset{rr.kichban}{DD6, LD12, LD7, QR6}
\soset{rr.trainkichban}{0}
\soset{dt.giu.dung}{86}
\soset{dt.giu.tong}{100}
\soset{dt.giu.tile}{86,0\%}
\soset{dt.giu.lo}{77,9\%}
\soset{dt.giu.hi}{91,5\%}
\soset{dt.bo.llm}{61,0\%}
\soset{dt.mau}{200}
\soset{cp.model}{deepseek/deepseek-v4.1-flash-1017}
\soset{cp.ganlai.goi}{503}
\soset{cp.ganlai.usd}{0,40}
\soset{cp.ganlai.phut}{7,8}
\soset{cp.sinh.goi}{2.140}
\soset{cp.sinh.usd}{1,35}
\soset{cp.sinh.phut}{31,3}
\soset{cp.loc.goi}{527}
\soset{cp.loc.usd}{0,53}
\soset{cp.loc.phut}{8,1}
\soset{cp.tong.goi}{3.170}
\soset{cp.tong.usd}{2,28}
\soset{cp.tong.vnd}{61.555}
\soset{cp.tong.phut}{47}
\soset{cp.tong.trieutoken}{17,0}
\soset{api.C0.p50}{3,46}
\soset{api.C0.p95}{9,55}
\soset{api.C0.usd1000}{0,085}
\soset{api.songsong}{4}
\soset{api.provider}{venice/int4}
\soset{api.reasoning}{low}
\soset{api.C1.p50}{3,16}
\soset{api.C1.p95}{8,22}
\soset{api.C1.usd1000}{0,085}
\soset{api.C2.p50}{2,83}
\soset{api.C2.p95}{4,85}
\soset{api.C2.usd1000}{0,085}
\soset{api.tong.usd}{0,14}
\soset{v2.p50}{39,4}
\soset{v2.p95}{42,7}
\soset{v2.trongso}{1.177}
\soset{v2.dinh}{1.218}
\soset{v2.T}{1,22}
\soset{v2.nguongbc}{0,30}
\soset{v2.nguongmottu}{0,40}
\soset{v2.maxtok}{512}
\soset{v2.epoch}{5}
\soset{v2.accent}{30,0\%}
\soset{v2.surface}{20,0\%}
\soset{v2.fgm}{0,5}
\soset{tb.a.cpu}{AMD Ryzen AI 7 H 350 w/ Radeon 860M}
\soset{tb.a.ort}{1.26.0}
\soset{tb.a.int8.8.n}{100}
\soset{tb.a.int8.8.p50}{24,7}
\soset{tb.a.int8.8.p95}{31,5}
\soset{tb.a.int8.8.ram}{986}
\soset{tb.a.int8.8.nap}{1,5}
\soset{tb.a.int8.4.n}{100}
\soset{tb.a.int8.4.p50}{28,5}
\soset{tb.a.int8.4.p95}{41,3}
\soset{tb.a.int8.4.ram}{982}
\soset{tb.a.int8.4.nap}{1,6}
\soset{tb.a.int8.1.n}{50}
\soset{tb.a.int8.1.p50}{61,5}
\soset{tb.a.int8.1.p95}{92,5}
\soset{tb.a.int8.1.ram}{982}
\soset{tb.a.int8.1.nap}{1,4}
\soset{tb.a.fp32.8.n}{60}
\soset{tb.a.fp32.8.p50}{44,0}
\soset{tb.a.fp32.8.p95}{58,5}
\soset{tb.a.fp32.8.ram}{2.419}
\soset{tb.a.fp32.8.nap}{2,0}
\soset{tb.a.torch.8.n}{40}
\soset{tb.a.torch.8.p50}{60,4}
\soset{tb.a.torch.8.p95}{80,4}
\soset{tb.a.torch.8.ram}{4.781}
\soset{tb.a.torch.8.nap}{20,4}
\soset{tb.int8.mib}{301}
\soset{tb.fp32.mib}{1.178}
\soset{tb.tok.mib}{33}
\soset{tb.nen}{3,9}
\soset{tb.int8.C0.khop}{97,9\%}
\soset{tb.int8.C1.khop}{97,8\%}
\soset{tb.int8.C2.khop}{95,8\%}
\soset{lt.C0.fp32.macro}{86,5}
\soset{lt.C0.fp32.bcy}{83,0}
\soset{lt.C0.fp32.fpr}{10,4\%}
\soset{lt.C0.int8.macro}{86,0}
\soset{lt.C0.int8.bcy}{84,0}
\soset{lt.C0.int8.fpr}{10,4\%}
\soset{lt.C0.d.macro}{$-$0,5}
\soset{lt.C0.d.bcy}{+1,0}
\soset{lt.C0.p}{1,00}
\soset{lt.C0.n}{334}
\soset{lt.C1.fp32.macro}{78,6}
\soset{lt.C1.fp32.bcy}{76,5}
\soset{lt.C1.fp32.fpr}{10,9\%}
\soset{lt.C1.int8.macro}{77,7}
\soset{lt.C1.int8.bcy}{76,7}
\soset{lt.C1.int8.fpr}{10,9\%}
\soset{lt.C1.d.macro}{$-$0,9}
\soset{lt.C1.d.bcy}{+0,2}
\soset{lt.C1.p}{0,58}
\soset{lt.C1.n}{668}
\soset{lt.C2.fp32.macro}{77,1}
\soset{lt.C2.fp32.bcy}{75,3}
\soset{lt.C2.fp32.fpr}{10,4\%}
\soset{lt.C2.int8.macro}{75,7}
\soset{lt.C2.int8.bcy}{76,3}
\soset{lt.C2.int8.fpr}{13,0\%}
\soset{lt.C2.d.macro}{$-$1,4}
\soset{lt.C2.d.bcy}{+1,0}
\soset{lt.C2.p}{0,12}
\soset{lt.C2.n}{668}
\soset{lt.ALL.fp32.macro}{79,6}
\soset{lt.ALL.fp32.bcy}{77,4}
\soset{lt.ALL.fp32.fpr}{10,6\%}
\soset{lt.ALL.int8.macro}{78,6}
\soset{lt.ALL.int8.bcy}{78,1}
\soset{lt.ALL.int8.fpr}{11,7\%}
\soset{lt.ALL.d.macro}{$-$1,0}
\soset{lt.ALL.d.bcy}{+0,7}
\soset{lt.ALL.p}{0,07}
\soset{lt.ALL.n}{1.670}
\soset{tb.int8.n}{1.670}
\soset{tb.int8.macro}{78,6}
\soset{tb.fp32.macro}{79,6}
\soset{tb.int8.bcy}{78,1}
\soset{tb.fp32.bcy}{77,4}
\soset{tb.int8.khop}{97,0\%}
\soset{tb.int8.mcnemar}{0,07}
\soset{tb.b.cpu}{Intel(R) Core(TM) Ultra 7 155H}
\soset{tb.b.ort}{1.32.0}
\soset{tb.b.n}{1.670}
\soset{tb.b.4.p50}{81,6}
\soset{tb.b.4.p95}{150,6}
\soset{tb.b.4.ram}{903}
\soset{tb.b.4.nap}{5,2}
\soset{tb.b.4.tl}{11,3}
\soset{tb.b.8.p50}{69,1}
\soset{tb.b.8.p95}{120,8}
\soset{tb.b.8.ram}{904}
\soset{tb.b.8.nap}{4,5}
\soset{tb.b.8.tl}{13,5}
\soset{tb.b.1.p50}{197,2}
\soset{tb.b.1.p95}{450,2}
\soset{tb.b.1.ram}{902}
\soset{tb.b.1.nap}{5,2}
\soset{tb.b.1.tl}{4,4}
\soset{tb.b.pnghi}{8,6}
\soset{tb.b.ptai}{14,3}
\soset{tb.b.mj}{510}
\soset{tb.b.mjtong}{1.285}
\soset{tb.b.wh}{0,14}
\soset{tb.b.giay}{300}
\soset{tb.b.khop.torch}{97,0\%}
\soset{tb.b.khop.int8}{97,0\%}
\soset{tb.b.macro}{78,3}
\soset{tb.b.bcy}{77,7}
\soset{tb.web.trinhduyet}{Edge 158.0.4370.33}
\soset{tb.web.ortweb}{1.31.1}
\soset{tb.web.backend}{WASM · 4 luồng}
\soset{tb.web.n}{1.670}
\soset{tb.web.p50}{337}
\soset{tb.web.p95}{605}
\soset{tb.web.tl}{2,8}
\soset{tb.web.nap}{7,2}
\soset{tb.web.taive}{348}
\soset{tb.web.ram.trong}{434}
\soset{tb.web.ram.dinh}{1.383}
\soset{tb.web.ram.nap}{1.823}
\soset{tb.web.ram.tang}{949}
\soset{tb.web.pnghi}{7,7}
\soset{tb.web.ptai}{10,3}
\soset{tb.web.mj}{966}
\soset{tb.web.mjtong}{3.744}
\soset{tb.web.wh}{0,27}
\soset{tb.web.khop.torch}{96,8\%}
\soset{tb.web.khop.cli}{97,2\%}
\soset{tb.web.khop.cli.bc}{72,1\%}
\soset{tb.web.macro}{78,4}
\soset{tb.web.bcy}{77,8}
\soset{cnt.C0.n}{334}
\soset{cnt.C0.raw}{86,5}
\soset{cnt.C0.norm}{85,2}
\soset{cnt.C0.gain}{$-$1,2}
\soset{cnt.C0.bc.raw}{83,0}
\soset{cnt.C0.bc.norm}{84,5}
\soset{cnt.C0.bc.gain}{+1,5}
\soset{cnt.C1.n}{668}
\soset{cnt.C1.raw}{78,6}
\soset{cnt.C1.norm}{82,8}
\soset{cnt.C1.gain}{+4,2}
\soset{cnt.C1.bc.raw}{76,5}
\soset{cnt.C1.bc.norm}{81,7}
\soset{cnt.C1.bc.gain}{+5,2}
\soset{cnt.C2.n}{668}
\soset{cnt.C2.raw}{77,1}
\soset{cnt.C2.norm}{75,7}
\soset{cnt.C2.gain}{$-$1,4}
\soset{cnt.C2.bc.raw}{75,3}
\soset{cnt.C2.bc.norm}{74,1}
\soset{cnt.C2.bc.gain}{$-$1,2}
\soset{cnt.chuathay.n}{764}
\soset{cnt.chuathay.raw}{78,4}
\soset{cnt.chuathay.norm}{77,3}
\soset{cnt.chuathay.gain}{$-$1,1}
\soset{cnt.chuathay.bc.raw}{76,6}
\soset{cnt.chuathay.bc.norm}{75,6}
\soset{cnt.chuathay.bc.gain}{$-$1,0}
\soset{cnt.kieu.bodau.gain}{+0,6}
\soset{cnt.kieu.chenkytu.gain}{+17,5}
\soset{cnt.kieu.dinhtu.gain}{+1,0}
\soset{cnt.kieu.gosai.gain}{$-$1,2}
\soset{cnt.kieu.keodai.gain}{+3,8}
\soset{cnt.kieu.sothaychu.gain}{+3,5}
\soset{cnt.kieu.teencode.gain}{+4,2}
\soset{cnt.kieu.cyrillichomoglyph.gain}{$-$1,4}
\soset{cnt.kieu.doubleelongation.gain}{+2,1}
\soset{cnt.kieu.exoticseparator.gain}{+1,8}
\soset{cnt.kieu.fullwidth.gain}{$-$0,7}
\soset{cnt.kieu.spaceinjection.gain}{$-$3,6}
\soset{cnt.kieu.zerowidth.gain}{$-$4,3}
\soset{v2.gpu}{NVIDIA L4}
\soset{v2.torch}{2.14.1+cu130}
\soset{v2.transformers}{4.57.6}
\soset{v2.datarev}{546300b}
\soset{dev.chia.fit}{10.713}
\soset{dev.chia.dev}{1.340}
\soset{dev.chia.cal}{1.339}
\soset{dev.mmbert.macro}{92,3}
\soset{dev.mmbert.bin}{92,6}
\soset{dev.mmbert.fpr}{8,4\%}
\soset{dev.mmbert.bcf}{69,7}
\soset{dev.mmbert.toan}{7,2\%}
\soset{dev.mmbert.gap}{1,5}
\soset{dev.mmbert.tieng}{74,8}
\soset{dev.mmbert.tieng.fpr}{14,2\%}
\soset{dev.mmbert.trongso}{1.177}
\soset{dev.mmbert.p50}{39,4}
\soset{dev.mmbert.trainphut}{44}
\soset{dev.mmbert.trainvram}{9.283}
\soset{dev.bamibert.macro}{92,2}
\soset{dev.bamibert.bin}{93,3}
\soset{dev.bamibert.fpr}{8,6\%}
\soset{dev.bamibert.bcf}{71,6}
\soset{dev.bamibert.toan}{6,9\%}
\soset{dev.bamibert.gap}{1,9}
\soset{dev.bamibert.tieng}{55,5}
\soset{dev.bamibert.tieng.fpr}{14,2\%}
\soset{dev.bamibert.trongso}{397}
\soset{dev.bamibert.p50}{18,6}
\soset{dev.bamibert.trainphut}{19}
\soset{dev.bamibert.trainvram}{5.904}
\soset{dev.visobert.macro}{92,9}
\soset{dev.visobert.bin}{93,5}
\soset{dev.visobert.fpr}{6,0\%}
\soset{dev.visobert.bcf}{70,6}
\soset{dev.visobert.toan}{5,7\%}
\soset{dev.visobert.gap}{1,8}
\soset{dev.visobert.tieng}{62,7}
\soset{dev.visobert.tieng.fpr}{11,5\%}
\soset{dev.visobert.trongso}{376}
\soset{dev.visobert.p50}{20,8}
\soset{dev.visobert.trainphut}{19}
\soset{dev.visobert.trainvram}{5.634}
\soset{dev.guard.macro}{90,0}
\soset{dev.guard.bin}{91,2}
\soset{dev.guard.fpr}{10,7\%}
\soset{dev.guard.bcf}{67,7}
\soset{dev.guard.toan}{9,9\%}
\soset{dev.guard.gap}{6,5}
\soset{dev.guard.tieng}{73,4}
\soset{dev.guard.tieng.fpr}{16,8\%}
\soset{dev.guard.trongso}{1.177}
\soset{dev.guard.p50}{82,3}
\soset{dev.guard.trainphut}{21}
\soset{dev.guard.trainvram}{7.303}
\soset{dev.tieng.n}{178}
\soset{dev.mc.guard}{0,023}
\soset{dev.mc.visobert}{0,29}
\soset{dev.mc.bamibert}{0,69}
\soset{ctx.n}{587}
\soset{ctx.full}{98,0}
\soset{ctx.last}{70,9}
\soset{ctx.gain}{+27,1}
\soset{ctx.lenonly}{30,2}
\soset{cn.cleanc0.n}{1.340}
\soset{cn.cleanc0.raw}{92,3}
\soset{cn.cleanc0.norm}{91,6}
\soset{cn.cleanc0.gain}{$-$0,8}
\soset{cn.accentstrip.n}{1.216}
\soset{cn.accentstrip.raw}{87,5}
\soset{cn.accentstrip.norm}{87,2}
\soset{cn.accentstrip.gain}{$-$0,3}
\soset{cn.accentstrip.goc}{92,0}
\soset{cn.accentstrip.tut}{$-$4,6}
\soset{cn.c1indist.n}{1.255}
\soset{cn.c1indist.raw}{88,8}
\soset{cn.c1indist.norm}{90,4}
\soset{cn.c1indist.gain}{+1,6}
\soset{cn.c1indist.goc}{92,3}
\soset{cn.c1indist.tut}{$-$3,5}
\soset{cn.c1heldoutcontrolled.n}{1.153}
\soset{cn.c1heldoutcontrolled.raw}{92,6}
\soset{cn.c1heldoutcontrolled.norm}{91,3}
\soset{cn.c1heldoutcontrolled.gain}{$-$1,3}
\soset{cn.c1heldoutcontrolled.goc}{92,7}
\soset{cn.c1heldoutcontrolled.tut}{$-$0,1}
\soset{cn.blindspothomoglyph.n}{1.203}
\soset{cn.blindspothomoglyph.raw}{91,0}
\soset{cn.blindspothomoglyph.norm}{91,0}
\soset{cn.blindspothomoglyph.gain}{0,0}
\soset{cn.blindspothomoglyph.goc}{93,1}
\soset{cn.blindspothomoglyph.tut}{$-$2,0}
\soset{cn.blindspotzerowidth.n}{1.330}
\soset{cn.blindspotzerowidth.raw}{90,1}
\soset{cn.blindspotzerowidth.norm}{89,3}
\soset{cn.blindspotzerowidth.gain}{$-$0,7}
\soset{cn.blindspotzerowidth.goc}{92,4}
\soset{cn.blindspotzerowidth.tut}{$-$2,3}
\soset{v1goc.C0}{85,3}
\soset{v1goc.C1}{82,7}
\soset{v1goc.C2}{66,5}
\soset{v1goc.n}{2.376}
\soset{v1.onnx}{1.039}
\soset{plot.bin.v1}{(C0,31.2) (C1,31.0) (C2,39.5)}
\soset{plot.bin.llm}{(C0,89.2) (C1,74.5) (C2,77.9)}
\soset{plot.bin.v2}{(C0,92.9) (C1,90.3) (C2,90.9)}


% Tên nhãn, đường dẫn tệp
\newcommand{\nhan}[1]{\texttt{\detokenize{#1}}}
\newcommand{\tep}[1]{\nolinkurl{#1}}
\expandafter\def\expandafter\UrlBreaks\expandafter{\UrlBreaks\do\_\do\-}
% Hộp ghi nhớ ngắn
\newcommand{\ghinho}[1]{\par\noindent\colorbox{nenxanh}{\parbox{\dimexpr\linewidth-2\fboxsep}{\small #1}}\par}

\pagestyle{fancy}
\fancyhf{}
\fancyhead[L]{\small SafeViChi v2}
\fancyhead[R]{\small Báo cáo vòng chung kết}
\fancyfoot[C]{\thepage}
\setlength{\headheight}{15pt}

\begin{document}

% ---------- Trang bìa ----------
\begin{titlepage}
\centering
\vspace*{30mm}
{\Huge\bfseries\color{navy} SafeViChi v2\par}
\vspace{6mm}
{\Large Phát hiện tin nhắn tiếng Việt nguy hại theo tình huống,\\ kèm bằng chứng, chạy trên thiết bị\par}
\vspace{12mm}
{\large Báo cáo vòng chung kết\par}
\vspace{40mm}
\begin{tabular}{ll}
Phiên bản mã & nhánh \tep{v2-final} \\
Dữ liệu train/test & commit \texttt{\so{v2.datarev}} \\
Ngày nộp & 26/11/2026 \\
\end{tabular}
\vfill
{\small Mọi con số trong báo cáo được sinh tự động từ thư mục \tep{results/v2/} bằng\\
\tep{python -m src.v2.eval.xuat_so_lieu_bao_cao}.\par}
\end{titlepage}

\tableofcontents
\newpage

% Phần 0: tóm tắt. Mọi số lấy qua \so{...}.
{% Giữ phần Tóm tắt gọn trong 1 trang, không ảnh hưởng các phần sau
\setlength{\parskip}{0.25em}
\setlength{\intextsep}{4pt plus 1pt minus 1pt}
\renewcommand{\arraystretch}{1.04}
\section*{Tóm tắt}
\addcontentsline{toc}{section}{Tóm tắt}

SafeViChi v2 đọc một đoạn hội thoại tiếng Việt (1--8 dòng) và trả về hai đầu ra: \textbf{tình huống} của đoạn đó
(một trong 5 nhãn \nhan{AN_TOAN}, \nhan{XUC_PHAM}, \nhan{QUAY_ROI}, \nhan{DE_DOA}, \nhan{LUA_DAO}) và
\textbf{bằng chứng} là các cụm từ trích nguyên văn từ tin nhắn. Mô hình chạy cục bộ trên thiết bị người dùng, không gửi
tin nhắn ra ngoài.

\textbf{Thay đổi chính so với SafeViChi v1:}
SafeViChi v1 chỉ phân loại nhị phân \nhan{HATE}/\nhan{CLEAN} nên chủ yếu bắt được nội dung xúc phạm: trên tập kiểm
thử mới (C0), v1 gắn \nhan{HATE} cho \so{C0.v1.xp.r} số đoạn \nhan{XUC_PHAM} nhưng chỉ \so{C0.v1.dd.r} số đoạn
\nhan{DE_DOA}, \so{C0.v1.qr.r} số đoạn \nhan{QUAY_ROI} và \so{C0.v1.ld.r} số đoạn \nhan{LUA_DAO}.
SafeViChi v2 dùng 5 nhãn theo tình huống, đọc cả đoạn hội thoại thay vì một câu và trích bằng chứng cùng lúc với phân
loại. Dữ liệu huấn luyện gồm hai nguồn đã lọc: bình luận thực tế (ViHSD, ViHOS, ViTHSD) được LLM gán lại nhãn, chỉ giữ mẫu đồng thuận với người gán gốc; và hội thoại tổng hợp theo kịch bản, được gán lại nhãn độc lập (ẩn nhãn yêu cầu) để lọc nhiễu. Tập kiểm thử do thành viên nhóm tự biên soạn toàn bộ.

\textbf{Kết quả trên tập kiểm thử} (\so{te.tong} đoạn gốc C0; \so{te.C1.n} bản biến dạng bề mặt thường
gặp C1, gồm cả kiểu đã thấy và chưa thấy; \so{te.C2.n} bản biến dạng đổi lớp ký tự Unicode ``điểm mù'' C2, không có
trong tập huấn luyện):

\begin{table}[H]
\centering
\small
\setlength{\tabcolsep}{4.5pt}
\caption{Kết quả chính trên tập kiểm thử (thang 0--100, tức F1 $\times$ 100). Macro-F1: trung bình F1 trên 5 nhãn;
Binary-F1: F1 phân biệt nguy hại/an toàn; BC F1: F1 bằng chứng theo ý (mức cụm). SafeViChi v1 chỉ có nhãn nhị phân,
không trích bằng chứng.}
\label{tab:p0-chinh}
\begin{tabular}{l c ccc ccc}
\toprule
& \textbf{SafeViChi v1} & \multicolumn{3}{c}{\textbf{gpt-oss-120b (API)}} & \multicolumn{3}{c}{\textbf{SafeViChi v2}} \\
\cmidrule(lr){2-2}\cmidrule(lr){3-5}\cmidrule(lr){6-8}
Tập & Binary-F1 & Macro-F1 & Binary-F1 & BC F1 & Macro-F1 & Binary-F1 & BC F1 \\
\midrule
C0 & \so{C0.v1.bin} & \so{C0.llm.macro} & \so{C0.llm.bin} & \so{C0.llm.bcf} & \textbf{\so{C0.v2.macro}} & \textbf{\so{C0.v2.bin}} & \textbf{\so{C0.v2.bcf}} \\
C1 & \so{C1.v1.bin} & \so{C1.llm.macro} & \so{C1.llm.bin} & \so{C1.llm.bcf} & \textbf{\so{C1.v2.macro}} & \textbf{\so{C1.v2.bin}} & \textbf{\so{C1.v2.bcf}} \\
C2 & \so{C2.v1.bin} & \so{C2.llm.macro} & \so{C2.llm.bin} & \so{C2.llm.bcf} & \textbf{\so{C2.v2.macro}} & \textbf{\so{C2.v2.bin}} & \textbf{\so{C2.v2.bcf}} \\
\bottomrule
\end{tabular}
\end{table}

\begin{itemize}[itemsep=0.05em,topsep=0.1em]
  \item SafeViChi v2 cao nhất trên cả ba tập và ổn định hơn khi tin nhắn bị biến dạng; gpt-oss-120b (API) giảm mạnh
        trên C1, C2, riêng BC F1 giảm từ \so{C0.llm.bcf} (C0) xuống \so{C2.llm.bcf} (C2).
  \item Trọng số FP32 nặng \so{v2.trongso}~MiB (dưới ngưỡng 2~GB khi chưa lượng tử hóa), độ trễ trung vị
        (p50, batch~1) \so{v2.p50}~ms mỗi đoạn trên GPU \so{v2.gpu}. Bản ONNX INT8 cho thiết bị đích còn
        \so{tb.int8.mib}~MiB (nhỏ hơn \so{tb.nen} lần), p50 \so{tb.a.int8.8.p50}~ms trên
        CPU laptop (8 luồng), \so{tb.mob.4.p50}~ms trên điện thoại Android (4 luồng, RAM đỉnh \so{tb.mob.4.ram}~MiB) và
        \so{tb.web.p50}~ms trong trình duyệt (4 luồng). Chi phí API LLM để xây dữ liệu huấn luyện tổng cộng
        \so{cp.tong.usd}~USD.
  \item Hạn chế chính: tỉ lệ báo nhầm (FPR) trên đoạn an toàn còn cao (\so{C0.v2.fpr} trên C0); hiệu năng giảm mạnh
        trên kịch bản chưa thấy khi huấn luyện (Macro-F1 trên C0 giảm từ \so{phan.C0.trongphanphoi.v2.macro} xuống
        \so{phan.C0.chuathay.v2.macro}); suy luận trong trình duyệt còn chậm hơn đáng kể so với CLI.
\end{itemize}

\begin{table}[H]
\centering
\small
\caption{Ba yêu cầu cải tiến của ban tổ chức và nội dung nhóm thực hiện.}
\label{tab:p0-yeucau}
\begin{tabularx}{\linewidth}{>{\raggedright\arraybackslash}p{0.28\linewidth}Y}
\toprule
\textbf{Yêu cầu của ban tổ chức} & \textbf{Nội dung thực hiện} \\
\midrule
1. Giảm giải thích sai và kiểm duyệt sai & Tách nhãn và bằng chứng thành hai đầu ra, hai ngưỡng và hai bộ độ đo riêng; phân tích các ca báo nhầm và bỏ sót (Mục~\ref{sec:p2-yc1}). \\
2. Khái quát trước kiểu né lọc mới & Đánh giá riêng theo kiểu biến dạng, kịch bản chưa thấy và giọng vùng miền; kiểm tra rò rỉ dữ liệu huấn luyện--kiểm thử; phân tích các ca thất bại (Mục~\ref{sec:yc2}). \\
3. So sánh phương pháp và chi phí & So sánh SafeViChi v1, LLM qua API, 4 ứng viên mô hình nhỏ và các thử nghiệm ablation; đo dung lượng, bộ nhớ, độ trễ, năng lượng và chi phí API (Mục~\ref{sec:yc3}). \\
\bottomrule
\end{tabularx}
\end{table}
}
\newpage

% Phần 1: thay đổi từ v1 sang v2, năm nhãn, dữ liệu. Số liệu lấy qua \so{...} (so_lieu.tex).

\section{Từ v1 đến v2: thay đổi gì và vì sao}
\label{sec:p1-v1v2}

\textbf{v1 (vòng loại)} phân loại một câu bình luận thành HATE hoặc CLEAN. Câu đi qua bộ chuẩn hóa 5 tầng
(Unicode, ký tự đồng dạng, ký tự phân cách, ký tự lặp, teencode), rồi vào ViHateT5 đã fine-tune đối kháng (FGM).
Phần giải thích dùng occlusion: che lần lượt từng cụm hai từ, chạy lại mô hình, cụm nào làm điểm rủi ro giảm mạnh
thì tô sáng. Trên test của vòng loại (\so{v1goc.n} câu bình luận ViHSD/ViHOS), F1 lớp HATE của v1 là
\so{v1goc.C0} với câu sạch, \so{v1goc.C1} với câu né lọc (C1) và \so{v1goc.C2} với kiểu điểm mù (C2).

\textbf{Vấn đề.} Khi chạy lại v1 trên test mới của v2 (hội thoại do người viết, mục~\ref{sec:p1-test}), v1 báo
HATE cho \so{C0.v1.xp.r} số đoạn \nhan{XUC_PHAM}, nhưng chỉ \so{C0.v1.dd.r} đoạn đe dọa, \so{C0.v1.qr.r} đoạn
quấy rối và \so{C0.v1.ld.r} đoạn lừa đảo. Lý do nằm ở định nghĩa: HATE trong các bộ dữ liệu cũ gần với ``chửi bậy,
miệt thị''. Còn tin nhắn nguy hại thật thì khác: tin lừa đảo thường rất lịch sự, lời dọa có thể không có từ tục
nào. Vì vậy v2 thay đổi ba thứ:

\begin{itemize}
  \item \textbf{Nhãn theo tình huống.} 5 nhãn thay cho HATE/CLEAN, vì mỗi tình huống cần ứng dụng phản ứng khác
        nhau: làm mờ câu chửi, gợi ý chặn kẻ quấy rối, cảnh báo mạnh khi bị dọa, vô hiệu link lừa đảo.
  \item \textbf{Đọc cả đoạn hội thoại.} Ý đồ xấu thường chỉ lộ ở vài dòng cuối, sau mấy câu chào hỏi bình thường.
        Mô hình nhận tối đa 8 dòng gần nhất thay vì một câu.
  \item \textbf{Bằng chứng ra cùng lúc với nhãn.} Mô hình tự trả về các cụm từ làm bằng chứng trong cùng một lượt
        chạy, không phải chạy lại nhiều lần như occlusion. Nhãn và bằng chứng có ngưỡng riêng, được đo riêng.
\end{itemize}

\begin{table}[H]
\centering
\caption{So sánh v1 và v2}
\label{tab:p1-v1v2}
\small
\begin{tabularx}{\linewidth}{@{}p{2.6cm}YY@{}}
\toprule
 & \textbf{v1 (vòng loại)} & \textbf{v2 (chung kết)} \\
\midrule
Nhãn & HATE / CLEAN & 5 nhãn theo tình huống (mục~\ref{sec:p1-nhan}) \\
Đầu vào & Một câu bình luận & Đoạn hội thoại hai người, 1--8 dòng, mỗi dòng là một lượt nói \\
Mô hình & ViHateT5 (encoder--decoder), fine-tune đối kháng FGM & mmBERT-base, một bộ mã hóa với hai đầu ra: nhãn và
  cụm bằng chứng \\
Chống né lọc & Chuẩn hóa 5 tầng + FGM & Tăng cường dữ liệu khi train (bỏ dấu, biến dạng bề mặt) + FGM + nhánh đọc
  theo byte; chuẩn hóa thành tùy chọn và được đo bật/tắt \\
Giải thích & Occlusion sau khi phân loại, chạy lại mô hình nhiều lần & Gắn thẻ cụm bằng chứng ngay trong lượt
  chạy; cụm luôn chép nguyên văn từ tin nhắn \\
Dữ liệu train & ViHSD, ViHOS với nhãn gốc & Câu thật được LLM gán lại + hội thoại LLM sinh theo kịch bản
  (mục~\ref{sec:p1-dulieu}) \\
Dữ liệu test & Split test của ViHSD/ViHOS & Hội thoại do người trong nhóm tự viết, kèm bản biến dạng C1, C2 \\
Cách đo & F1 lớp HATE trên C0/C1/C2 & F1 từng nhãn; theo giọng vùng miền và kịch bản chưa thấy; precision/recall
  bằng chứng đo riêng; so với LLM lớn qua API \\
Dung lượng (FP32) & ONNX \so{v1.onnx} MiB & Trọng số \so{v2.trongso} MiB \\
\bottomrule
\end{tabularx}
\end{table}

\section{Năm nhãn mới}
\label{sec:p1-nhan}

Mỗi mẫu được gán đúng một nhãn trong bảng~\ref{tab:p1-nhan}. Định nghĩa đầy đủ và các ca khó nằm trong
\tep{docs/v2/HUONG_DAN_GAN_NHAN.md}; người gán, LLM gán sơ bộ và LLM baseline đều dùng chung tài liệu này.

\begin{table}[H]
\centering
\caption{Năm nhãn của v2}
\label{tab:p1-nhan}
\small
\begin{tabularx}{\linewidth}{@{}p{2.1cm}YYp{2.9cm}@{}}
\toprule
\textbf{Nhãn} & \textbf{Nghĩa} & \textbf{Ví dụ} & \textbf{Ứng dụng làm gì} \\
\midrule
\nhan{AN_TOAN} & Không thuộc 4 nhãn dưới & ``Đơn hàng của bạn đang được giao, vui lòng giữ máy'' & Hiện bình thường \\
\nhan{XUC_PHAM} & Lăng mạ, miệt thị một người hoặc một nhóm (kể cả vùng miền) & ``Đúng là não cá vàng, nói mãi
  không hiểu'' & Làm mờ, bấm mới xem \\
\nhan{QUAY_ROI} & Gạ gẫm không mong muốn, bám riết, bêu riếu đời tư & ``Sao seen mà không rep, anh nhắn 50 tin rồi
  đó'' & Gợi ý chặn, báo cáo \\
\nhan{DE_DOA} & Dọa gây hại thân thể, tài sản, danh dự; tống tiền & ``Mai tao chặn đầu mày ở cổng trường'' & Cảnh
  báo mạnh, gợi ý lưu bằng chứng \\
\nhan{LUA_DAO} & Thông tin giả để dụ chuyển tiền, đưa OTP, bấm link, cài ứng dụng & ``Chúc mừng bạn trúng iPhone,
  nộp phí ship 300k để nhận quà'' & Cảnh báo, vô hiệu link \\
\bottomrule
\end{tabularx}
\end{table}

\textbf{Khi một đoạn dính nhiều nhãn}, người gán hỏi lần lượt và dừng ở câu trả lời ``có'' đầu tiên:
\begin{enumerate}
  \item Có dọa gây hại, hoặc dọa để ép ai làm gì không? $\to$ \nhan{DE_DOA}
  \item Có dụ đưa tiền, OTP, mật khẩu, bấm link, cài ứng dụng bằng thông tin giả không? $\to$ \nhan{LUA_DAO}
  \item Có gạ gẫm, bình phẩm cơ thể, bám riết hoặc bêu riếu một người cụ thể không? $\to$ \nhan{QUAY_ROI}
  \item Có chửi, lăng mạ, miệt thị người hoặc nhóm người không? $\to$ \nhan{XUC_PHAM}
  \item Còn lại $\to$ \nhan{AN_TOAN}
\end{enumerate}
Ví dụ: ``Chuyển 10 triệu trước 8h không tao gửi clip cho cả nhà mày'' có đòi tiền, nhưng là tống tiền nên thuộc
\nhan{DE_DOA}, không phải \nhan{LUA_DAO}. Với đoạn nhiều dòng, nhãn của cả đoạn là tình huống nặng nhất xuất hiện
trong đoạn.

\textbf{Ca khó: trông nguy hại nhưng an toàn.} Đây là nguồn báo nhầm chính, nên được ghi rõ trong hướng dẫn và
được sinh thêm làm dữ liệu train (mục~\ref{sec:p1-sinh}):

\begin{table}[H]
\centering
\caption{Một số ca \nhan{AN_TOAN} dễ nhầm}
\label{tab:p1-cakho}
\small
\begin{tabularx}{\linewidth}{@{}p{6.4cm}Y@{}}
\toprule
\textbf{Tin nhắn} & \textbf{Vì sao an toàn} \\
\midrule
``Vãi cả chưởng, kèo này căng đét'' & Chửi thề cảm thán, không nhắm vào ai \\
``Tao giết mày luôn á, cười đau bụng'' & Nói quá kiểu đùa, dấu hiệu đùa nằm ngay trong câu \\
``Mã OTP của bạn là 123456. Không chia sẻ mã này cho bất kỳ ai'' & Tin OTP thật: không đòi gì, còn dặn đừng chia sẻ \\
``Ai bảo người miền Tây lười là phân biệt vùng miền đấy'' & Trích lại để phản bác \\
\bottomrule
\end{tabularx}
\end{table}

\textbf{Bằng chứng} là các cụm chép nguyên văn từ tin nhắn, cho biết vì sao đoạn đó nguy hại. Mỗi cụm nằm gọn
trong một dòng, lấy cụm ngắn nhất còn đủ nghĩa, được lấy nhiều cụm. \nhan{AN_TOAN} luôn rỗng; câu nguy hại cũng
được để rỗng nếu cái hại nằm ở cả câu (ví dụ mỉa mai). Ví dụ dưới đây là một đoạn thật trong test (giọng
\so{vd.dung.vung}, kịch bản \so{vd.dung.kichban}: dọa tìm ra ngoài đời trong game) và đầu ra của v2:

\ghinho{\textbf{Đầu vào} (mỗi dòng là một lượt nói, không ghi ai nói; nhãn đúng: \so{vd.dung.gold}):\\[2pt]
\so{vd.dung.dong}\\[4pt]
\textbf{Đầu ra của v2:}\\[2pt]
{\ttfamily \so{vd.dung.json}}}

\section{Dữ liệu}
\label{sec:p1-dulieu}

Dữ liệu v2 có ba phần, đi theo hai đường riêng (hình~\ref{fig:p1-luong}). Train gồm câu thật được LLM gán lại
theo 5 nhãn và hội thoại do LLM sinh theo kịch bản. Test do người trong nhóm tự viết, không có câu nào do LLM sinh.

\begin{figure}[H]
\centering
\begin{tikzpicture}[
  hop/.style={draw=navy, fill=nenxanh, rounded corners=2pt, align=center, font=\footnotesize,
              text width=2.55cm, minimum height=1.3cm, inner sep=3pt},
  hopdich/.style={hop, fill=navy!15, text width=1.9cm, font=\footnotesize\bfseries},
  hoptest/.style={hop, fill=nenxam, draw=chuxam},
  mui/.style={-{Stealth[length=2mm]}, thick, navy},
  node distance=5mm and 4.5mm]
% Nhánh 1: câu thật
\node[hop] (a1) {ViHSD, ViHOS, ViTHSD\\ \so{dl.goc.tong} bản ghi};
\node[hop, right=of a1] (a2) {Gộp, khử trùng\\ \so{dl.khacnhau} câu khác nhau};
\node[hop, right=of a2] (a3) {Tách kho test \so{dl.khotest} câu\\ (không gửi LLM)};
\node[hop, right=of a3] (a4) {LLM gán lại \so{dl.ganlai} câu; giữ \so{dl.that.giu} câu khớp người gán gốc};
% Nhánh 2: sinh theo kịch bản
\node[hop, below=of a1] (b1) {Kịch bản $\times$ kênh $\times$ giọng $\times$ độ lộ};
\node[hop, right=of b1] (b2) {LLM sinh \so{dl.sinh.tong} đoạn, kèm sinh đôi \nhan{AN_TOAN}};
\node[hop, right=of b2] (b3) {LLM gán lại, không biết nhãn yêu cầu};
\node[hop, right=of b3] (b4) {Giữ \so{dl.sinh.giu} đoạn trùng nhãn};
% Gộp thành train
\path (a4.east) -- (b4.east) coordinate[midway] (giua);
\node[hopdich, right=5mm of giua] (tr) {Train\\ \so{tr.tong} đoạn\\ {\normalfont(chọn cân nhãn)}};
\draw[mui] (a1) -- (a2); \draw[mui] (a2) -- (a3); \draw[mui] (a3) -- (a4);
\draw[mui] (b1) -- (b2); \draw[mui] (b2) -- (b3); \draw[mui] (b3) -- (b4);
\draw[mui] (a4.east) -- (tr.north west);
\draw[mui] (b4.east) -- (tr.south west);
% Nhánh 3: test do người viết
\node[hoptest, below=10mm of b1] (c1) {Người trong nhóm tự viết, trước khi đọc câu sinh};
\node[hoptest, right=of c1] (c2) {Test C0: \so{te.tong} đoạn, 3 giọng vùng};
\node[hoptest, right=of c2] (c3) {Sinh biến dạng C1 (\so{te.C1.n}) và C2 (\so{te.C2.n})};
\node[hoptest, right=of c3] (c4) {Kiểm tra rò rỉ: \so{rr.trunggan} cặp gần trùng với train};
\draw[-{Stealth[length=2mm]}, thick, chuxam] (c1) -- (c2);
\draw[-{Stealth[length=2mm]}, thick, chuxam] (c2) -- (c3);
\draw[-{Stealth[length=2mm]}, thick, chuxam] (c3) -- (c4);
\draw[dashed, chuxam] (tr.south) |- (c4.east);
\end{tikzpicture}
\caption{Luồng tạo dữ liệu v2. Hàng trên: câu thật; hàng giữa: hội thoại sinh; hàng dưới (xám): test do người viết.}
\label{fig:p1-luong}
\end{figure}

\subsection{Câu thật: gán lại bằng LLM}
\label{sec:p1-that}

\begin{itemize}
  \item \textbf{Gộp nguồn.} Ba bộ ViHSD (\so{dl.goc.vihsd}), ViHOS (\so{dl.goc.vihos}), ViTHSD (\so{dl.goc.vithsd})
        có tổng \so{dl.goc.tong} bản ghi. ViHOS và ViTHSD phần lớn dùng lại câu của ViHSD, nên sau khi gộp theo nội
        dung còn \so{dl.khacnhau} câu khác nhau.
  \item \textbf{Tách kho test trước khi gọi LLM.} Sau khử trùng gần đúng, các nhóm câu chưa từng nằm trong split train
        của ViHSD/ViHOS (phần ViHateT5 của v1 đã học) và chưa từng gửi LLM được tách thành kho test
        (\so{dl.khotest} câu). Kho này không bao giờ gửi cho LLM và không vào train.
  \item \textbf{Chọn câu cần gán lại.} Từ \so{dl.khotrain} câu còn lại, chọn \so{dl.ganlai} câu:
        \so{dl.ganlai.nguyhai} câu nhãn gốc nguy hại, \so{dl.ganlai.blacklist} câu sạch nhưng có từ trong blacklist
        của v1 (ca dễ báo nhầm), \so{dl.ganlai.tukhoa} câu khớp từ khóa của 3 nhãn hiếm, và \so{dl.ganlai.ngaunhien}
        câu sạch chọn ngẫu nhiên (seed 42).
  \item \textbf{Gán lại.} DeepSeek V4.1 Flash (qua OpenRouter, tắt chế độ suy nghĩ) đọc hướng dẫn gán nhãn và trả về
        nhãn kèm bằng chứng dạng JSON. Script tự kiểm nhãn hợp lệ và bỏ cụm bằng chứng không có nguyên văn trong câu.
  \item \textbf{Lọc theo đồng thuận.} Chỉ giữ câu mà nhãn LLM (xét sạch/nguy hại) khớp với ít nhất một người gán
        gốc. Giữ \so{dl.that.giu} câu, bỏ \so{dl.that.bo} câu. Phần lớn câu bị bỏ là \so{dl.that.bo.hate.antoan} câu
        gốc HATE mà LLM cho là an toàn và \so{dl.that.bo.clean.xucpham} câu gốc CLEAN mà LLM cho là xúc phạm. Nhóm
        không đủ người để xét tay từng câu mâu thuẫn, nên chọn bỏ chúng thay vì giữ một nhãn có thể sai.
\end{itemize}

\subsection{Hội thoại sinh theo kịch bản}
\label{sec:p1-sinh}

Câu thật gần như không có ba nhãn hiếm: sau khi gán lại chỉ còn \so{tr.dd.that} câu \nhan{DE_DOA},
\so{tr.qr.that} câu \nhan{QUAY_ROI} và \so{tr.ld.that} câu \nhan{LUA_DAO}. Vì vậy nhóm cho LLM sinh hội thoại theo
kịch bản (\tep{docs/v2/KICH_BAN.md}).

\begin{itemize}
  \item \textbf{Kịch bản.} 26 kịch bản xoay quanh học sinh, sinh viên: 12 lừa đảo (nạp game giá rẻ, việc làm thêm
        phải nộp cọc, giveaway giả...), 6 đe dọa (hẹn đánh, đòi nợ, tống tiền bằng ảnh...), 6 quấy rối (xin ảnh
        riêng tư, nhắn dồn dập sau khi bị từ chối...), 2 xúc phạm. Bốn kịch bản (\so{rr.kichban}) được giữ ngoài
        train, dùng làm phần ``kịch bản chưa thấy'' khi test.
  \item \textbf{Ô sinh.} Mỗi lượt gọi LLM ứng với một ô = kịch bản $\times$ kênh $\times$ giọng $\times$ độ lộ
        $\times$ kiểu lộ, và xin 5 đoạn khác nhau. Kênh: Messenger/Zalo, nhóm lớp, chat game (2--6 dòng); bình
        luận, SMS (1 dòng). Giọng: teencode, không dấu, viết tắt, giọng Bắc, Trung, Nam. Độ lộ: lộ liễu hoặc ẩn
        ý. Kiểu lộ: lộ ngay dòng đầu hoặc lộ dần ở cuối.
  \item \textbf{Sinh đôi an toàn.} Mỗi ô nguy hại có một ô \nhan{AN_TOAN} cùng kênh, cùng giọng: đoạn trông giống
        nguy hại nhưng vô hại (bạn thân đùa, nhắc trả tiền trà sữa lịch sự). Mục đích là để mô hình không học mẹo
        ``cứ có chữ dọa là nguy hại''. \nhan{XUC_PHAM} chỉ sinh dạng nhiều dòng, vì câu \nhan{XUC_PHAM} thật đều một
        dòng; như vậy mô hình không học được mẹo ``nhiều dòng là nhãn hiếm''.
  \item \textbf{An toàn nội dung.} Không có nội dung tình dục tường minh, không tên, số điện thoại, tài khoản thật;
        link là tên miền giả dạng \tep{xxx[.]com}. Ô nào LLM từ chối thì bỏ qua.
  \item \textbf{Kết quả sinh.} \so{cp.sinh.goi} lượt gọi cho \so{dl.sinh.tong} đoạn, trung bình \so{dl.sinh.dongtb}
        dòng mỗi đoạn.
  \item \textbf{Lọc lại.} LLM gán nhãn lại từng đoạn mà không biết nhãn đã yêu cầu khi sinh; chỉ giữ đoạn trùng
        nhãn: \so{dl.sinh.giu} đoạn (\so{dl.sinh.tilegiu}). \nhan{QUAY_ROI} giữ ít nhất (\so{dl.sinh.qr.tilegiu}),
        chủ yếu do kịch bản bình phẩm ngoại hình (QR3) chỉ giữ \so{dl.sinh.qrba.tilegiu}: các đoạn này hay bị gán lại
        thành \nhan{XUC_PHAM}. Ranh giới giữa hai nhãn này là chỗ hướng dẫn gán nhãn còn chưa rõ.
  \item Đoạn sinh được khử trùng gần đúng, và đoạn nào gần trùng câu test người viết thì bị loại.
\end{itemize}

\subsection{Tập train và tập test}
\label{sec:p1-test}

\begin{figure}[H]
\centering
% Sinh tự động bằng `python -m src.v2.eval.xuat_so_lieu_bao_cao`. Không sửa tay.
\begin{tikzpicture}
\begin{axis}[xbar stacked, bar width=11pt, width=0.9\linewidth, height=5.2cm,
  symbolic y coords={ld,dd,qr,xp,at}, ytick=data, yticklabels={\nhan{LUA_DAO},\nhan{DE_DOA},\nhan{QUAY_ROI},\nhan{XUC_PHAM},\nhan{AN_TOAN}},
  xmin=0, xlabel={Số mẫu train}, axis x line*=bottom, axis y line*=left, xmajorgrids,
  grid style={draw=black!12}, tick label style={font=\small}, label style={font=\small},
  scaled x ticks=false, xticklabel style={/pgf/number format/.cd, 1000 sep={.}},
  legend style={at={(0.5,1.03)}, anchor=south, legend columns=2, draw=none, font=\small,
    /tikz/every even column/.append style={column sep=0.4cm}}]
\addplot[fill=navy, draw=white, area legend] coordinates {(0,ld) (156,dd) (62,qr) (2527,xp) (3967,at)};
\addlegendentry{Câu thật (gán lại)}
\addplot[fill=maullm, draw=white, area legend] coordinates {(1751,ld) (1361,dd) (1098,qr) (470,xp) (2000,at)};
\addlegendentry{Đoạn LLM sinh}
\end{axis}
\end{tikzpicture}

\caption{Thành phần tập train theo nhãn. Ba nhãn hiếm (\nhan{QUAY_ROI}, \nhan{DE_DOA}, \nhan{LUA_DAO}) gần như toàn
là đoạn LLM sinh, nên điểm của chúng phải đo lại trên câu người viết.}
\label{fig:p1-train}
\end{figure}

\begin{table}[H]
\centering
\caption{Số mẫu theo nhãn. Train lưu bản gốc C0; test C0 do người viết}
\label{tab:p1-train}
\small
\begin{tabular}{@{}lrrrr@{}}
\toprule
\textbf{Nhãn} & \textbf{Train: câu thật} & \textbf{Train: đoạn sinh} & \textbf{Train: tổng} & \textbf{Test C0} \\
\midrule
\nhan{AN_TOAN}  & \so{tr.at.that} & \so{tr.at.sinh} & \so{tr.at} & \so{te.at} \\
\nhan{XUC_PHAM} & \so{tr.xp.that} & \so{tr.xp.sinh} & \so{tr.xp} & \so{te.xp} \\
\nhan{QUAY_ROI} & \so{tr.qr.that} & \so{tr.qr.sinh} & \so{tr.qr} & \so{te.qr} \\
\nhan{DE_DOA}   & \so{tr.dd.that} & \so{tr.dd.sinh} & \so{tr.dd} & \so{te.dd} \\
\nhan{LUA_DAO}  & \so{tr.ld.that} & \so{tr.ld.sinh} & \so{tr.ld} & \so{te.ld} \\
\midrule
Tổng & \so{tr.that} & \so{tr.sinh} & \so{tr.tong} & \so{te.tong} \\
\bottomrule
\end{tabular}
\end{table}

\textbf{Train} có \so{tr.tong} đoạn. Ba nhãn hiếm gần như toàn là đoạn sinh, nên điểm thật của chúng chỉ đo được
trên test. Câu thật đều một dòng (\so{tr.motdong} mẫu một dòng), còn đoạn sinh phần lớn nhiều dòng; phần kết quả
kiểm tra mô hình không đoán nhãn theo số dòng.

\textbf{Test} có \so{te.tong} đoạn do \so{te.nguoiviet} người trong nhóm tự viết theo
\tep{docs/v2/HUONG_DAN_VIET_TEST.md}, viết xong trước khi đọc bất kỳ đoạn LLM sinh nào để không vô tình bắt chước.
\begin{itemize}
  \item \textbf{Giọng vùng miền:} Bắc \so{te.bac}, Trung \so{te.trung}, Nam \so{te.nam} đoạn, viết đúng cách người
        mỗi miền nhắn (``tau nỏ cần mô'', ``tui hổng cần đâu'').
  \item \textbf{Kịch bản:} \so{te.trongpp} đoạn thuộc kịch bản đã sinh vào train; \so{te.chuathay} đoạn thuộc 4 kịch
        bản giữ ngoài train và các tình huống mới do người viết tự nghĩ. Phần chưa thấy chỉ có \so{te.chuathay.xp}
        đoạn \nhan{XUC_PHAM}.
  \item \textbf{Nhãn:} người viết gán nhãn và bằng chứng; một LLM gán độc lập để đối chiếu. Hai bên khớp ở
        \so{te.chon.trung} đoạn. Ở các đoạn lệch, người duyệt giữ nhãn người viết cho \so{te.chon.nguoi} đoạn và đổi
        theo LLM cho \so{te.chon.llm} đoạn (ghi \tep{cach_gan=llm_nguoi_duyet}). Test chưa có lượt gán chéo giữa hai
        người nên chưa có kappa.
\end{itemize}

\textbf{Bản biến dạng C1, C2.} Mỗi đoạn test được biến dạng thành 2 bản cho mỗi nhóm, kiểu bốc ngẫu nhiên (seed 42),
được \so{te.C1.n} mẫu C1 và \so{te.C2.n} mẫu C2. Biến dạng làm theo từng từ, khoảng \so{te.C1.tiletu} số từ, luôn
trúng cụm bằng chứng, và vị trí bằng chứng đi theo từ bị đổi.

\begin{table}[H]
\centering
\caption{Các kiểu biến dạng trong test (số mẫu)}
\label{tab:p1-bien-dang}
\small
\begin{tabularx}{\linewidth}{@{}p{3.6cm}Y@{}}
\toprule
\textbf{Nhóm} & \textbf{Kiểu} \\
\midrule
C1, đã thấy (xuất hiện tự nhiên trong train) & bỏ dấu (\so{te.C1.bodau}), teencode (\so{te.C1.teencode}),
  gõ sai (\so{te.C1.gosai}), kéo dài ký tự (\so{te.C1.keodai}) \\
C1, giữ ngoài train & dính từ (\so{te.C1.dinhtu}), số thay chữ (\so{te.C1.sothaychu}),
  chèn ký tự (\so{te.C1.chenkytu}) \\
C2, điểm mù, không có trong train & chữ toàn độ rộng (\so{te.C2.fullwidth}), ký tự vô hình (\so{te.C2.zerowidth}),
  chèn dấu cách giữa các chữ (\so{te.C2.spaceinjection}), ký tự phân cách lạ (\so{te.C2.exoticseparator}),
  chữ Cyrillic giống chữ Latin (\so{te.C2.cyrillichomoglyph}), kéo dài kép (\so{te.C2.doubleelongation}) \\
\bottomrule
\end{tabularx}
\end{table}

\textbf{Chống rò rỉ.} Script \tep{src.v2.data.kiem_tra_ro_ri} so test với train và kho test: \so{rr.trunggan} cặp gần
trùng (điểm trùng 3-gram từ \so{rr.nguong} trở lên, so cả đoạn và từng dòng), \so{rr.trainkichban} câu train thuộc
4 kịch bản giữ ngoài, không trùng id. Các kiểu C2 và nhóm C1 giữ ngoài không được dùng khi tăng cường dữ liệu lúc
train.

\subsection{Chất lượng nhãn train: điều chưa làm được}
\label{sec:p1-chatluong}

Train chủ yếu do LLM gán hoặc sinh, nên cần ước tỉ lệ nhãn sai. Kế hoạch là hai người gán độc lập mẫu
\so{dt.mau} câu (một nửa câu được giữ, một nửa câu bị bỏ, che nhãn). \textbf{Thực tế mẫu này do một mô hình thứ
hai (Claude) gán, chưa phải người.} Vì vậy chưa có kappa giữa người, và các số sau chỉ là mức đồng ý giữa hai mô hình:

\begin{itemize}
  \item Ở nhóm câu được giữ, mô hình thứ hai đồng ý với nhãn của DeepSeek ở \so{dt.giu.dung}/\so{dt.giu.tong} câu
        (\so{dt.giu.tile}, khoảng tin cậy 95\%: \so{dt.giu.lo} -- \so{dt.giu.hi}). Như vậy vẫn còn một phần nhãn train
        có thể sai.
  \item Ở nhóm câu bị bỏ, mô hình thứ hai đứng về phía LLM (thay vì nhãn gốc) trong \so{dt.bo.llm} trường hợp: bỏ
        câu mâu thuẫn cũng làm mất một số câu mà LLM gán đúng.
\end{itemize}
Do đó điểm trên test người viết là thước đo chính, và nhãn train nhiễu được ghi là một hạn chế.

\subsection{Chi phí tạo dữ liệu}
\label{sec:p1-chiphi}

\begin{table}[H]
\centering
\caption{Chi phí gọi LLM để tạo dữ liệu train (DeepSeek V4.1 Flash, tính từ số token $\times$ giá endpoint)}
\label{tab:p1-chiphi}
\small
\begin{tabular}{@{}lrrr@{}}
\toprule
\textbf{Bước} & \textbf{Lượt gọi} & \textbf{USD} & \textbf{Phút} \\
\midrule
Gán lại câu thật & \so{cp.ganlai.goi} & \so{cp.ganlai.usd} & \so{cp.ganlai.phut} \\
Sinh hội thoại & \so{cp.sinh.goi} & \so{cp.sinh.usd} & \so{cp.sinh.phut} \\
Lọc lại đoạn sinh & \so{cp.loc.goi} & \so{cp.loc.usd} & \so{cp.loc.phut} \\
\midrule
Tổng & \so{cp.tong.goi} & \so{cp.tong.usd} & \so{cp.tong.phut} \\
\bottomrule
\end{tabular}
\end{table}

Toàn bộ phần dữ liệu do LLM tạo và gán tốn \so{cp.tong.usd} USD (khoảng \so{cp.tong.vnd} đồng theo tỷ giá giả định
ghi trong \tep{chi_phi.json}), khoảng \so{cp.tong.phut} phút chạy và \so{cp.tong.trieutoken} triệu token. Chi phí
thấp nên có thể sinh lại hoặc mở rộng kịch bản khi xuất hiện kiểu lừa đảo mới.

% Phần 2: Mô hình v2, cách đánh giá, kết quả chung, Yêu cầu 1.
% Mọi con số lấy qua \so{...} (so_lieu.tex, sinh tự động). Không gõ số tay.

\section{Mô hình v2}
\label{sec:p2-mo-hinh}

\subsection{Vì sao chọn encoder thay cho mô hình sinh}

Kế hoạch ban đầu là fine-tune một mô hình sinh nhỏ (Qwen3-1.7B, QLoRA) để viết ra JSON. Khi làm thật, nhóm huấn luyện
song song 4 ứng viên trên cùng dữ liệu và cùng seed, rồi chọn trên tập dev (so sánh đầy đủ ở phần Yêu cầu 3). Ứng viên
dựa trên LLM 0,6B (GuardMerge) thua encoder có ý nghĩa thống kê (kiểm định McNemar, $p = \so{dev.mc.guard}$) và chậm
gấp hơn hai lần (\so{dev.guard.p50} ms so với \so{dev.mmbert.p50} ms mỗi câu trên GPU). Vì vậy v2 dùng encoder.

\subsection{Cấu trúc}

Mô hình chốt có tên nội bộ \textbf{SafeVi-DualORC-CRF}, xây trên \textbf{mmBERT-base} (encoder đa ngôn ngữ). Nó đọc
cả đoạn hội thoại một lần (tối đa \so{v2.maxtok} token), không chỉ dòng cuối, và có hai đầu ra:

\begin{enumerate}
  \item \textbf{Đầu phân loại:} trả xác suất cho 5 nhãn. Nhãn có xác suất cao nhất là kết quả.
  \item \textbf{Đầu gắn thẻ:} gán cho từng token một thẻ B (bắt đầu cụm), I (trong cụm) hoặc O (ngoài cụm), dùng thêm
        tầng CRF để các thẻ nối với nhau hợp lệ. Các đoạn B--I liền nhau là cụm bằng chứng.
\end{enumerate}

Vì cụm bằng chứng được \emph{cắt ra từ chính văn bản}, nó luôn khớp nguyên văn. Ngoài encoder, mô hình có thêm một nhánh CNN đọc từng byte của văn bản, nhằm nhận ra chữ
bị thay bằng ký tự lạ mà bộ tách từ không quen. Hai đầu ra được ghép thành JSON \texttt{\{"nhan": ..., "bang\_chung":
[...]\}} đúng định dạng đã nêu ở phần nhãn.

\begin{figure}[H]
\centering
\begin{tikzpicture}[
  font=\small,
  hop/.style={draw=navy, rounded corners=2pt, fill=nenxanh, align=center, text width=2.5cm, minimum height=1.1cm},
  tuychon/.style={hop, dashed, fill=nenxam},
  mui/.style={-{Stealth[length=2mm]}, thick, navy}
]
  \node[hop] (vao) {Đoạn hội thoại\\(1--8 dòng)};
  \node[tuychon, right=0.5cm of vao] (chuan) {Chuẩn hóa\\(tùy chọn)};
  \node[hop, right=0.5cm of chuan, text width=2.7cm] (enc) {mmBERT-base\\+ CNN đọc byte};
  \node[hop, right=0.7cm of enc, yshift=0.8cm, text width=3.0cm] (pl) {Đầu phân loại\\xác suất 5 nhãn};
  \node[hop, right=0.7cm of enc, yshift=-0.8cm, text width=3.0cm] (gt) {Đầu gắn thẻ\\BIO + CRF $\to$ cụm};
  % Hộp JSON đặt hẳn dưới đầu gắn thẻ, mép phải thẳng hàng, để không chồng lên hộp phía trên
  \node[hop, anchor=north east, text width=4.6cm] (ra) at ([yshift=-0.9cm]gt.south east)
    {JSON: \texttt{nhan} + \texttt{bang\_chung}};
  \draw[mui] (vao) -- (chuan);
  \draw[mui] (chuan) -- (enc);
  \draw[mui] (enc.east) -- ++(0.3,0) |- (pl.west);
  \draw[mui] (enc.east) -- ++(0.3,0) |- (gt.west);
  \draw[mui] (pl.east) -- ++(0.4,0) |- (ra.east);
  \draw[mui] (gt.east) -- ++(0.4,0) |- (ra.east);
\end{tikzpicture}
\caption{Luồng xử lý của v2. Hai đầu ra dùng chung encoder nhưng có ngưỡng và cách chấm riêng.}
\label{fig:p2-luong}
\end{figure}

\subsection{Huấn luyện}

Tệp train lưu bản gốc C0. Khi huấn luyện, mỗi câu được biến dạng ngẫu nhiên ngay lúc đọc: bỏ dấu với xác suất
\so{v2.accent}, biến dạng bề mặt thường gặp (kiểu C1 đã thấy) với xác suất \so{v2.surface}. Mô hình còn được phạt nếu dự
đoán trên bản gốc và bản biến dạng khác nhau, và được huấn luyện đối kháng FGM như v1 (cộng nhiễu nhỏ vào embedding,
$\varepsilon = \so{v2.fgm}$). Các kiểu điểm mù C2 và nhóm C1 giữ ngoài train không được dùng ở bước này.
Huấn luyện với \so{v2.epoch} epoch, seed 42.

\subsection{Chia dữ liệu và các ngưỡng}

Tập train được chia thành ba phần: \textbf{fit} (\so{dev.chia.fit} mẫu) để học, \textbf{dev} (\so{dev.chia.dev} mẫu) để
chọn mô hình và checkpoint, \textbf{cal} (\so{dev.chia.cal} mẫu) để hiệu chỉnh. Dev và cal tách từ train nên nhãn của
chúng cũng do LLM gán; không câu test nào được dùng ở đây. Trên cal, nhóm dò và khóa ba tham số:

\begin{itemize}
  \item nhiệt độ $T = \so{v2.T}$ để xác suất phân loại sát với tỉ lệ đúng thật (chia logit cho $T$ trước softmax);
  \item ngưỡng bằng chứng \so{v2.nguongbc}: token có xác suất thuộc cụm từ ngưỡng này trở lên mới được tô;
  \item ngưỡng riêng \so{v2.nguongmottu} cho cụm chỉ có một từ, cao hơn để tránh tô lẻ tẻ.
\end{itemize}

Cả ba được ghi vào \tep{calibration.json} và đóng băng trước khi chạy test.

% ---------------------------------------------------------------------------
\section{Cách đánh giá và kết quả chung}
\label{sec:p2-danh-gia}

\subsection{Ba hệ được so sánh}

\begin{itemize}
  \item \textbf{v1:} SafeViChi vòng loại (ViHateT5, HATE/CLEAN). Chỉ có nhãn nhị phân, không có bằng chứng, nên chỉ chấm
        được phần nguy hại/an toàn. HATE được tính là ``nguy hại'', CLEAN là \nhan{AN_TOAN}.
  \item \textbf{gpt-oss-120b qua API:} LLM lớn, zero-shot. Prompt chính là hướng dẫn gán nhãn của nhóm (định nghĩa 5 nhãn,
        thứ tự quyết định, ví dụ, quy tắc bằng chứng). Mức suy luận \texttt{\so{api.reasoning}}.
  \item \textbf{v2:} mô hình ở Mục~\ref{sec:p2-mo-hinh}, không bật chuẩn hóa đầu vào.
\end{itemize}

Cả ba chạy trên cùng bộ test: C0 (\so{C0.v2.n} đoạn gốc do người viết), C1 (\so{C1.v2.n} bản biến dạng thường gặp như
bỏ dấu, teencode, gõ sai) và C2 (\so{C2.v2.n} bản biến dạng ``điểm mù'' chưa từng xuất hiện khi train v2).

\subsection{Độ đo}

\begin{itemize}
  \item \textbf{Macro-F1:} trung bình F1 của 5 nhãn, mỗi nhãn quan trọng như nhau. Đo khả năng gán nhãn đúng tình huống
  \item \textbf{Binary-F1:} gộp 4 nhãn nguy hại thành một, đo khả năng phân biệt nguy hại với an toàn.
  \item \textbf{FPR:} tỉ lệ câu \nhan{AN_TOAN} bị báo nguy hại (báo nhầm). Càng thấp càng tốt.
  \item \textbf{BC P/R/F1 theo ý:} chấm bằng chứng. Mỗi cụm người gán là một ý, precision là phần cụm mô hình tô
        trúng ý đúng, recall là phần ý đúng được tô ra.
\end{itemize}

Trong toàn báo cáo, F1, precision và recall được nhân 100 (ví dụ 86,5); các tỉ lệ ghi kèm dấu \%. Khoảng tin cậy 95\%
tính bằng bootstrap 1.000 lần, seed 42.

\subsection{Kết quả}

\begin{table}[H]
\centering
\small
\caption{Kết quả chính trên bộ test. ``--'': hệ không có đầu ra tương ứng.}
\label{tab:p2-chinh}
\begin{tabular}{llccccc}
\toprule
Tập & Hệ & Macro-F1 & KTC 95\% Macro-F1 & Binary-F1 & FPR & BC F1 theo ý \\
\midrule
\multirow{3}{*}{C0} & v1 & -- & -- & \so{C0.v1.bin} & \so{C0.v1.fpr} & -- \\
 & gpt-oss-120b & \so{C0.llm.macro} & \so{C0.llm.macro.lo} -- \so{C0.llm.macro.hi} & \so{C0.llm.bin} & \so{C0.llm.fpr} & \so{C0.llm.bcf} \\
 & \textbf{v2} & \textbf{\so{C0.v2.macro}} & \so{C0.v2.macro.lo} -- \so{C0.v2.macro.hi} & \textbf{\so{C0.v2.bin}} & \so{C0.v2.fpr} & \textbf{\so{C0.v2.bcf}} \\
\midrule
\multirow{3}{*}{C1} & v1 & -- & -- & \so{C1.v1.bin} & \so{C1.v1.fpr} & -- \\
 & gpt-oss-120b & \so{C1.llm.macro} & \so{C1.llm.macro.lo} -- \so{C1.llm.macro.hi} & \so{C1.llm.bin} & \so{C1.llm.fpr} & \so{C1.llm.bcf} \\
 & \textbf{v2} & \textbf{\so{C1.v2.macro}} & \so{C1.v2.macro.lo} -- \so{C1.v2.macro.hi} & \textbf{\so{C1.v2.bin}} & \so{C1.v2.fpr} & \textbf{\so{C1.v2.bcf}} \\
\midrule
\multirow{3}{*}{C2} & v1 & -- & -- & \so{C2.v1.bin} & \so{C2.v1.fpr} & -- \\
 & gpt-oss-120b & \so{C2.llm.macro} & \so{C2.llm.macro.lo} -- \so{C2.llm.macro.hi} & \so{C2.llm.bin} & \so{C2.llm.fpr} & \so{C2.llm.bcf} \\
 & \textbf{v2} & \textbf{\so{C2.v2.macro}} & \so{C2.v2.macro.lo} -- \so{C2.v2.macro.hi} & \textbf{\so{C2.v2.bin}} & \so{C2.v2.fpr} & \textbf{\so{C2.v2.bcf}} \\
\bottomrule
\end{tabular}
\end{table}

\begin{figure}[H]
\centering
% Sinh tự động bằng `python -m src.v2.eval.xuat_so_lieu_bao_cao`. Không sửa tay.
\begin{tikzpicture}
\begin{axis}[ybar=3pt, bar width=18pt, width=\linewidth, height=6cm, ymin=0, ymax=100,
  symbolic x coords={C0,C1,C2}, xtick=data, enlarge x limits=0.25, ylabel={Binary-F1},
  axis x line*=bottom, axis y line*=left, ymajorgrids, grid style={draw=black!12},
  tick label style={font=\small}, label style={font=\small},
  nodes near coords={\pgfmathprintnumber[fixed,fixed zerofill,precision=1,use comma]{\pgfplotspointmeta}},
  every node near coord/.append style={font=\scriptsize, color=black!75},
  legend style={at={(0.5,1.03)}, anchor=south, legend columns=3, draw=none, font=\small,
    /tikz/every even column/.append style={column sep=0.4cm}}]
\addplot[fill=mauvmot, draw=white, area legend] coordinates {(C0,31.2) (C1,31.0) (C2,39.5)};
\addlegendentry{SafeViChi v1}
\addplot[fill=maullm, draw=white, area legend] coordinates {(C0,89.2) (C1,74.5) (C2,77.9)};
\addlegendentry{gpt-oss-120b (API)}
\addplot[fill=mauvhai, draw=white, area legend] coordinates {(C0,92.9) (C1,90.3) (C2,90.9)};
\addlegendentry{SafeViChi v2}
\end{axis}
\end{tikzpicture}

\caption{Binary-F1 (nguy hại/an toàn) của ba hệ trên C0, C1, C2. Số đầy đủ ở Bảng~\ref{tab:p2-chinh}.}
\label{fig:p2-binary}
\end{figure}

\textbf{Nhận xét.}
\begin{itemize}
  \item \textbf{v2 tốt nhất trên cả ba tập.} Khoảng tin cậy 95\% của Macro-F1 giữa v2 và gpt-oss-120b không chồng nhau ở
        cả C0, C1 và C2, nên khoảng cách không phải do may rủi của mẫu.
  \item \textbf{v2 giữ điểm khi chữ bị biến dạng.} Binary-F1 của v2 đi từ \so{C0.v2.bin} (C0) xuống \so{C2.v2.bin} (C2);
        của gpt-oss-120b từ \so{C0.llm.bin} xuống \so{C2.llm.bin}.
  % \item \textbf{v1 gần như chỉ bắt được xúc phạm.} Ở C0, v1 báo HATE cho \so{C0.v1.xp.r} câu \nhan{XUC_PHAM}, nhưng chỉ
  %       \so{C0.v1.dd.r} câu \nhan{DE_DOA}, \so{C0.v1.qr.r} câu \nhan{QUAY_ROI} và \so{C0.v1.ld.r} câu \nhan{LUA_DAO}.
  %       Precision \so{C0.v1.p} của v1 chỉ phản ánh việc nó hiếm khi báo. Đây là lý do phải đổi bộ nhãn.
  \item \textbf{Đổi lại, v2 báo nhầm nhiều hơn.} Ở C0, FPR của v2 là \so{C0.v2.fpr}, của gpt-oss-120b là
        \so{C0.llm.fpr}. gpt-oss-120b thận trọng hơn: precision nguy hại \so{C0.llm.p} nhưng recall chỉ \so{C0.llm.r}
        (v2: \so{C0.v2.p} và \so{C0.v2.r}).
\end{itemize}

% ---------------------------------------------------------------------------
\section{Yêu cầu 1: giảm kiểm duyệt sai và giải thích sai}
\label{sec:p2-yc1}

\subsection{Tách kết quả phân loại khỏi bằng chứng}

Ở v1, bằng chứng tìm bằng occlusion \emph{sau} khi đã phân loại: che từng cụm rồi xem điểm giảm bao nhiêu. Cách này phải
chạy lại mô hình nhiều lần và chất lượng bằng chứng gắn chặt với bộ phân loại. v2 tách hẳn hai việc:

\begin{itemize}
  \item \textbf{Hai đầu ra riêng}, mỗi đầu có hàm mất mát riêng khi train.
  \item \textbf{Hai ngưỡng riêng:} nhiệt độ $T$ cho xác suất nhãn; ngưỡng \so{v2.nguongbc} và \so{v2.nguongmottu} cho
        cụm bằng chứng. Đổi ngưỡng bằng chứng không làm đổi nhãn.
  \item \textbf{Hai bộ độ đo riêng:} nhãn chấm bằng P/R/F1 theo nhãn; bằng chứng chấm bằng P/R/F1 theo ý.
  \item \textbf{Vai trò khác nhau trong ứng dụng:} nhãn quyết định ứng dụng làm gì (làm mờ, cảnh báo, gợi ý chặn); bằng
        chứng chỉ để người dùng hiểu vì sao, nên chỉ có nghĩa khi nhãn là nguy hại.
\end{itemize}

\subsection{Kết quả theo từng nhãn}

\begin{table}[H]
\centering
\small
\caption{Kết quả theo nhãn trên C0. n: số câu có nhãn đúng là nhãn đó.}
\label{tab:p2-theo-nhan}
\begin{tabular}{lcccccccc}
\toprule
 & & \multicolumn{4}{c}{v2} & \multicolumn{2}{c}{gpt-oss-120b} \\
\cmidrule(lr){3-6}\cmidrule(lr){7-8}
Nhãn & n & P & R & F1 & BC F1 & F1 & BC F1 \\
\midrule
\nhan{AN_TOAN}  & \so{C0.v2.at.n} & \so{C0.v2.at.p} & \so{C0.v2.at.r} & \so{C0.v2.at.f} & -- & \so{C0.llm.at.f} & -- \\
\nhan{XUC_PHAM} & \so{C0.v2.xp.n} & \so{C0.v2.xp.p} & \so{C0.v2.xp.r} & \so{C0.v2.xp.f} & \so{C0.v2.xp.bcf} & \so{C0.llm.xp.f} & \so{C0.llm.xp.bcf} \\
\nhan{QUAY_ROI} & \so{C0.v2.qr.n} & \so{C0.v2.qr.p} & \so{C0.v2.qr.r} & \so{C0.v2.qr.f} & \so{C0.v2.qr.bcf} & \so{C0.llm.qr.f} & \so{C0.llm.qr.bcf} \\
\nhan{DE_DOA}   & \so{C0.v2.dd.n} & \so{C0.v2.dd.p} & \so{C0.v2.dd.r} & \so{C0.v2.dd.f} & \so{C0.v2.dd.bcf} & \so{C0.llm.dd.f} & \so{C0.llm.dd.bcf} \\
\nhan{LUA_DAO}  & \so{C0.v2.ld.n} & \so{C0.v2.ld.p} & \so{C0.v2.ld.r} & \so{C0.v2.ld.f} & \so{C0.v2.ld.bcf} & \so{C0.llm.ld.f} & \so{C0.llm.ld.bcf} \\
\bottomrule
\end{tabular}
\end{table}

\nhan{LUA_DAO} và \nhan{XUC_PHAM} là hai nhãn tốt nhất. \nhan{QUAY_ROI} yếu nhất, chủ yếu do recall thấp: mô hình bỏ sót
khá nhiều câu quấy rối. So với gpt-oss-120b, v2 hơn rõ ở \nhan{XUC_PHAM}, \nhan{QUAY_ROI}, \nhan{DE_DOA} và ngang ở
\nhan{LUA_DAO}; bằng chứng của v2 tốt hơn ở cả bốn nhãn nguy hại.

\begin{figure}[H]
\centering
% Sinh tự động bằng `python -m src.v2.eval.xuat_so_lieu_bao_cao`. Không sửa tay.
\begin{tikzpicture}
\fill[mauvmot!90] (0.0,0.0) rectangle ++(1.45,-1.45);\node[font=\small\bfseries,text=white] at (0.725,-0.725) {86};
\fill[mauvmot!0] (1.45,0.0) rectangle ++(1.45,-1.45);\node[font=\small\bfseries,text=black] at (2.175,-0.725) {0};
\fill[mauvmot!1] (2.9,0.0) rectangle ++(1.45,-1.45);\node[font=\small\bfseries,text=black] at (3.625,-0.725) {1};
\fill[mauvmot!6] (4.35,0.0) rectangle ++(1.45,-1.45);\node[font=\small\bfseries,text=black] at (5.074999999999999,-0.725) {6};
\fill[mauvmot!3] (5.8,0.0) rectangle ++(1.45,-1.45);\node[font=\small\bfseries,text=black] at (6.5249999999999995,-0.725) {3};
\node[anchor=east,font=\footnotesize] at (-0.1,-0.725) {\nhan{AN_TOAN}};
\fill[mauvmot!4] (0.0,-1.45) rectangle ++(1.45,-1.45);\node[font=\small\bfseries,text=black] at (0.725,-2.175) {2};
\fill[mauvmot!91] (1.45,-1.45) rectangle ++(1.45,-1.45);\node[font=\small\bfseries,text=white] at (2.175,-2.175) {48};
\fill[mauvmot!6] (2.9,-1.45) rectangle ++(1.45,-1.45);\node[font=\small\bfseries,text=black] at (3.625,-2.175) {3};
\fill[mauvmot!0] (4.35,-1.45) rectangle ++(1.45,-1.45);\node[font=\small\bfseries,text=black] at (5.074999999999999,-2.175) {0};
\fill[mauvmot!0] (5.8,-1.45) rectangle ++(1.45,-1.45);\node[font=\small\bfseries,text=black] at (6.5249999999999995,-2.175) {0};
\node[anchor=east,font=\footnotesize] at (-0.1,-2.175) {\nhan{XUC_PHAM}};
\fill[mauvmot!16] (0.0,-2.9) rectangle ++(1.45,-1.45);\node[font=\small\bfseries,text=black] at (0.725,-3.625) {8};
\fill[mauvmot!0] (1.45,-2.9) rectangle ++(1.45,-1.45);\node[font=\small\bfseries,text=black] at (2.175,-3.625) {0};
\fill[mauvmot!71] (2.9,-2.9) rectangle ++(1.45,-1.45);\node[font=\small\bfseries,text=white] at (3.625,-3.625) {36};
\fill[mauvmot!14] (4.35,-2.9) rectangle ++(1.45,-1.45);\node[font=\small\bfseries,text=black] at (5.074999999999999,-3.625) {7};
\fill[mauvmot!0] (5.8,-2.9) rectangle ++(1.45,-1.45);\node[font=\small\bfseries,text=black] at (6.5249999999999995,-3.625) {0};
\node[anchor=east,font=\footnotesize] at (-0.1,-3.625) {\nhan{QUAY_ROI}};
\fill[mauvmot!15] (0.0,-4.35) rectangle ++(1.45,-1.45);\node[font=\small\bfseries,text=black] at (0.725,-5.074999999999999) {9};
\fill[mauvmot!2] (1.45,-4.35) rectangle ++(1.45,-1.45);\node[font=\small\bfseries,text=black] at (2.175,-5.074999999999999) {1};
\fill[mauvmot!2] (2.9,-4.35) rectangle ++(1.45,-1.45);\node[font=\small\bfseries,text=black] at (3.625,-5.074999999999999) {1};
\fill[mauvmot!82] (4.35,-4.35) rectangle ++(1.45,-1.45);\node[font=\small\bfseries,text=white] at (5.074999999999999,-5.074999999999999) {51};
\fill[mauvmot!0] (5.8,-4.35) rectangle ++(1.45,-1.45);\node[font=\small\bfseries,text=black] at (6.5249999999999995,-5.074999999999999) {0};
\node[anchor=east,font=\footnotesize] at (-0.1,-5.074999999999999) {\nhan{DE_DOA}};
\fill[mauvmot!6] (0.0,-5.8) rectangle ++(1.45,-1.45);\node[font=\small\bfseries,text=black] at (0.725,-6.5249999999999995) {4};
\fill[mauvmot!0] (1.45,-5.8) rectangle ++(1.45,-1.45);\node[font=\small\bfseries,text=black] at (2.175,-6.5249999999999995) {0};
\fill[mauvmot!0] (2.9,-5.8) rectangle ++(1.45,-1.45);\node[font=\small\bfseries,text=black] at (3.625,-6.5249999999999995) {0};
\fill[mauvmot!0] (4.35,-5.8) rectangle ++(1.45,-1.45);\node[font=\small\bfseries,text=black] at (5.074999999999999,-6.5249999999999995) {0};
\fill[mauvmot!94] (5.8,-5.8) rectangle ++(1.45,-1.45);\node[font=\small\bfseries,text=white] at (6.5249999999999995,-6.5249999999999995) {68};
\node[anchor=east,font=\footnotesize] at (-0.1,-6.5249999999999995) {\nhan{LUA_DAO}};
\node[anchor=west,rotate=40,font=\footnotesize] at (0.5249999999999999,0.15) {\nhan{AN_TOAN}};
\node[anchor=west,rotate=40,font=\footnotesize] at (1.9749999999999999,0.15) {\nhan{XUC_PHAM}};
\node[anchor=west,rotate=40,font=\footnotesize] at (3.425,0.15) {\nhan{QUAY_ROI}};
\node[anchor=west,rotate=40,font=\footnotesize] at (4.874999999999999,0.15) {\nhan{DE_DOA}};
\node[anchor=west,rotate=40,font=\footnotesize] at (6.324999999999999,0.15) {\nhan{LUA_DAO}};
\draw[black!25] (0,0) grid[step=1.45] (7.25,-7.25);
\node[font=\small,rotate=90] at (-3.0,-3.625) {Nhãn đúng};
\node[font=\small] at (3.625,2.0) {Nhãn SafeViChi v2 dự đoán};
\end{tikzpicture}

\caption{Ma trận nhầm lẫn của v2 trên C0. Số trong ô là số câu; màu đậm theo tỉ lệ trong hàng. Ngoài đường chéo,
lỗi dồn vào cột \nhan{AN_TOAN} (bỏ sót) và cột \nhan{DE_DOA} (báo nhầm, lẫn quấy rối với đe dọa).}
\label{fig:p2-nham-lan}
\end{figure}

Ngoài nhầm với \nhan{AN_TOAN}, lỗi lớn nhất là lẫn \nhan{QUAY_ROI} thành \nhan{DE_DOA} (\so{C0.v2.cm.qr.dd} câu).
Hai nhãn này giáp nhau thật: câu bám riết thường kèm một lời ép (``không gửi ảnh thì\dots''), nên ranh giới giữa
``quấy rối'' và ``dọa'' khó cả với người gán.

\subsection{Báo nhầm câu an toàn}

Ở C0, \so{C0.v2.fp.tong} trên \so{C0.at.n} câu \nhan{AN_TOAN} bị báo nguy hại: \so{C0.v2.cm.at.dd} câu thành
\nhan{DE_DOA}, \so{C0.v2.cm.at.ld} câu thành \nhan{LUA_DAO}, \so{C0.v2.cm.at.qr} câu thành \nhan{QUAY_ROI}. Xem các câu
này, mô hình thường bám vào \emph{từ khóa của tình huống nguy hại} mà bỏ qua ý của cả câu: chữ ``OTP'', ``phí ship'',
``báo cô'', ``chặn đường''. Bảng~\ref{tab:p2-bao-nham} là hai ví dụ.

\begin{table}[H]
\centering
\small
\caption{Hai câu an toàn bị v2 báo nhầm (C0).}
\label{tab:p2-bao-nham}
\begin{tabularx}{\linewidth}{>{\hsize=1.45\hsize}Y >{\hsize=0.55\hsize}Y}
\toprule
Tin nhắn & Kết quả \\
\midrule
\so{vd.otp.text} &
Đúng: \so{vd.otp.gold}\newline v2: \so{vd.otp.pred} (\so{vd.otp.tincay})\newline gpt-oss-120b: \so{vd.otp.llm}\newline v2 tô: \so{vd.otp.bc} \\
\midrule
\so{vd.baoco.text} &
Đúng: \so{vd.baoco.gold}\newline v2: \so{vd.baoco.pred} (\so{vd.baoco.tincay})\newline gpt-oss-120b: \so{vd.baoco.llm}\newline v2 tô: \so{vd.baoco.bc} \\
\bottomrule
\end{tabularx}
\end{table}

Câu đầu là tin OTP thật: nó không đòi gì, chỉ bảo nhập mã vào ứng dụng. Hướng dẫn gán nhãn ghi rõ loại tin này là
\nhan{AN_TOAN}, nhưng trong train, chữ ``OTP'' gần như chỉ xuất hiện ở đoạn lừa đảo. Câu sau là bạn bè khuyên nhau
\emph{báo cô giáo}, tức là cách xử lý đúng khi bị bắt nạt; mô hình lại thấy ``gửi cho cô'', ``gặp nhóm kia'' và đoán
dọa. Cả hai đều bị báo với độ tin cậy cao, nên chỉ hạ ngưỡng chung không đủ sửa.

\textbf{Tiếng lóng.} Câu tiếng lóng dễ bị báo nhầm vì chứa từ tục mà không nhắm ai. Trên dev, ở lát câu có thẻ tiếng
lóng (n = \so{dev.tieng.n}), mô hình chốt có FPR \so{dev.mmbert.tieng.fpr} ở lát này, ngang BamiBERT
(\so{dev.bamibert.tieng.fpr}), thấp hơn GuardMerge (\so{dev.guard.tieng.fpr}) nhưng cao hơn ViSoBERT
(\so{dev.visobert.tieng.fpr}). Bù lại, Macro-F1 của nó trên lát này cao nhất trong 4 ứng viên (\so{dev.mmbert.tieng} so với
\so{dev.visobert.tieng} của ViSoBERT), nên tổng thể vẫn là một lý do chọn mmBERT-base. Dù vậy, mức báo nhầm này vẫn cao.

\subsection{Bỏ sót câu nguy hại}

Ở C0, \so{C0.v2.fn.tong} câu nguy hại bị đoán \nhan{AN_TOAN}: \so{C0.v2.cm.dd.at} câu \nhan{DE_DOA}, \so{C0.v2.cm.qr.at}
câu \nhan{QUAY_ROI}, \so{C0.v2.cm.ld.at} câu \nhan{LUA_DAO} và \so{C0.v2.cm.xp.at} câu \nhan{XUC_PHAM}. Hơn nửa là câu
mang tính ẩn ý: lời dọa nói tránh (``mai đi đường đông người nhé''), kiểm soát người yêu bằng giọng quan tâm. Phần lớn
thuộc kịch bản chưa có trong train; ví dụ cụ thể ở phần Yêu cầu 2.

\subsection{Hiệu chỉnh ngưỡng: đã làm gì, chưa làm gì}

\begin{itemize}
  \item \textbf{Đã làm:} dò nhiệt độ $T$ và hai ngưỡng bằng chứng trên tập cal, khóa trước khi chạy test
        (Mục~\ref{sec:p2-mo-hinh}).
  \item \textbf{Chưa làm:} dò ngưỡng riêng cho từng nhãn (ví dụ đòi xác suất \nhan{DE_DOA} cao hơn mới báo). Đây là cách
        trực tiếp nhất để hạ FPR, và là việc tiếp theo. Cần một tập validation do người gán; dò trên cal (nhãn LLM) có
        nguy cơ khớp theo lỗi của LLM.
\end{itemize}

Kết quả thực nghiệm: FPR của v2 ở C0 (\so{C0.v2.fpr}) cao hơn gpt-oss-120b (\so{C0.llm.fpr}). Theo vùng, FPR cao nhất ở giọng Nam
(\so{C0.nam.v2.fpr}), so với Bắc \so{C0.bac.v2.fpr} và Trung \so{C0.trung.v2.fpr}. Mỗi vùng chỉ có khoảng
\so{C0.nam.v2.at.n}--\so{C0.bac.v2.at.n} câu an toàn, nên một câu sai làm FPR của vùng đổi vài điểm phần trăm; chênh lệch
vùng cần đọc kèm cỡ mẫu này.

\subsection{Chất lượng bằng chứng}

\begin{figure}[H]
\centering
% Sinh tự động bằng `python -m src.v2.eval.xuat_so_lieu_bao_cao`. Không sửa tay.
\begin{tikzpicture}
\begin{axis}[width=0.75\linewidth, height=5.8cm, ymin=0, ymax=100,
  symbolic x coords={C0,C1,C2}, xtick=data, enlarge x limits=0.15, ylabel={F1 bằng chứng (theo ý)},
  axis x line*=bottom, axis y line*=left, ymajorgrids, grid style={draw=black!12},
  tick label style={font=\small}, label style={font=\small}, clip=false,
  xticklabels={C0 (gốc),C1 (biến dạng thường),C2 (điểm mù)}]
\addplot[mauvhai, line width=2pt, mark=*, mark size=3pt] coordinates {(C0,83.0) (C1,76.5) (C2,75.3)};
\addplot[maullm, line width=2pt, mark=*, mark size=3pt] coordinates {(C0,70.7) (C1,56.1) (C2,32.2)};
\node[anchor=west,font=\small] at (axis cs:C2,75.3) {\ SafeViChi v2: 75,3};
\node[anchor=west,font=\small] at (axis cs:C2,32.2) {\ gpt-oss-120b: 32,2};
\node[anchor=south,font=\small] at (axis cs:C0,83.0) {83,0};
\node[anchor=north,font=\small] at (axis cs:C0,70.7) {70,7};
\end{axis}
\end{tikzpicture}

\caption{F1 bằng chứng (theo ý) khi tin nhắn bị biến dạng dần. gpt-oss-120b vẫn đoán được nhãn nhưng không còn chỉ
đúng chỗ; v2 cắt cụm ngay trên văn bản gốc nên giữ được.}
\label{fig:p2-bc-bien-dang}
\end{figure}

\begin{table}[H]
\centering
\small
\caption{Bằng chứng tô sáng của v2 và gpt-oss-120b. n BC: số câu nguy hại được chấm bằng chứng. Tô oan: câu
\nhan{AN_TOAN} bị tô ít nhất một cụm. Rỗng: câu nguy hại được báo đúng là nguy hại nhưng không có cụm nào. Sai chữ:
cụm không có nguyên văn trong tin nhắn.}
\label{tab:p2-bang-chung}
\footnotesize
\begin{tabular}{llccccccc}
\toprule
Tập & Hệ & n BC & BC P & BC R & BC F1 & Tô oan & Rỗng & Sai chữ \\
\midrule
\multirow{2}{*}{C0} & gpt-oss-120b & \so{C0.llm.nbc} & \so{C0.llm.bcp} & \so{C0.llm.bcr} & \so{C0.llm.bcf} & \so{C0.llm.toan} & \so{C0.llm.rong} & \so{C0.llm.khongnguyenvan} \\
 & v2 & \so{C0.v2.nbc} & \so{C0.v2.bcp} & \so{C0.v2.bcr} & \so{C0.v2.bcf} & \so{C0.v2.toan} & \so{C0.v2.rong} & \so{C0.v2.khongnguyenvan} \\
\midrule
\multirow{2}{*}{C1} & gpt-oss-120b & \so{C1.llm.nbc} & \so{C1.llm.bcp} & \so{C1.llm.bcr} & \so{C1.llm.bcf} & \so{C1.llm.toan} & \so{C1.llm.rong} & \so{C1.llm.khongnguyenvan} \\
 & v2 & \so{C1.v2.nbc} & \so{C1.v2.bcp} & \so{C1.v2.bcr} & \so{C1.v2.bcf} & \so{C1.v2.toan} & \so{C1.v2.rong} & \so{C1.v2.khongnguyenvan} \\
\midrule
\multirow{2}{*}{C2} & gpt-oss-120b & \so{C2.llm.nbc} & \so{C2.llm.bcp} & \so{C2.llm.bcr} & \so{C2.llm.bcf} & \so{C2.llm.toan} & \so{C2.llm.rong} & \so{C2.llm.khongnguyenvan} \\
 & v2 & \so{C2.v2.nbc} & \so{C2.v2.bcp} & \so{C2.v2.bcr} & \so{C2.v2.bcf} & \so{C2.v2.toan} & \so{C2.v2.rong} & \so{C2.v2.khongnguyenvan} \\
\bottomrule
\end{tabular}
\end{table}

\begin{itemize}
  \item Trên C0, bằng chứng của v2 vừa đúng hơn vừa đủ hơn gpt-oss-120b: P/R/F1 theo ý \so{C0.v2.bcp}/\so{C0.v2.bcr}/%
        \so{C0.v2.bcf} so với \so{C0.llm.bcp}/\so{C0.llm.bcr}/\so{C0.llm.bcf}.
  \item Khi chữ bị biến dạng, gpt-oss-120b gần như không trả được bằng chứng dùng được: ở C2, BC F1 chỉ còn
        \so{C2.llm.bcf} và \so{C2.llm.khongnguyenvan} số cụm không có nguyên văn trong tin. Lý do: LLM đọc hiểu được chữ
        lạ nhưng khi chép lại thì ``sửa'' chúng thành chữ thường, nên cụm không còn khớp tin gốc và không tô được.
        v2 cắt cụm từ chính văn bản nên vẫn giữ \so{C2.v2.bcf}.
  \item Điểm yếu của v2: \so{C0.v2.toan} câu an toàn bị tô oan (cao hơn \so{C0.llm.toan} của gpt-oss-120b). Con số này
        đi theo FPR: hễ câu an toàn bị báo nguy hại thì gần như luôn có cụm được tô.
\end{itemize}

\ghinho{\textbf{Cách chấm bằng chứng.} Điểm bằng chứng chỉ tính trên câu nguy hại mà hệ đã báo là nguy hại; câu nguy hại
bị đoán \nhan{AN_TOAN} đã bị phạt ở phần phân loại. Vì v2 có recall cao hơn, nó được chấm bằng chứng trên nhiều câu hơn
(C0: \so{C0.v2.nbc} so với \so{C0.llm.nbc} câu của gpt-oss-120b), tức là trên cả những câu khó mà LLM bỏ qua.}

% Phần 3 -- Yêu cầu 2: khái quát trước kiểu né lọc mới. Mọi số lấy qua \so{} (so_lieu.tex).
\section{Yêu cầu 2: khái quát trước kiểu né lọc mới}
\label{sec:yc2}

Nhóm đánh giá khả năng khái quát theo ba trục trên tập kiểm thử (Mục~\ref{sec:p1-test}): kiểu biến dạng chưa thấy
(C1, C2), kịch bản chưa thấy (\so{te.chuathay} đoạn C0) và giọng vùng miền. Kiểm tra rò rỉ giữa hai tập ghi nhận \so{rr.trunggan} cặp gần trùng (tương đồng 3-gram $\ge \so{rr.nguong}$) và \so{rr.trainkichban} đoạn
huấn luyện thuộc các kịch bản giữ ngoài.

C1 tác động lên chính tả, âm tiết nên có thể làm mất thông tin (``đánh'' và ``dành'' bỏ dấu đều thành ``danh''); C2
giữ nguyên âm tiết nhưng đổi lớp ký tự Unicode. Nhờ nhánh CNN đọc theo byte, v2 đạt Binary-F1 \so{C2.v2.bin} trên C2,
nhỉnh hơn C1 (\so{C1.v2.bin}); Macro-F1 trên hai tập gần như nhau (\so{C1.v2.macro} và \so{C2.v2.macro}).

\subsection{Theo từng kiểu biến dạng}

\begin{figure}[H]
\centering
% Sinh tự động bằng `python -m src.v2.eval.xuat_so_lieu_bao_cao`. Không sửa tay.
\begin{tikzpicture}
\draw[black!12] (0.00,0.3) -- (0.00,-7.80);
\node[font=\scriptsize,anchor=north] at (0.00,-7.80) {40};
\draw[black!12] (1.64,0.3) -- (1.64,-7.80);
\node[font=\scriptsize,anchor=north] at (1.64,-7.80) {50};
\draw[black!12] (3.27,0.3) -- (3.27,-7.80);
\node[font=\scriptsize,anchor=north] at (3.27,-7.80) {60};
\draw[black!12] (4.91,0.3) -- (4.91,-7.80);
\node[font=\scriptsize,anchor=north] at (4.91,-7.80) {70};
\draw[black!12] (6.55,0.3) -- (6.55,-7.80);
\node[font=\scriptsize,anchor=north] at (6.55,-7.80) {80};
\draw[black!12] (8.18,0.3) -- (8.18,-7.80);
\node[font=\scriptsize,anchor=north] at (8.18,-7.80) {90};
\node[font=\small,anchor=north] at (4.50,-8.25) {Macro-F1 5 nhãn};
\node[anchor=west,font=\small\bfseries] at (-5.2,0.00) {C1: hỏng chính con chữ};
\node[anchor=east,font=\footnotesize] at (-0.15,-0.55) {Dính từ};
\draw[black!30, line width=1.2pt] (2.84,-0.55) -- (7.71,-0.55);
\fill[maullm] (2.84,-0.55) circle (3.2pt);
\fill[mauvhai] (7.71,-0.55) circle (3.2pt);
\node[anchor=west,font=\scriptsize] at (7.83,-0.55) {87,1};
\node[anchor=east,font=\scriptsize,text=black!70] at (2.72,-0.55) {57,4};
\node[anchor=east,font=\footnotesize] at (-0.15,-1.05) {Kéo dài ký tự};
\draw[black!30, line width=1.2pt] (2.48,-1.05) -- (6.71,-1.05);
\fill[maullm] (2.48,-1.05) circle (3.2pt);
\fill[mauvhai] (6.71,-1.05) circle (3.2pt);
\node[anchor=west,font=\scriptsize] at (6.83,-1.05) {81,0};
\node[anchor=east,font=\scriptsize,text=black!70] at (2.36,-1.05) {55,1};
\node[anchor=east,font=\footnotesize] at (-0.15,-1.55) {Bỏ dấu};
\draw[black!30, line width=1.2pt] (2.96,-1.55) -- (6.56,-1.55);
\fill[maullm] (2.96,-1.55) circle (3.2pt);
\fill[mauvhai] (6.56,-1.55) circle (3.2pt);
\node[anchor=west,font=\scriptsize] at (6.68,-1.55) {80,1};
\node[anchor=east,font=\scriptsize,text=black!70] at (2.84,-1.55) {58,1};
\node[anchor=east,font=\footnotesize] at (-0.15,-2.05) {Số thay chữ};
\draw[black!30, line width=1.2pt] (4.50,-2.05) -- (6.47,-2.05);
\fill[maullm] (4.50,-2.05) circle (3.2pt);
\fill[mauvhai] (6.47,-2.05) circle (3.2pt);
\node[anchor=west,font=\scriptsize] at (6.59,-2.05) {79,5};
\node[anchor=east,font=\scriptsize,text=black!70] at (4.38,-2.05) {67,5};
\node[anchor=east,font=\footnotesize] at (-0.15,-2.55) {Chèn ký tự};
\draw[black!30, line width=1.2pt] (1.70,-2.55) -- (6.01,-2.55);
\fill[maullm] (1.70,-2.55) circle (3.2pt);
\fill[mauvhai] (6.01,-2.55) circle (3.2pt);
\node[anchor=west,font=\scriptsize] at (6.13,-2.55) {76,7};
\node[anchor=east,font=\scriptsize,text=black!70] at (1.58,-2.55) {50,4};
\node[anchor=east,font=\footnotesize] at (-0.15,-3.05) {Gõ sai};
\draw[black!30, line width=1.2pt] (2.19,-3.05) -- (5.36,-3.05);
\fill[maullm] (2.19,-3.05) circle (3.2pt);
\fill[mauvhai] (5.36,-3.05) circle (3.2pt);
\node[anchor=west,font=\scriptsize] at (5.48,-3.05) {72,8};
\node[anchor=east,font=\scriptsize,text=black!70] at (2.07,-3.05) {53,4};
\node[anchor=east,font=\footnotesize] at (-0.15,-3.55) {Teencode};
\draw[black!30, line width=1.2pt] (2.62,-3.55) -- (5.06,-3.55);
\fill[maullm] (2.62,-3.55) circle (3.2pt);
\fill[mauvhai] (5.06,-3.55) circle (3.2pt);
\node[anchor=west,font=\scriptsize] at (5.18,-3.55) {70,9};
\node[anchor=east,font=\scriptsize,text=black!70] at (2.50,-3.55) {56,0};
\node[anchor=west,font=\small\bfseries] at (-5.2,-4.25) {C2: đổi lớp ký tự (chưa thấy khi train)};
\node[anchor=east,font=\footnotesize] at (-0.15,-4.80) {Chữ toàn khổ};
\draw[black!30, line width=1.2pt] (5.52,-4.80) -- (7.02,-4.80);
\fill[maullm] (5.52,-4.80) circle (3.2pt);
\fill[mauvhai] (7.02,-4.80) circle (3.2pt);
\node[anchor=west,font=\scriptsize] at (7.14,-4.80) {82,9};
\node[anchor=east,font=\scriptsize,text=black!70] at (5.40,-4.80) {73,7};
\node[anchor=east,font=\footnotesize] at (-0.15,-5.30) {Kéo dài kép};
\draw[black!30, line width=1.2pt] (4.42,-5.30) -- (6.94,-5.30);
\fill[maullm] (4.42,-5.30) circle (3.2pt);
\fill[mauvhai] (6.94,-5.30) circle (3.2pt);
\node[anchor=west,font=\scriptsize] at (7.06,-5.30) {82,4};
\node[anchor=east,font=\scriptsize,text=black!70] at (4.30,-5.30) {67,0};
\node[anchor=east,font=\footnotesize] at (-0.15,-5.80) {Dấu phân cách lạ};
\draw[black!30, line width=1.2pt] (2.59,-5.80) -- (6.24,-5.80);
\fill[maullm] (2.59,-5.80) circle (3.2pt);
\fill[mauvhai] (6.24,-5.80) circle (3.2pt);
\node[anchor=west,font=\scriptsize] at (6.36,-5.80) {78,2};
\node[anchor=east,font=\scriptsize,text=black!70] at (2.47,-5.80) {55,8};
\node[anchor=east,font=\footnotesize] at (-0.15,-6.30) {Chữ Cyrillic giống Latin};
\draw[black!30, line width=1.2pt] (5.30,-6.30) -- (6.23,-6.30);
\fill[maullm] (5.30,-6.30) circle (3.2pt);
\fill[mauvhai] (6.23,-6.30) circle (3.2pt);
\node[anchor=west,font=\scriptsize] at (6.35,-6.30) {78,1};
\node[anchor=east,font=\scriptsize,text=black!70] at (5.18,-6.30) {72,4};
\node[anchor=east,font=\footnotesize] at (-0.15,-6.80) {Ký tự vô hình};
\draw[black!30, line width=1.2pt] (2.92,-6.80) -- (5.30,-6.80);
\fill[maullm] (2.92,-6.80) circle (3.2pt);
\fill[mauvhai] (5.30,-6.80) circle (3.2pt);
\node[anchor=west,font=\scriptsize] at (5.42,-6.80) {72,4};
\node[anchor=east,font=\scriptsize,text=black!70] at (2.80,-6.80) {57,9};
\node[anchor=east,font=\footnotesize] at (-0.15,-7.30) {Chèn dấu cách};
\draw[black!30, line width=1.2pt] (0.83,-7.30) -- (4.43,-7.30);
\fill[maullm] (0.83,-7.30) circle (3.2pt);
\fill[mauvhai] (4.43,-7.30) circle (3.2pt);
\node[anchor=west,font=\scriptsize] at (4.55,-7.30) {67,1};
\node[anchor=east,font=\scriptsize,text=black!70] at (0.71,-7.30) {45,1};
\fill[mauvhai] (1.0,0.75) circle (3.2pt); \node[anchor=west,font=\small] at (1.15,0.75) {v2};\fill[maullm] (2.6,0.75) circle (3.2pt); \node[anchor=west,font=\small] at (2.75,0.75) {gpt-oss-120b (API)};
\end{tikzpicture}

\caption{Macro-F1 5 nhãn theo 13 kiểu biến dạng, sắp giảm dần theo điểm của SafeViChi v2: gpt-oss-120b (API) (cam) và
SafeViChi v2 (xanh).}
\label{fig:p3-kieu}
\end{figure}

\begin{table}[H]
\centering
\footnotesize
\setlength{\tabcolsep}{3pt}
\caption{Kết quả theo từng kiểu biến dạng trên tập kiểm thử (Macro-F1, Bin-F1, BC F1 theo ý trên thang 0--100; FPR tính
theo \%; $n$: số đoạn; v1 là SafeViChi v1). Cột \textbf{Train C0}: kiểu có xuất hiện tự nhiên trong tập huấn luyện gốc
hay không.}
\label{tab:p3-kieu}
\begin{tabularx}{\linewidth}{@{}Y c c c c c c c c c@{}}
\toprule
 & & & \textbf{v1} & \multicolumn{2}{c}{\textbf{gpt-oss-120b (API)}} & \multicolumn{4}{c}{\textbf{SafeViChi v2}} \\
\cmidrule(lr){4-4}\cmidrule(lr){5-6}\cmidrule(l){7-10}
\textbf{Kiểu biến dạng} & \textbf{Train C0} & \textbf{$n$} & \textbf{Bin-F1} & \textbf{Macro} & \textbf{Bin-F1} & \textbf{Macro} & \textbf{Bin-F1} & \textbf{FPR} & \textbf{BC} \\
\midrule
\multicolumn{10}{@{}l}{\textit{C1: biến dạng bề mặt}} \\
Bỏ dấu & có & \so{kieu.C1.bodau.n} & \so{kieu.C1.bodau.v1.bin} & \so{kieu.C1.bodau.llm.macro} & \so{kieu.C1.bodau.llm.bin} & \textbf{\so{kieu.C1.bodau.v2.macro}} & \textbf{\so{kieu.C1.bodau.v2.bin}} & \so{kieu.C1.bodau.v2.fpr} & \so{kieu.C1.bodau.v2.bcf} \\
Teencode & có & \so{kieu.C1.teencode.n} & \so{kieu.C1.teencode.v1.bin} & \so{kieu.C1.teencode.llm.macro} & \so{kieu.C1.teencode.llm.bin} & \textbf{\so{kieu.C1.teencode.v2.macro}} & \textbf{\so{kieu.C1.teencode.v2.bin}} & \so{kieu.C1.teencode.v2.fpr} & \so{kieu.C1.teencode.v2.bcf} \\
Gõ sai & có & \so{kieu.C1.gosai.n} & \so{kieu.C1.gosai.v1.bin} & \so{kieu.C1.gosai.llm.macro} & \so{kieu.C1.gosai.llm.bin} & \textbf{\so{kieu.C1.gosai.v2.macro}} & \textbf{\so{kieu.C1.gosai.v2.bin}} & \so{kieu.C1.gosai.v2.fpr} & \so{kieu.C1.gosai.v2.bcf} \\
Kéo dài ký tự & có & \so{kieu.C1.keodai.n} & \so{kieu.C1.keodai.v1.bin} & \so{kieu.C1.keodai.llm.macro} & \so{kieu.C1.keodai.llm.bin} & \textbf{\so{kieu.C1.keodai.v2.macro}} & \textbf{\so{kieu.C1.keodai.v2.bin}} & \so{kieu.C1.keodai.v2.fpr} & \so{kieu.C1.keodai.v2.bcf} \\
Chèn ký tự & không & \so{kieu.C1.chenkytu.n} & \so{kieu.C1.chenkytu.v1.bin} & \so{kieu.C1.chenkytu.llm.macro} & \so{kieu.C1.chenkytu.llm.bin} & \textbf{\so{kieu.C1.chenkytu.v2.macro}} & \textbf{\so{kieu.C1.chenkytu.v2.bin}} & \so{kieu.C1.chenkytu.v2.fpr} & \so{kieu.C1.chenkytu.v2.bcf} \\
Số thay chữ & không & \so{kieu.C1.sothaychu.n} & \so{kieu.C1.sothaychu.v1.bin} & \so{kieu.C1.sothaychu.llm.macro} & \so{kieu.C1.sothaychu.llm.bin} & \textbf{\so{kieu.C1.sothaychu.v2.macro}} & \textbf{\so{kieu.C1.sothaychu.v2.bin}} & \so{kieu.C1.sothaychu.v2.fpr} & \so{kieu.C1.sothaychu.v2.bcf} \\
Dính từ & không & \so{kieu.C1.dinhtu.n} & \so{kieu.C1.dinhtu.v1.bin} & \so{kieu.C1.dinhtu.llm.macro} & \so{kieu.C1.dinhtu.llm.bin} & \textbf{\so{kieu.C1.dinhtu.v2.macro}} & \textbf{\so{kieu.C1.dinhtu.v2.bin}} & \so{kieu.C1.dinhtu.v2.fpr} & \so{kieu.C1.dinhtu.v2.bcf} \\
\midrule
\multicolumn{10}{@{}l}{\textit{C2: đổi lớp ký tự Unicode (điểm mù, hoàn toàn ngoài huấn luyện)}} \\
Chữ toàn khổ {\scriptsize(fullwidth)} & không & \so{kieu.C2.fullwidth.n} & \so{kieu.C2.fullwidth.v1.bin} & \so{kieu.C2.fullwidth.llm.macro} & \so{kieu.C2.fullwidth.llm.bin} & \textbf{\so{kieu.C2.fullwidth.v2.macro}} & \textbf{\so{kieu.C2.fullwidth.v2.bin}} & \so{kieu.C2.fullwidth.v2.fpr} & \so{kieu.C2.fullwidth.v2.bcf} \\
Ký tự vô hình {\scriptsize(zero-width)} & không & \so{kieu.C2.zerowidth.n} & \so{kieu.C2.zerowidth.v1.bin} & \so{kieu.C2.zerowidth.llm.macro} & \so{kieu.C2.zerowidth.llm.bin} & \textbf{\so{kieu.C2.zerowidth.v2.macro}} & \textbf{\so{kieu.C2.zerowidth.v2.bin}} & \so{kieu.C2.zerowidth.v2.fpr} & \so{kieu.C2.zerowidth.v2.bcf} \\
Chèn dấu cách {\scriptsize(space injection)} & không & \so{kieu.C2.spaceinjection.n} & \so{kieu.C2.spaceinjection.v1.bin} & \so{kieu.C2.spaceinjection.llm.macro} & \so{kieu.C2.spaceinjection.llm.bin} & \textbf{\so{kieu.C2.spaceinjection.v2.macro}} & \textbf{\so{kieu.C2.spaceinjection.v2.bin}} & \so{kieu.C2.spaceinjection.v2.fpr} & \so{kieu.C2.spaceinjection.v2.bcf} \\
Dấu phân cách lạ {\scriptsize(exotic separator)} & không & \so{kieu.C2.exoticseparator.n} & \so{kieu.C2.exoticseparator.v1.bin} & \so{kieu.C2.exoticseparator.llm.macro} & \so{kieu.C2.exoticseparator.llm.bin} & \textbf{\so{kieu.C2.exoticseparator.v2.macro}} & \textbf{\so{kieu.C2.exoticseparator.v2.bin}} & \so{kieu.C2.exoticseparator.v2.fpr} & \so{kieu.C2.exoticseparator.v2.bcf} \\
Chữ Cyrillic giống Latin {\scriptsize(homoglyph)} & không & \so{kieu.C2.cyrillichomoglyph.n} & \so{kieu.C2.cyrillichomoglyph.v1.bin} & \so{kieu.C2.cyrillichomoglyph.llm.macro} & \so{kieu.C2.cyrillichomoglyph.llm.bin} & \textbf{\so{kieu.C2.cyrillichomoglyph.v2.macro}} & \textbf{\so{kieu.C2.cyrillichomoglyph.v2.bin}} & \so{kieu.C2.cyrillichomoglyph.v2.fpr} & \so{kieu.C2.cyrillichomoglyph.v2.bcf} \\
Kéo dài kép {\scriptsize(double elongation)} & không & \so{kieu.C2.doubleelongation.n} & \so{kieu.C2.doubleelongation.v1.bin} & \so{kieu.C2.doubleelongation.llm.macro} & \so{kieu.C2.doubleelongation.llm.bin} & \textbf{\so{kieu.C2.doubleelongation.v2.macro}} & \textbf{\so{kieu.C2.doubleelongation.v2.bin}} & \so{kieu.C2.doubleelongation.v2.fpr} & \so{kieu.C2.doubleelongation.v2.bcf} \\
\bottomrule
\end{tabularx}
\end{table}

\textbf{Nhận xét:}
\begin{itemize}
  \item \textbf{SafeViChi v2 đứng đầu ở cả 13 kiểu biến dạng} về Macro-F1 và Binary-F1, kể cả kiểu chưa có
        khi huấn luyện như \emph{dính từ} (\so{kieu.C1.dinhtu.v2.macro} so với \so{kieu.C1.dinhtu.llm.macro} của LLM).
        Nhóm C1 giữ ngoài tập huấn luyện gốc còn đạt Macro-F1 cao hơn nhóm đã thấy (\so{nhom.C1.chuathay.v2.macro} so
        với \so{nhom.C1.dathay.v2.macro}).
  \item \textbf{Bằng chứng trên biến dạng điểm mù:} với \emph{chữ toàn khổ} và \emph{ký tự vô hình}, BC F1 của
        gpt-oss-120b (API) chỉ còn \so{kieu.C2.fullwidth.llm.bcf} và \so{kieu.C2.zerowidth.llm.bcf}, so với
        \so{kieu.C2.fullwidth.v2.bcf} và \so{kieu.C2.zerowidth.v2.bcf} của v2.
  \item \textbf{FPR tăng ở các kiểu phá vỡ ranh giới từ.} Phần lớn các kiểu giữ FPR thấp (\emph{bỏ dấu}
        \so{kieu.C1.bodau.v2.fpr}, \emph{chữ toàn khổ} \so{kieu.C2.fullwidth.v2.fpr}); FPR tăng chủ yếu ở \emph{chèn dấu cách} (\so{kieu.C2.spaceinjection.v2.fpr}), \emph{gõ
        sai} (\so{kieu.C1.gosai.v2.fpr}) và \emph{ký tự vô hình} (\so{kieu.C2.zerowidth.v2.fpr}). Do mỗi kiểu chỉ có
        khoảng \so{kieu.C2.spaceinjection.nat} đoạn \nhan{AN_TOAN}, một đoạn sai đã làm FPR đổi hơn 3 điểm phần trăm.
\end{itemize}

\subsection{Khái quát trên kịch bản chưa thấy}

\begin{figure}[H]
\centering
% Sinh tự động bằng `python -m src.v2.eval.xuat_so_lieu_bao_cao`. Không sửa tay.
\begin{tikzpicture}
\begin{axis}[ybar=2pt, bar width=13pt, width=\linewidth, height=6cm, ymin=0, ymax=100,
  symbolic x coords={C0,C1,C2}, xtick=data, enlarge x limits=0.22, ylabel={Macro-F1 5 nhãn},
  axis x line*=bottom, axis y line*=left, ymajorgrids, grid style={draw=black!12},
  tick label style={font=\small}, label style={font=\small},
  nodes near coords={\pgfmathprintnumber[fixed,fixed zerofill,precision=1,use comma]{\pgfplotspointmeta}},
  every node near coord/.append style={font=\tiny, color=black!75}, legend cell align=left,
  legend style={at={(0.5,1.03)}, anchor=south, legend columns=2, draw=none, font=\footnotesize,
    /tikz/every even column/.append style={column sep=0.5cm}}]
\addplot[fill=mauvhai, draw=white, area legend] coordinates {(C0,93.0) (C1,86.7) (C2,86.0)};
\addlegendentry{v2, kịch bản đã có trong train}
\addplot[fill=mauvhai!35, draw=mauvhai, area legend] coordinates {(C0,62.1) (C1,54.4) (C2,53.2)};
\addlegendentry{v2, kịch bản chưa thấy}
\addplot[fill=maullm, draw=white, area legend] coordinates {(C0,84.8) (C1,63.3) (C2,67.6)};
\addlegendentry{gpt-oss-120b, kịch bản đã có trong train v2}
\addplot[fill=maullm!35, draw=maullm, area legend] coordinates {(C0,53.0) (C1,39.0) (C2,45.5)};
\addlegendentry{gpt-oss-120b, kịch bản chưa thấy}
\end{axis}
\end{tikzpicture}

\caption{Macro-F1 5 nhãn trên kịch bản đã có trong tập huấn luyện (màu đậm) và kịch bản chưa thấy (màu nhạt).}
\label{fig:p3-kichban}
\end{figure}

\begin{table}[H]
\centering
\small
\setlength{\tabcolsep}{4.5pt}
\caption{Kết quả trên kịch bản trong phân phối huấn luyện và kịch bản chưa thấy (Macro-F1, Bin-F1, BC F1 theo ý
trên thang 0--100; FPR tính theo \%; $n$: số đoạn).}
\label{tab:p3-kichban}
\begin{tabular}{@{}l l c c c c c c@{}}
\toprule
 & & & \multicolumn{4}{c}{\textbf{SafeViChi v2}} & \textbf{gpt-oss-120b (API)} \\
\cmidrule(lr){4-7}\cmidrule(l){8-8}
\textbf{Tập} & \textbf{Phân nhóm} & \textbf{$n$} & \textbf{Macro-F1} & \textbf{Bin-F1} & \textbf{FPR} & \textbf{BC F1} & \textbf{Macro-F1} \\
\midrule
\multirow{2}{*}{C0} & Trong phân phối & \so{phan.C0.trongphanphoi.n} & \textbf{\so{phan.C0.trongphanphoi.v2.macro}} & \so{phan.C0.trongphanphoi.v2.bin} & \so{phan.C0.trongphanphoi.v2.fpr} & \so{phan.C0.trongphanphoi.v2.bcf} & \so{phan.C0.trongphanphoi.llm.macro} \\
 & Chưa thấy & \so{phan.C0.chuathay.n} & \textbf{\so{phan.C0.chuathay.v2.macro}} & \so{phan.C0.chuathay.v2.bin} & \so{phan.C0.chuathay.v2.fpr} & \so{phan.C0.chuathay.v2.bcf} & \so{phan.C0.chuathay.llm.macro} \\
\midrule
\multirow{2}{*}{C1} & Trong phân phối & \so{phan.C1.trongphanphoi.n} & \textbf{\so{phan.C1.trongphanphoi.v2.macro}} & \so{phan.C1.trongphanphoi.v2.bin} & \so{phan.C1.trongphanphoi.v2.fpr} & \so{phan.C1.trongphanphoi.v2.bcf} & \so{phan.C1.trongphanphoi.llm.macro} \\
 & Chưa thấy & \so{phan.C1.chuathay.n} & \textbf{\so{phan.C1.chuathay.v2.macro}} & \so{phan.C1.chuathay.v2.bin} & \so{phan.C1.chuathay.v2.fpr} & \so{phan.C1.chuathay.v2.bcf} & \so{phan.C1.chuathay.llm.macro} \\
\midrule
\multirow{2}{*}{C2} & Trong phân phối & \so{phan.C2.trongphanphoi.n} & \textbf{\so{phan.C2.trongphanphoi.v2.macro}} & \so{phan.C2.trongphanphoi.v2.bin} & \so{phan.C2.trongphanphoi.v2.fpr} & \so{phan.C2.trongphanphoi.v2.bcf} & \so{phan.C2.trongphanphoi.llm.macro} \\
 & Chưa thấy & \so{phan.C2.chuathay.n} & \textbf{\so{phan.C2.chuathay.v2.macro}} & \so{phan.C2.chuathay.v2.bin} & \so{phan.C2.chuathay.v2.fpr} & \so{phan.C2.chuathay.v2.bcf} & \so{phan.C2.chuathay.llm.macro} \\
\bottomrule
\end{tabular}
\end{table}

Trên C0, Macro-F1 của v2 giảm từ \so{phan.C0.trongphanphoi.v2.macro} (kịch bản đã thấy) xuống
\so{phan.C0.chuathay.v2.macro} (chưa thấy) nhưng vẫn cao hơn gpt-oss-120b (API) ở nhóm chưa thấy trên cả ba tập
(Bảng~\ref{tab:p3-kichban}). gpt-oss-120b (API) cũng giảm mạnh ở nhóm này (\so{phan.C0.trongphanphoi.llm.macro} $\to$
\so{phan.C0.chuathay.llm.macro}), cho thấy tình huống ẩn ý như đe dọa trả thù (DD6) hay kiểm soát người yêu (QR6)
vốn khó phân biệt giữa \nhan{DE_DOA}, \nhan{QUAY_ROI} và \nhan{AN_TOAN}. Nhóm chưa thấy ở C0 chỉ có \so{te.chuathay.xp}
đoạn \nhan{XUC_PHAM} nên Macro-F1 dao động mạnh theo từng mẫu.

\subsection{Theo giọng vùng miền}

\begin{table}[H]
\centering
\footnotesize
\setlength{\tabcolsep}{3.5pt}
\caption{Kết quả theo giọng vùng miền trên C0 (Macro-F1, Bin-F1 trên thang 0--100; FPR tính theo \%), kèm FPR của SafeViChi v2
trên C1 và C2. Chênh lệch là hiệu giữa vùng cao nhất và vùng thấp nhất.}
\label{tab:p3-vung}
\begin{tabular}{@{}l c c c c c c c c c@{}}
\toprule
 & & \multicolumn{3}{c}{\textbf{SafeViChi v2 (C0)}} & \multicolumn{3}{c}{\textbf{gpt-oss-120b (API, C0)}} & \multicolumn{2}{c}{\textbf{FPR v2}} \\
\cmidrule(lr){3-5}\cmidrule(lr){6-8}\cmidrule(l){9-10}
\textbf{Giọng} & \textbf{$n$} & \textbf{Macro} & \textbf{Bin-F1} & \textbf{FPR} & \textbf{Macro} & \textbf{Bin-F1} & \textbf{FPR} & \textbf{C1} & \textbf{C2} \\
\midrule
Bắc & \so{C0.bac.v2.n} & \textbf{\so{C0.bac.v2.macro}} & \textbf{\so{C0.bac.v2.bin}} & \so{C0.bac.v2.fpr} & \so{C0.bac.llm.macro} & \so{C0.bac.llm.bin} & \so{C0.bac.llm.fpr} & \so{C1.bac.v2.fpr} & \so{C2.bac.v2.fpr} \\
Trung & \so{C0.trung.v2.n} & \textbf{\so{C0.trung.v2.macro}} & \textbf{\so{C0.trung.v2.bin}} & \so{C0.trung.v2.fpr} & \so{C0.trung.llm.macro} & \so{C0.trung.llm.bin} & \so{C0.trung.llm.fpr} & \so{C1.trung.v2.fpr} & \so{C2.trung.v2.fpr} \\
Nam & \so{C0.nam.v2.n} & \textbf{\so{C0.nam.v2.macro}} & \textbf{\so{C0.nam.v2.bin}} & \so{C0.nam.v2.fpr} & \so{C0.nam.llm.macro} & \so{C0.nam.llm.bin} & \so{C0.nam.llm.fpr} & \so{C1.nam.v2.fpr} & \so{C2.nam.v2.fpr} \\
\midrule
Chênh lệch & & \textbf{\so{C0.v2.gap.macro}} & \textbf{\so{C0.v2.gap.bin}} & \so{C0.v2.gap.fpr} & \so{C0.llm.gap.macro} & \so{C0.llm.gap.bin} & \so{C0.llm.gap.fpr} & \so{C1.v2.gap.fpr} & \so{C2.v2.gap.fpr} \\
\bottomrule
\end{tabular}
\end{table}

SafeViChi v2 có Macro-F1 và Binary-F1 cao hơn gpt-oss-120b (API) ở cả ba giọng (Bảng~\ref{tab:p3-vung}). Chênh lệch
Macro-F1 giữa các vùng của v2 là \so{C0.v2.gap.macro} điểm, nhỏ hơn \so{C0.llm.gap.macro} của gpt-oss-120b (API).
Ngược lại, FPR của v2 chênh lớn giữa các vùng (Trung \so{C0.trung.v2.fpr}, Bắc \so{C0.bac.v2.fpr}, Nam
\so{C0.nam.v2.fpr}; chênh \so{C0.v2.gap.fpr}) và lặp lại trên C1, C2, trong khi FPR của gpt-oss-120b (API) gần như
bằng nhau (chênh \so{C0.llm.gap.fpr}). Tuy vậy, mỗi vùng chỉ có
\so{C0.nam.v2.at.n}--\so{C0.bac.v2.at.n} đoạn \nhan{AN_TOAN} nên mỗi đoạn sai làm FPR đổi khoảng 3 điểm phần trăm;
cần thêm mẫu an toàn theo vùng trước khi kết luận v2 thiên lệch theo giọng.

\subsection{Các lát cắt khác}

\begin{table}[H]
\centering
\small
\caption{Kết quả trên một số lát cắt ngữ cảnh, độ lộ và kênh giao tiếp ở tập kiểm thử C0 (thang 0--100; $n$: số đoạn).}
\label{tab:p3-latcat}
\begin{tabular}{@{}l c c c c@{}}
\toprule
 & & \multicolumn{2}{c}{\textbf{SafeViChi v2}} & \textbf{gpt-oss-120b (API)} \\
\cmidrule(lr){3-4}\cmidrule(l){5-5}
\textbf{Lát cắt} & \textbf{$n$} & \textbf{Macro-F1} & \textbf{Bin-F1} & \textbf{Macro-F1} \\
\midrule
Hội thoại 1 dòng & \so{luot.C0.1dong.n} & \textbf{\so{luot.C0.1dong.v2.macro}} & \so{luot.C0.1dong.v2.bin} & \so{luot.C0.1dong.llm.macro} \\
Hội thoại 2--8 dòng & \so{luot.C0.nhieudong.n} & \textbf{\so{luot.C0.nhieudong.v2.macro}} & \so{luot.C0.nhieudong.v2.bin} & \so{luot.C0.nhieudong.llm.macro} \\
Đoạn nguy hại lộ liễu & \so{dolo.C0.lolieu.n} & \textbf{\so{dolo.C0.lolieu.v2.macro}} & \so{dolo.C0.lolieu.v2.bin} & \so{dolo.C0.lolieu.llm.macro} \\
Đoạn nguy hại ẩn ý & \so{dolo.C0.any.n} & \textbf{\so{dolo.C0.any.v2.macro}} & \so{dolo.C0.any.v2.bin} & \so{dolo.C0.any.llm.macro} \\
Kênh Messenger/Zalo & \so{kenh.C0.messengerzalo.n} & \textbf{\so{kenh.C0.messengerzalo.v2.macro}} & \so{kenh.C0.messengerzalo.v2.bin} & \so{kenh.C0.messengerzalo.llm.macro} \\
Kênh nhóm lớp & \so{kenh.C0.nhomlop.n} & \textbf{\so{kenh.C0.nhomlop.v2.macro}} & \so{kenh.C0.nhomlop.v2.bin} & \so{kenh.C0.nhomlop.llm.macro} \\
Kênh SMS & \so{kenh.C0.sms.n} & \so{kenh.C0.sms.v2.macro} & \so{kenh.C0.sms.v2.bin} & \textbf{\so{kenh.C0.sms.llm.macro}} \\
\bottomrule
\end{tabular}
\end{table}

Theo Bảng~\ref{tab:p3-latcat}, khoảng cách giữa v2 và gpt-oss-120b (API) lớn nhất ở hội thoại 2--8 dòng và đoạn nguy hại
ẩn ý. Kênh SMS (\so{kenh.C0.sms.n} đoạn, nhiều tin một dòng thiếu ngữ cảnh như thông báo OTP, giao hàng hợp lệ) là lát
cắt duy nhất gpt-oss-120b (API) cao hơn.

\subsection{Phân tích một số trường hợp thất bại}

\begin{table}[H]
\centering
\footnotesize
\caption{Bốn đoạn nguy hại thuộc kịch bản chưa thấy bị SafeViChi v2 dự đoán là \nhan{AN_TOAN} trên C0 (mỗi dòng trong ô là một
lượt nói).}
\label{tab:p3-vidu}
\begin{tabularx}{\linewidth}{@{}>{\raggedright\arraybackslash}p{0.41\linewidth} >{\raggedright\arraybackslash}p{0.20\linewidth} Y@{}}
\toprule
\textbf{Đoạn tin nhắn} & \textbf{Kết quả} & \textbf{Nguyên nhân khó} \\
\midrule
\so{vd.ghiam.text} &
Nhãn chuẩn: \so{vd.ghiam.gold}\newline v2: \so{vd.ghiam.pred}\newline LLM (API): \so{vd.ghiam.llm} &
Kịch bản \so{vd.ghiam.kichban} (đe dọa trả thù vì bị báo thầy cô). Đe dọa ẩn trong câu ``nhớ đi đường đông người nhé'', không có từ thô tục; gpt-oss-120b (API) cũng bỏ sót. \\
\midrule
\so{vd.vay.text} &
Nhãn chuẩn: \so{vd.vay.gold}\newline v2: \so{vd.vay.pred}\newline LLM (API): \so{vd.vay.llm} &
Kịch bản \so{vd.vay.kichban} (vay trực tuyến chỉ cần căn cước) giữ ngoài tập huấn luyện. Đòi giấy tờ thay vì chuyển tiền ngay nên khác mẫu lừa đảo quen thuộc. \\
\midrule
\so{vd.nham.text} &
Nhãn chuẩn: \so{vd.nham.gold}\newline v2: \so{vd.nham.pred}\newline LLM (API): \so{vd.nham.llm} &
Chiêu ``chuyển nhầm tiền'' hoàn toàn mới, lời lẽ lịch sự như nhờ giúp đỡ thông thường. \\
\midrule
\so{vd.kiemsoat.text} &
Nhãn chuẩn: \so{vd.kiemsoat.gold}\newline v2: \so{vd.kiemsoat.pred}\newline LLM (API): \so{vd.kiemsoat.llm} &
Kịch bản \so{vd.kiemsoat.kichban} (kiểm soát người yêu) giữ ngoài tập huấn luyện. Không có từ xúc phạm hay đe dọa thân thể; nguy hại nằm ở việc ép kiểm soát tài khoản và đòi gửi ảnh chụp màn hình. \\
\bottomrule
\end{tabularx}
\end{table}

Điểm chung của bốn trường hợp: mô hình học tốt mẫu đã thấy nhưng chưa nắm khái niệm tổng quát về ``ép buộc''
hay ``dụ chuyển tiền'' khi cách diễn đạt hoàn toàn mới.

% Phần 4: Yêu cầu 3 (so sánh phương pháp, chi phí) và hạn chế.
% Mọi con số lấy qua \so{...} từ so_lieu.tex.

\section{Yêu cầu 3: so sánh phương pháp và chi phí}
\label{sec:yc3}

\subsection{Chọn mô hình: bốn ứng viên trên tập dev}
\label{sec:yc3-chon}

Nhóm huấn luyện song song bốn ứng viên trên cùng tập dữ liệu, cùng phân hoạch phân tầng theo nhóm trùng lặp
(\so{dev.chia.fit} mẫu \textbf{fit}, \so{dev.chia.dev} mẫu \textbf{dev}, \so{dev.chia.cal} mẫu \textbf{cal}), cùng
seed 42 và cùng tiêu chí chọn checkpoint. Ba ứng viên sử dụng chung kiến trúc bộ mã hóa \textbf{SafeVi-DualORC-CRF}
(phân loại 5 nhãn kết hợp gắn thẻ BIO + CRF mức từ), chỉ khác bộ mã hóa nền: mmBERT-base (đa ngôn ngữ), BamiBERT và
ViSoBERT (tiếng Việt). Ứng viên thứ tư (\textbf{GuardMerge}) chuyển đổi từ nền mô hình ngôn ngữ lớn hai chiều khoảng
0,6 tỉ tham số và tinh chỉnh bằng LoRA. Toàn bộ quy trình chọn mô hình được thực hiện độc lập trên tập \textbf{dev}
trước khi mở khóa tập kiểm thử (Hình~\ref{fig:p4-ung-vien}, Bảng~\ref{tab:p4-ung-vien}).

\begin{figure}[H]
\centering
% Sinh tự động bằng `python -m src.v2.eval.xuat_so_lieu_bao_cao`. Không sửa tay.
\begin{tikzpicture}
\begin{axis}[width=0.85\linewidth, height=6cm, xmin=0, xmax=100, ymin=89, ymax=93.6,
  xlabel={Độ trễ p50 mỗi đoạn trên GPU (ms, batch 1)}, ylabel={Macro-F1 trên dev},
  axis x line*=bottom, axis y line*=left, grid=major, grid style={draw=black!12},
  tick label style={font=\small}, label style={font=\small}, clip=false,
  yticklabel style={/pgf/number format/.cd, use comma, fixed}]
\addplot[only marks, mark=*, mark size=4pt, mauvhai] coordinates {(37.96,92.29)};
\node[anchor=west,xshift=4pt,font=\footnotesize] at (axis cs:37.96,92.29) {mmBERT-base (chọn): 92,3; 1.177 MiB};
\addplot[only marks, mark=*, mark size=4pt, navy!70] coordinates {(18.65,92.24)};
\node[anchor=north,yshift=-4pt,font=\footnotesize] at (axis cs:18.65,92.24) {BamiBERT: 92,2; 397 MiB};
\addplot[only marks, mark=*, mark size=4pt, navy!70] coordinates {(20.82,92.89)};
\node[anchor=south,yshift=4pt,font=\footnotesize] at (axis cs:20.82,92.89) {ViSoBERT: 92,9; 376 MiB};
\addplot[only marks, mark=*, mark size=4pt, navy!70] coordinates {(82.31,90.03)};
\node[anchor=south,yshift=4pt,font=\footnotesize] at (axis cs:82.31,90.03) {GuardMerge 0,6B: 90,0; 1.177 MiB};
\end{axis}
\end{tikzpicture}

\caption{Bốn ứng viên trên tập dev: Macro-F1 theo độ trễ GPU (batch~1), kèm dung lượng trọng số FP32. Ba bộ mã hóa
SafeVi-DualORC-CRF vượt trội GuardMerge cả về chất lượng lẫn tốc độ.}
\label{fig:p4-ung-vien}
\end{figure}

\begin{table}[H]
\centering
\small
\caption{So sánh bốn ứng viên trên tập dev (F1 trên thang 0--100; FPR và Tô oan tính theo \%). ``Tô oan'': tỉ lệ đoạn
\nhan{AN_TOAN} bị trích xuất bằng chứng. ``Chênh vùng'': hiệu Macro-F1 giữa vùng cao nhất và thấp nhất. Độ trễ đo trên
GPU \so{v2.gpu}, batch~1.}
\label{tab:p4-ung-vien}
\setlength{\tabcolsep}{4pt}
\begin{tabular}{lcccc}
\toprule
\textbf{Chỉ số} & \textbf{mmBERT-base} & \textbf{BamiBERT} & \textbf{ViSoBERT} & \textbf{GuardMerge} \\
\midrule
Macro-F1 (5 nhãn) & \so{dev.mmbert.macro} & \so{dev.bamibert.macro} & \textbf{\so{dev.visobert.macro}} & \so{dev.guard.macro} \\
Binary-F1 (nguy hại/an toàn) & \so{dev.mmbert.bin} & \so{dev.bamibert.bin} & \textbf{\so{dev.visobert.bin}} & \so{dev.guard.bin} \\
FPR đoạn \nhan{AN_TOAN} & \so{dev.mmbert.fpr} & \so{dev.bamibert.fpr} & \textbf{\so{dev.visobert.fpr}} & \so{dev.guard.fpr} \\
Bằng chứng: BC F1 theo ý & \so{dev.mmbert.bcf} & \textbf{\so{dev.bamibert.bcf}} & \so{dev.visobert.bcf} & \so{dev.guard.bcf} \\
Tô oan đoạn \nhan{AN_TOAN} & \so{dev.mmbert.toan} & \so{dev.bamibert.toan} & \textbf{\so{dev.visobert.toan}} & \so{dev.guard.toan} \\
Chênh vùng (Macro-F1) & \textbf{\so{dev.mmbert.gap}} & \so{dev.bamibert.gap} & \so{dev.visobert.gap} & \so{dev.guard.gap} \\
Macro-F1 tiếng lóng ($n = \so{dev.tieng.n}$) & \textbf{\so{dev.mmbert.tieng}} & \so{dev.bamibert.tieng} & \so{dev.visobert.tieng} & \so{dev.guard.tieng} \\
\midrule
Trọng số FP32 (MiB) & \so{dev.mmbert.trongso} & \so{dev.bamibert.trongso} & \textbf{\so{dev.visobert.trongso}} & \so{dev.guard.trongso} \\
Độ trễ p50 (ms) & \so{dev.mmbert.p50} & \textbf{\so{dev.bamibert.p50}} & \so{dev.visobert.p50} & \so{dev.guard.p50} \\
Thời gian huấn luyện (phút) & \so{dev.mmbert.trainphut} & \textbf{\so{dev.bamibert.trainphut}} & \textbf{\so{dev.visobert.trainphut}} & \so{dev.guard.trainphut} \\
\bottomrule
\end{tabular}
\end{table}

Kiểm định ghép cặp McNemar trên từng đoạn của tập dev cho thấy:
\begin{itemize}
  \item mmBERT-base vượt GuardMerge có ý nghĩa thống kê ($p = \so{dev.mc.guard}$) và nhanh hơn hơn hai lần
        (\so{dev.mmbert.p50}~ms so với \so{dev.guard.p50}~ms). Hướng dùng nền LLM 0,6 tỉ tham số tinh chỉnh LoRA bị
        loại.
  \item Ba bộ mã hóa SafeVi-DualORC-CRF hòa nhau về mặt thống kê trên chỉ số tổng (mmBERT-base so với ViSoBERT:
        $p = \so{dev.mc.visobert}$; so với BamiBERT: $p = \so{dev.mc.bamibert}$).
\end{itemize}

Khi tổng điểm hòa về mặt thống kê, nhóm lựa chọn theo các lát cắt sát nhất với mục tiêu cốt lõi của v2:
\textbf{mmBERT-base} dẫn đầu trên hội thoại nhiều lượt, vượt xa hai bộ mã hóa tiếng Việt trên lát cắt tiếng lóng
(\so{dev.mmbert.tieng} so với \so{dev.visobert.tieng} của ViSoBERT và \so{dev.bamibert.tieng} của BamiBERT), đạt độ đồng
đều vùng miền tốt nhất (chênh lệch chỉ \so{dev.mmbert.gap}), đồng thời nền đa ngôn ngữ giúp tổng quát hóa tốt trước ký
tự Unicode biến dạng (C2). Lợi thế nhẹ của ViSoBERT trên tổng điểm dev chủ yếu nằm ở bình luận một dòng, trùng khớp kho
tiền huấn luyện của nó. Dù trọng số FP32 của mmBERT-base lớn hơn (\so{dev.mmbert.trongso}~MiB), sau lượng tử hóa ONNX
INT8 mô hình chỉ còn \so{tb.int8.mib}~MiB (khoảng 15\% giới hạn 2~GB) và suy luận \so{tb.a.int8.8.p50}~ms trên CPU
laptop; BamiBERT được lưu làm phương án dự phòng siêu nhẹ cho thiết bị cấu hình rất yếu.

\subsection{Đóng góp của từng mô-đun (ablation)}
\label{sec:yc3-ablation}

\paragraph{Đọc toàn bộ đoạn hội thoại.}
Trên \so{ctx.n} đoạn nhiều dòng của tập dev, v2 đạt Macro-F1 \so{ctx.full} khi đọc toàn bộ hội thoại nhưng chỉ còn
\so{ctx.last} khi chỉ đọc dòng cuối (tăng \so{ctx.gain} điểm). Kết quả này khẳng định mô hình thực sự khai thác ngữ cảnh
liên lượt để phát hiện các tình huống bộc lộ dần về cuối đoạn. Một bộ phân loại chỉ dựa vào độ dài đoạn chỉ đạt Macro-F1
\so{ctx.lenonly}, chứng minh v2 không học lối tắt dựa vào số dòng dù bình luận thực tế đều dài một dòng còn hội thoại
tổng hợp phần lớn nhiều dòng.

\paragraph{Bộ chuẩn hóa đầu vào v1.}
Khác với v1 phụ thuộc vào bộ chuẩn hóa 5 tầng, v2 học trực tiếp trên văn bản biến dạng thông qua tăng cường trực tuyến,
phạt nhất quán, FGM và nhánh CNN đọc theo byte. Nhóm đo ảnh hưởng khi bật/tắt bộ chuẩn hóa v1 trên các bộ thử của tập
dev (Bảng~\ref{tab:p4-chuan-hoa}) và xác nhận lại sau khi đóng băng mô hình trên toàn bộ tập kiểm thử
(Bảng~\ref{tab:p4-chuan-hoa-test}).

\begin{table}[H]
\centering
\footnotesize
\setlength{\tabcolsep}{3.5pt}
\caption{Ảnh hưởng khi bật/tắt bộ chuẩn hóa v1 trên các bộ thử của tập dev (Macro-F1, thang 0--100). ``Mức tụt'':
Macro-F1 bản biến dạng trừ bản gốc khi tắt chuẩn hóa; trong ngoặc là chênh lệch khi bật so với khi tắt.}
\label{tab:p4-chuan-hoa}
\begin{tabularx}{\linewidth}{@{}Ycccccc@{}}
\toprule
\textbf{Bộ thử} & \textbf{Đã thấy?} & \textbf{$n$} & \textbf{Bản gốc} & \textbf{Mức tụt} & \textbf{Tắt} & \textbf{Bật} \\
\midrule
Đoạn gốc C0 & -- & \so{cn.cleanc0.n} & -- & -- & \textbf{\so{cn.cleanc0.raw}} & \so{cn.cleanc0.norm} (\so{cn.cleanc0.gain}) \\
Bỏ dấu & Có & \so{cn.accentstrip.n} & \so{cn.accentstrip.goc} & \so{cn.accentstrip.tut} & \textbf{\so{cn.accentstrip.raw}} & \so{cn.accentstrip.norm} (\so{cn.accentstrip.gain}) \\
C1 đã thấy & Có & \so{cn.c1indist.n} & \so{cn.c1indist.goc} & \so{cn.c1indist.tut} & \so{cn.c1indist.raw} & \textbf{\so{cn.c1indist.norm}} (\so{cn.c1indist.gain}) \\
C1 giữ ngoài huấn luyện & Không & \so{cn.c1heldoutcontrolled.n} & \so{cn.c1heldoutcontrolled.goc} & \so{cn.c1heldoutcontrolled.tut} & \textbf{\so{cn.c1heldoutcontrolled.raw}} & \so{cn.c1heldoutcontrolled.norm} (\so{cn.c1heldoutcontrolled.gain}) \\
Chữ đồng dạng (homoglyph) & Không & \so{cn.blindspothomoglyph.n} & \so{cn.blindspothomoglyph.goc} & \so{cn.blindspothomoglyph.tut} & \textbf{\so{cn.blindspothomoglyph.raw}} & \so{cn.blindspothomoglyph.norm} (\so{cn.blindspothomoglyph.gain}) \\
Ký tự vô hình & Không & \so{cn.blindspotzerowidth.n} & \so{cn.blindspotzerowidth.goc} & \so{cn.blindspotzerowidth.tut} & \textbf{\so{cn.blindspotzerowidth.raw}} & \so{cn.blindspotzerowidth.norm} (\so{cn.blindspotzerowidth.gain}) \\
\bottomrule
\end{tabularx}
\end{table}

\begin{table}[H]
\centering
\small
\caption{Kiểm chứng bật/tắt bộ chuẩn hóa v1 trên toàn bộ tập kiểm thử sau khi đóng băng mô hình và ngưỡng (thang 0--100;
trong ngoặc là chênh lệch khi bật so với khi tắt).}
\label{tab:p4-chuan-hoa-test}
\begin{tabular}{lccccc}
\toprule
 & & \multicolumn{2}{c}{\textbf{Macro-F1}} & \multicolumn{2}{c}{\textbf{BC F1 theo ý}} \\
\cmidrule(lr){3-4}\cmidrule(lr){5-6}
\textbf{Tập} & \textbf{$n$} & \textbf{Tắt} & \textbf{Bật} & \textbf{Tắt} & \textbf{Bật} \\
\midrule
C0 & \so{cnt.C0.n} & \textbf{\so{cnt.C0.raw}} & \so{cnt.C0.norm} (\so{cnt.C0.gain}) & \so{cnt.C0.bc.raw} & \textbf{\so{cnt.C0.bc.norm}} (\so{cnt.C0.bc.gain}) \\
C1 & \so{cnt.C1.n} & \so{cnt.C1.raw} & \textbf{\so{cnt.C1.norm}} (\so{cnt.C1.gain}) & \so{cnt.C1.bc.raw} & \textbf{\so{cnt.C1.bc.norm}} (\so{cnt.C1.bc.gain}) \\
C2 & \so{cnt.C2.n} & \textbf{\so{cnt.C2.raw}} & \so{cnt.C2.norm} (\so{cnt.C2.gain}) & \textbf{\so{cnt.C2.bc.raw}} & \so{cnt.C2.bc.norm} (\so{cnt.C2.bc.gain}) \\
Biến dạng chưa thấy khi huấn luyện & \so{cnt.chuathay.n} & \textbf{\so{cnt.chuathay.raw}} & \so{cnt.chuathay.norm} (\so{cnt.chuathay.gain}) & \textbf{\so{cnt.chuathay.bc.raw}} & \so{cnt.chuathay.bc.norm} (\so{cnt.chuathay.bc.gain}) \\
\bottomrule
\end{tabular}
\end{table}

Kết quả trên cả tập dev lẫn tập kiểm thử hoàn toàn nhất quán: bộ chuẩn hóa v1 chỉ hỗ trợ một số biến dạng C1 đã biết
(trên dev \so{cn.c1indist.gain}; trên test C1 chủ yếu nhờ chèn ký tự \so{cnt.kieu.chenkytu.gain}), nhưng lại làm giảm
Macro-F1 trên đoạn gốc C0 (\so{cnt.C0.gain}), nhóm điểm mù C2 (\so{cnt.C2.gain}, như ký tự vô hình
\so{cnt.kieu.zerowidth.gain}, chèn dấu cách \so{cnt.kieu.spaceinjection.gain}) và nhóm biến dạng chưa thấy
(\so{cnt.chuathay.gain}) do đôi khi chuyển đổi sai các từ mô hình vốn đã hiểu đúng. Vì vậy, v2 mặc định \textbf{tắt}
bộ chuẩn hóa v1 và suy luận trực tiếp trên văn bản gốc; năng lực chống né lọc nay nằm trực tiếp trong biểu diễn của mô
hình thay vì bộ luật cứng.

\paragraph{Bộ mã hóa nền.}
Từ Bảng~\ref{tab:p4-ung-vien}, giữ nguyên kiến trúc SafeVi-DualORC-CRF và thay đổi bộ mã hóa nền giữ điểm tổng trong
khoảng hòa thống kê, nhưng tạo khác biệt rõ về khả năng xử lý tiếng lóng, cân bằng vùng miền và chi phí triển khai.

\paragraph{Lượng tử hóa FP32 $\to$ INT8.}
Để tối ưu triển khai cục bộ, mô hình xuất sang ONNX được lượng tử hóa INT8 kiểu lai: bảng embedding nén INT8 đối xứng
theo hàng; các lớp attention và lớp vào MLP lượng tử hóa động cả trọng số lẫn đầu vào; lớp ra MLP chỉ lượng tử hóa trọng
số INT8 (tránh cắt ngọn ngoại lai); đầu phân loại, đầu gắn thẻ và tầng CRF giữ nguyên FP32 để bảo toàn hiệu chuẩn xác
suất. Nhờ đó, dung lượng tệp giảm từ \so{tb.fp32.mib}~MiB xuống \so{tb.int8.mib}~MiB (nén khoảng \so{tb.nen} lần), trong
khi chất lượng trên toàn bộ \so{lt.ALL.n} đoạn kiểm thử gần như bảo toàn tuyệt đối (Bảng~\ref{tab:p4-luong-tu}):
Macro-F1 chỉ chênh \so{lt.ALL.d.macro} điểm (không có ý nghĩa thống kê, McNemar $p = \so{lt.ALL.p}$), BC F1 theo ý thậm
chí tăng nhẹ \so{lt.ALL.d.bcy} điểm, và FPR trên C0, C1 giữ nguyên hoàn toàn.

\begin{table}[H]
\centering
\small
\caption{Đánh đổi chất lượng khi lượng tử hóa ONNX FP32 $\to$ INT8 trên tập kiểm thử (thang 0--100; FPR tính theo \%;
cùng tập và cùng ngưỡng đã đóng băng; trong ngoặc là INT8 trừ FP32).}
\label{tab:p4-luong-tu}
\begin{tabular}{lrcccccc}
\toprule
 & & \multicolumn{2}{c}{\textbf{Macro-F1}} & \multicolumn{2}{c}{\textbf{BC F1 theo ý}} & \multicolumn{2}{c}{\textbf{FPR}} \\
\cmidrule(lr){3-4}\cmidrule(lr){5-6}\cmidrule(lr){7-8}
\textbf{Tập} & \textbf{$n$} & \textbf{FP32} & \textbf{INT8} & \textbf{FP32} & \textbf{INT8} & \textbf{FP32} & \textbf{INT8} \\
\midrule
C0 & \so{lt.C0.n} & \so{lt.C0.fp32.macro} & \so{lt.C0.int8.macro} (\so{lt.C0.d.macro}) & \so{lt.C0.fp32.bcy} & \textbf{\so{lt.C0.int8.bcy}} (\so{lt.C0.d.bcy}) & \so{lt.C0.fp32.fpr} & \so{lt.C0.int8.fpr} \\
C1 & \so{lt.C1.n} & \so{lt.C1.fp32.macro} & \so{lt.C1.int8.macro} (\so{lt.C1.d.macro}) & \so{lt.C1.fp32.bcy} & \textbf{\so{lt.C1.int8.bcy}} (\so{lt.C1.d.bcy}) & \so{lt.C1.fp32.fpr} & \so{lt.C1.int8.fpr} \\
C2 & \so{lt.C2.n} & \so{lt.C2.fp32.macro} & \so{lt.C2.int8.macro} (\so{lt.C2.d.macro}) & \so{lt.C2.fp32.bcy} & \textbf{\so{lt.C2.int8.bcy}} (\so{lt.C2.d.bcy}) & \so{lt.C2.fp32.fpr} & \so{lt.C2.int8.fpr} \\
\midrule
\textbf{Tất cả} & \so{lt.ALL.n} & \so{lt.ALL.fp32.macro} & \so{lt.ALL.int8.macro} (\so{lt.ALL.d.macro}) & \so{lt.ALL.fp32.bcy} & \textbf{\so{lt.ALL.int8.bcy}} (\so{lt.ALL.d.bcy}) & \so{lt.ALL.fp32.fpr} & \so{lt.ALL.int8.fpr} \\
\bottomrule
\end{tabular}
\end{table}

\subsection{Chi phí vận hành trên thiết bị đích}
\label{sec:yc3-chi-phi}

Bảng~\ref{tab:p4-chi-phi} đối chiếu tổng quan chi phí triển khai và vận hành giữa ba hệ thống.

\begin{table}[H]
\centering
\small
\caption{Tổng hợp chi phí vận hành của ba hệ thống. ``--'': chưa đo hoặc không đo được từ phía thiết bị người dùng.}
\label{tab:p4-chi-phi}
\begin{tabularx}{\linewidth}{lYYY}
\toprule
\textbf{Chỉ tiêu} & \textbf{v1 (ViHateT5)} & \textbf{gpt-oss-120b (API)} & \textbf{v2 (mmBERT-base)} \\
\midrule
Môi trường chạy & Trình duyệt (ONNX) & Máy chủ bên ngoài & CPU laptop, điện thoại Android hoặc trình duyệt \\
Dung lượng mô hình & \so{v1.onnx}~MiB (ONNX FP32) & -- & \textbf{\so{tb.int8.mib}~MiB} (ONNX INT8; FP32 \so{v2.trongso}~MiB) \\
Bộ nhớ khi chạy & -- & Không đo được phía client & RAM đỉnh \so{tb.a.int8.8.ram}~MiB (laptop CPU, 8 luồng); \so{tb.mob.4.ram}~MiB (Android, 4 luồng); trình duyệt tăng \so{tb.web.ram.tang}~MiB \\
Độ trễ mỗi đoạn & -- & p50 \so{api.C0.p50}~s, p95 \so{api.C0.p95}~s (\so{api.songsong} yêu cầu song song) & Laptop CPU (8 luồng): \textbf{p50 \so{tb.a.int8.8.p50}~ms, p95 \so{tb.a.int8.8.p95}~ms}; Android (4 luồng): \textbf{p50 \so{tb.mob.4.p50}~ms, p95 \so{tb.mob.4.p95}~ms}; trình duyệt (4 luồng): p50 \so{tb.web.p50}~ms, p95 \so{tb.web.p95}~ms \\
Năng lượng mỗi đoạn & -- & Không đo được (máy chủ) & \so{tb.b.mj}~mJ (CPU, đo pin), \so{tb.web.mj}~mJ (trình duyệt) \\
Chi phí API & 0 & \so{api.C0.usd1000}~USD / 1.000 đoạn & 0 \\
Gửi tin nhắn ra ngoài & Không & Có & Không \\
Chi phí huấn luyện & (vòng loại) & Không huấn luyện & \so{dev.mmbert.trainphut}~phút, đỉnh \so{dev.mmbert.trainvram}~MiB VRAM (\so{v2.gpu}) \\
\bottomrule
\end{tabularx}
\end{table}

\begin{itemize}
  \item \textbf{LLM qua API rẻ nhưng không phù hợp mục tiêu sản phẩm:} dù chi phí mỗi 1.000 đoạn là
        \so{api.C0.usd1000}~USD, phương án API buộc phải gửi nội dung tin nhắn riêng tư ra máy chủ bên ngoài, có độ trễ
        lớn (p50 \so{api.C0.p50}~s, chậm hơn trên một trăm lần so với v2 trên CPU) và phụ thuộc kết nối mạng.
  \item \textbf{Chi phí xây dựng dữ liệu huấn luyện thấp:} toàn bộ quy trình gán nhãn lại, sinh hội thoại và lọc độc lập
        chỉ tốn \so{cp.tong.usd}~USD (Mục~\ref{sec:p1-chiphi}), là chi phí một lần trong giai đoạn phát triển.
\end{itemize}

\paragraph{Đo hiệu năng, năng lượng và độ nhất quán trên thiết bị đích.}
Mô hình ONNX INT8 được đánh giá trực tiếp ở chế độ tương tác \textbf{batch~1} (xử lý từng đoạn khi người dùng nhận tin)
trên ba môi trường đích: CPU laptop (\so{tb.a.cpu}, \texttt{onnxruntime}), điện thoại Android (\so{tb.mob.may},
\texttt{onnxruntime-android}) và trình duyệt (\so{tb.web.trinhduyet}, \so{tb.web.backend}, \texttt{onnxruntime-web}) qua
ứng dụng web v2 (Bảng~\ref{tab:p4-thiet-bi}, Bảng~\ref{tab:p4-nang-luong} và Bảng~\ref{tab:p4-khop}).

\begin{table}[H]
\centering
\small
\setlength{\tabcolsep}{4pt}
\caption{Độ trễ mỗi đoạn (ms, batch~1), RAM đỉnh và thời gian nạp mô hình của v2 trên thiết bị đích.}
\label{tab:p4-thiet-bi}
\begin{tabularx}{\linewidth}{@{}Ycrrrrr@{}}
\toprule
\textbf{Cấu hình} & \textbf{Luồng} & \textbf{$n$} & \textbf{p50 (ms)} & \textbf{p95 (ms)} & \textbf{RAM đỉnh (MiB)} & \textbf{Nạp (s)} \\
\midrule
PyTorch FP32 (CPU laptop) & 8 & \so{tb.a.torch.8.n} & \so{tb.a.torch.8.p50} & \so{tb.a.torch.8.p95} & \so{tb.a.torch.8.ram} & \so{tb.a.torch.8.nap} \\
ONNX FP32 (CPU laptop) & 8 & \so{tb.a.fp32.8.n} & \so{tb.a.fp32.8.p50} & \so{tb.a.fp32.8.p95} & \so{tb.a.fp32.8.ram} & \so{tb.a.fp32.8.nap} \\
\textbf{ONNX INT8 (CPU laptop)} & \textbf{8} & \so{tb.a.int8.8.n} & \textbf{\so{tb.a.int8.8.p50}} & \textbf{\so{tb.a.int8.8.p95}} & \textbf{\so{tb.a.int8.8.ram}} & \textbf{\so{tb.a.int8.8.nap}} \\
ONNX INT8 (CPU laptop) & 4 & \so{tb.a.int8.4.n} & \so{tb.a.int8.4.p50} & \so{tb.a.int8.4.p95} & \so{tb.a.int8.4.ram} & \so{tb.a.int8.4.nap} \\
ONNX INT8 (CPU laptop) & 1 & \so{tb.a.int8.1.n} & \so{tb.a.int8.1.p50} & \so{tb.a.int8.1.p95} & \so{tb.a.int8.1.ram} & \so{tb.a.int8.1.nap} \\
\midrule
\textbf{ONNX INT8 (Android)} & \textbf{4} & \so{tb.mob.4.n} & \textbf{\so{tb.mob.4.p50}} & \textbf{\so{tb.mob.4.p95}} & \textbf{\so{tb.mob.4.ram}}$^{\dagger}$ & \textbf{\so{tb.mob.4.nap}} \\
ONNX INT8 (trình duyệt) & 4 & \so{tb.web.n} & \so{tb.web.p50} & \so{tb.web.p95} & +\so{tb.web.ram.tang}$^{*}$ & \so{tb.web.nap} \\
\bottomrule
\end{tabularx}

\smallskip
\parbox{\linewidth}{\footnotesize $^{\dagger}$ Đo trên \so{tb.mob.may} ở cấu hình tối ưu (Total~PSS thực chiếm
\so{tb.mob.4.pss}~MiB).\\
$^{*}$ Phần RAM trình duyệt tăng thêm khi suy luận so với trang trống (\so{tb.web.ram.trong}~MiB); đỉnh toàn trình duyệt
là \so{tb.web.ram.dinh}~MiB (tạm đạt \so{tb.web.ram.nap}~MiB lúc nạp rồi giảm). Dung lượng tải lần đầu là
\so{tb.web.taive}~MiB, sau đó lưu trong bộ nhớ đệm để chạy ngoại tuyến.}
\end{table}

\begin{table}[H]
\centering
\footnotesize
\setlength{\tabcolsep}{3.5pt}
\caption{Năng lượng tiêu thụ của bản ONNX INT8 đo từ cảm biến xả pin Windows khi chạy liên tục trên laptop (4 luồng).
``Tăng thêm'': công suất đang chạy trừ công suất lúc nghỉ.}
\label{tab:p4-nang-luong}
\begin{tabularx}{\linewidth}{@{}lCCCCC@{}}
\toprule
\textbf{Môi trường} & \textbf{Nghỉ (W)} & \textbf{Đang chạy (W)} & \textbf{mJ/đoạn (tăng thêm)} & \textbf{mJ/đoạn (cả máy)} & \textbf{Wh/1.000 đoạn (tăng thêm)} \\
\midrule
CPU (4 luồng) & \so{tb.b.pnghi} & \so{tb.b.ptai} & \textbf{\so{tb.b.mj}} & \textbf{\so{tb.b.mjtong}} & \textbf{\so{tb.b.wh}} \\
Trình duyệt (4 luồng) & \so{tb.web.pnghi} & \so{tb.web.ptai} & \so{tb.web.mj} & \so{tb.web.mjtong} & \so{tb.web.wh} \\
\bottomrule
\end{tabularx}
\end{table}

\begin{table}[H]
\centering
\small
\caption{Chất lượng suy luận của bản triển khai ONNX INT8 trên toàn bộ \so{tb.int8.n} đoạn kiểm thử (C0, C1 và C2; thang
0--100).}
\label{tab:p4-khop}
\begin{tabularx}{\linewidth}{YCCC}
\toprule
\textbf{Môi trường suy luận} & \textbf{Macro-F1} & \textbf{BC F1 theo ý} & \textbf{Khớp nhãn với FP32} \\
\midrule
PyTorch FP32 (bản đánh giá chính thức) & \so{tb.fp32.macro} & \so{tb.fp32.bcy} & -- \\
ONNX INT8 (CPU, Python) & \so{tb.int8.macro} & \textbf{\so{tb.int8.bcy}} & \so{tb.int8.khop} \\
ONNX INT8 (trình duyệt, WASM 4 luồng) & \so{tb.web.macro} & \so{tb.web.bcy} & \so{tb.web.khop.torch} \\
\bottomrule
\end{tabularx}
\end{table}

Một số kết luận chính về vận hành thực tế:
\begin{itemize}
  \item \textbf{Hiệu năng cao và tiết kiệm tài nguyên trên CPU laptop:} so với PyTorch FP32 (8 luồng), bản ONNX INT8 rút
        ngắn thời gian nạp mô hình hơn 13 lần (\so{tb.a.torch.8.nap}~s $\to$ \so{tb.a.int8.8.nap}~s), giảm RAM đỉnh gần
        5 lần (\so{tb.a.torch.8.ram}~MiB $\to$ \so{tb.a.int8.8.ram}~MiB) và đạt độ trễ p50 chỉ \so{tb.a.int8.8.p50}~ms
        (nhanh hơn cả FP32 trên GPU \so{v2.gpu} ở batch~1 và vượt xa yêu cầu phản hồi thời gian thực).
  \item \textbf{Triển khai trực tiếp trên điện thoại Android đạt hiệu năng tương đương laptop:} trên \so{tb.mob.may} (4
        luồng), bản ONNX INT8 đạt độ trễ p50 \so{tb.mob.4.p50}~ms, p95 \so{tb.mob.4.p95}~ms (nhanh hơn laptop CPU ở cùng
        mức 4 luồng là \so{tb.a.int8.4.p50}~ms), nạp mô hình trong \so{tb.mob.4.nap}~s và RAM đỉnh chỉ
        \so{tb.mob.4.ram}~MiB (Total~PSS \so{tb.mob.4.pss}~MiB, thấp hơn 31\% so với laptop), đồng thời khớp nhãn
        \so{tb.mob.parity.khop} với bản FP32 trên tập đối chiếu \so{tb.mob.parity.n} đoạn C0.
  \item \textbf{Chất lượng ổn định giữa các runtime cục bộ:} trên trọn bộ \so{tb.int8.n} đoạn kiểm thử, bản ONNX INT8
        trên CPU và trình duyệt đạt Macro-F1 và BC F1 gần như tương đương PyTorch FP32 (Bảng~\ref{tab:p4-khop}), khớp
        nhãn lần lượt \so{tb.int8.khop} và \so{tb.web.khop.torch} (chênh lệch nhỏ phát sinh từ sai khác số học của lượng
        tử hóa động ở các đoạn nằm sát ranh giới hai nhãn; riêng trên C0/C1 độ khớp đạt \so{tb.int8.C0.khop} và
        \so{tb.int8.C1.khop}).
  \item \textbf{Suy luận trực tiếp trong trình duyệt và tiêu thụ năng lượng:} bản trình duyệt chạy ngoại tuyến qua
        WebAssembly đạt p50 \so{tb.web.p50}~ms (chậm hơn bản CPU gốc do tập lệnh vector WASM chỉ rộng 128 bit, nhưng vẫn
        đáp ứng tốt khi xử lý tin nhắn đến theo từng lượt). Phép đo năng lượng sử dụng cảm biến pin Windows qua cửa sổ
        chạy liên tục tối thiểu \so{tb.b.giay}~giây kèm hai cửa sổ nghỉ trước--sau, ghi nhận mức tiêu thụ tăng thêm rất
        thấp (\so{tb.b.wh}~Wh trên CPU và \so{tb.web.wh}~Wh trên trình duyệt cho mỗi 1.000 đoạn).
\end{itemize}

\section{Hạn chế và hướng phát triển}
\label{sec:han-che}

\paragraph{Hạn chế hiện tại.}
\begin{itemize}
  \item \textbf{Cỡ mẫu tập kiểm thử còn khiêm tốn} (\so{te.tong} đoạn gốc do thành viên nhóm tự biên soạn), khiến các
        lát cắt nhỏ theo phương ngữ hoặc kênh giao tiếp chịu biên độ dao động lớn hơn.
  \item \textbf{Tỉ lệ báo nhầm trên nhóm kịch bản chưa thấy còn cao,} chủ yếu do ngưỡng phân loại hiện dùng chung cho
        toàn bộ mô hình và thiếu các đoạn đối chứng an toàn tương ứng cho các tình huống ngoài tập huấn luyện.
  \item \textbf{Hiệu năng giảm trên kịch bản hoàn toàn mới} (Mục~\ref{sec:yc2}) và độ trễ trên trình duyệt WebAssembly
        còn cao hơn bản CPU gốc.
\end{itemize}

\paragraph{Hướng phát triển tiếp theo.}
\begin{itemize}
  \item Bổ sung các kịch bản chưa thấy (vay trực tuyến, giả giao hàng, kiểm soát người yêu, chuyển nhầm tiền) cùng tin
        nhắn thực tế đã ẩn danh và các đoạn đối chứng an toàn để nâng cao độ khái quát.
  \item Xây dựng tập kiểm định (validation) độc lập do người gán để hiệu chỉnh ngưỡng quyết định riêng cho từng nhãn
        nguy hại, tối ưu đồng thời FPR và độ phủ phát hiện.
  \item Tăng tốc suy luận trong trình duyệt bằng tăng tốc phần cứng (WebGPU) và đánh giá mở rộng trên các dòng điện thoại
        tầm trung, cấu hình thấp (kết hợp phương án dự phòng BamiBERT).
\end{itemize}


\end{document}
